% Appendix scaffold for main_iclr.tex.
% Order: A proofs for Section 4; B additional theory (appendix_math.tex,
% reproduced verbatim, not edited); C the effective gradient ratio;
% D/E experiment details (separate file when it exists); F verification record.

\section{Proofs for Section 4}
\label{app:proofs}

Theorem~\ref{thm:hessian} bounds a ratio of two top eigenvalues, and its
proof separates cleanly into a lower bound on the numerator and an upper
bound on the denominator. Appendix~\ref{app:witness} proves the first by
an explicit witness direction in the classifier head; Appendix~\ref{app:cap}
proves the second, which is Lemma~\ref{lem:opcap} of the main text;
Appendix~\ref{app:instance} exhibits a concrete architecture for which
\emph{both} structural hypotheses are theorems rather than assumptions, so
that the disparity law itself is proved outright for at least one member of
each side of the pipeline; Appendix~\ref{app:thm1} assembles the two bounds.
Appendix~\ref{app:classical} records how the disparity relates to the
classical causes of ill-conditioning, and Appendix~\ref{app:thm2_open} states
what remains open around Theorem~\ref{thm:twoscale}. Throughout, Grams are
normalized as empirical means over the $N := N_{\mathrm{data}} HW$ pixel
positions; the unnormalized (summed) Grams differ by the fixed factor $N$,
which does not affect any statement in $M$.

\subsection{The spatial block: an explicit witness}
\label{app:witness}

\begin{lemma}[Spatial block scaling]
\label{lem:phi_scaling}
Let $M$ denote the characteristic width of $g_\phip$ (channel count for a
CNN, embedding dimension for a ViT, hidden layer width for an MLP), let
$N_{\mathrm{cls}} \geq 2$, and suppose $g_\phip$ ends in a dense classifier
head over penultimate features $h_n \in \mathbb{R}^{M_h}$ with
$M_h = \Theta(M)$ whose softmax-weighted mean feature Gram
\begin{equation}
\label{eq:feature_gram}
    S_h^{p} \;:=\; \frac{1}{N} \sum_{n=1}^{N}
        p^{(n)}_{\min}\, h_n h_n^\top,
    \qquad
    p^{(n)}_{\min} := \min_i p^{(n)}_i ,
\end{equation}
satisfies $\lambda_{\max}(S_h^{p}) \geq q_h M_h$ with probability at least
$1 - \delta_h$ over the random initialization, for constants $q_h > 0$ and
$\delta_h < 1$ independent of $M$ (the \emph{head hypothesis}: the features
carry a common direction of $\Theta(M)$ energy across inputs, on pixels
whose initial predictions are non-degenerate). Then the spatial
Gauss--Newton block of~\eqref{eq:ggn} satisfies
\[
    \lambda_{\max}(G_{\phip\phip}) \;=\; \Omega(M)
\]
in expectation over random Kaiming (variance-scaled Gaussian)
initialization. The logit Gram $J_\phip^\top J_\phip \succeq G_{\phip\phip}$
inherits the bound unconditionally; the residual Gram
$J_\phip^{\mathrm{res}\top} J_\phip^{\mathrm{res}}$ satisfies it under the
analogous hypothesis on the $p^2$-weighted Gram $S_h^{p^2}$. For
architecture families with $\Cg = \Theta(M^2)$ (fixed depth with at least
two consecutive width-proportional layers), equivalently
$\lambda_{\max}(G_{\phip\phip}) = \Omega(\sqrt{\Cg})$.
\end{lemma}

The proof has three steps: (1)~context---the exact width-linear law of
Karakida et al.~\citep{karakida2019universal} for the top Fisher
eigenvalue in their model class, which we use as the asymptotic benchmark
but do \emph{not} invoke directly; (2)~an exact spectral lower bound on the
softmax inner matrix at initialization; (3)~an elementary explicit witness
direction in the head weights that yields $\Omega(M)$ self-containedly.

\begin{proof}
\emph{Step 1: The width-linear law of Karakida et al.\ (context).}
Karakida, Akaho, and Amari~\citep{karakida2019universal} analyze fully
connected networks of depth $L \geq 2$ with layer widths
$M_\ell = \alpha_\ell M$, linear outputs, and variance-scaled Gaussian
(Kaiming-type) initialization
$W^\ell_{ij} \sim \mathcal{N}(0, \sigma_w^2 / M_{\ell-1})$,
$b^\ell_i \sim \mathcal{N}(0, \sigma_b^2)$. For the empirical Fisher
information of the squared-loss model,
$F = \tfrac{1}{T} \sum_{t=1}^{T} \sum_{k} \nabla f_k(x_t) \nabla
f_k(x_t)^\top$ (the empirical-mean Gram of the output Jacobian), their
Theorem~4 gives, in the limit $M \gg 1$,
\[
    \lambda_{\max}(F)
    \;=\; \alpha \left( \tfrac{T-1}{T}\,\kappa_2
        + \tfrac{1}{T}\,\kappa_1 \right) M
    \;=\; \Theta(M),
\]
where $\kappa_1, \kappa_2 = \Theta(1)$ are macroscopic constants computed
from the layer-wise signal-propagation recursions and
$\alpha = \sum_\ell \alpha_\ell \alpha_{\ell - 1}$. Their Theorem~1
simultaneously shows the \emph{mean} eigenvalue is
$m_\lambda = C \kappa_1 / M = \mathcal{O}(1/M)$: as the network widens, the
spectrum becomes extremely skewed---almost all directions flatten while the
top eigenvalue grows linearly in width. We record two cautions. First,
Theorem~4 fixes the \emph{output} layer at $C = \mathcal{O}(1)$ but requires
the input and hidden widths to scale proportionally with $M$; a classifier
reading a fixed bottleneck $K$ violates the input-side requirement (its
coefficient $\alpha$ degenerates there), so the theorem cannot be applied
verbatim to the spatial modules considered here. Second, both statements are
asymptotic mean-field equalities with no finite-$M$ rate. Steps~2--3
therefore prove the $\Omega(M)$ bound self-containedly; Karakida's law
serves as the sharp benchmark within its model class, validated numerically
in their Figure~1.

\emph{Step 2: The softmax inner matrix at initialization.}
The loss is softmax cross-entropy, not squared loss, so the spatial
Gauss--Newton block carries the softmax inner matrix
\[
    G_{\phip\phip}
    \;=\; \frac{1}{N} \sum_{n=1}^{N}
        J^{(n)\top} H_\tau^{(n)} J^{(n)},
    \qquad
    H_\tau \;:=\; \mathrm{diag}(p) - p p^\top,
\]
where $J^{(n)} := \partial \hat{y}^{(n)} / \partial \phip$ is the per-pixel
logit Jacobian and $p = \mathrm{softmax}(\hat{y})$. For any
$v \in \mathbb{R}^{N_{\mathrm{cls}}}$, the quadratic form is a variance:
$v^\top H_\tau v = \sum_i p_i v_i^2 - (\sum_i p_i v_i)^2 =
\mathrm{Var}_{i \sim p}(v_i) \geq p_{\min} \sum_i (v_i - \bar{v}_p)^2
\geq p_{\min} \norm{P_\perp v}^2$, where $p_{\min} := \min_i p_i$,
$\bar{v}_p := \sum_i p_i v_i$, the last inequality uses that the unweighted
mean minimizes the sum of squared deviations, and
$P_\perp := I - \tfrac{1}{N_{\mathrm{cls}}} \mathbf{1}\mathbf{1}^\top$
projects onto the mean-zero subspace $\mathbf{1}^\perp$. Hence
\[
    H_\tau \;\succeq\; p_{\min} \, P_\perp .
\]
The bound is used per pixel in Step~3; no minimum over pixels is taken, so
no uniform-over-pixels lower bound on $p_{\min}$ is needed anywhere in the
proof. (For orientation: at Kaiming initialization the logits are
$\mathcal{O}(1)$ uniformly in $M$, so $p^{(n)}_{\min}$ is of order
$1/N_{\mathrm{cls}}$ for typical $n$.)

\emph{Step 3: An explicit witness direction.}
Write the classifier head of $g_\phip$ as $\hat{y}_n = W h_n + b$ with
penultimate features $h_n = h(x_n) \in \mathbb{R}^{M_h}$,
$M_h = \Theta(M)$. A perturbation of the head weights is a matrix
$V \in \mathbb{R}^{N_{\mathrm{cls}} \times M_h}$; it is a unit vector in
parameter space precisely when $\norm{V}_F = 1$, under any vectorization
convention, and the logit perturbation it induces at pixel $n$ is $V h_n$.
Let $u$ be a top unit eigenvector of $S_h^{p}$ in~\eqref{eq:feature_gram},
and fix a unit contrast $c \in \mathbb{R}^{N_{\mathrm{cls}}}$ with
$c \perp \mathbf{1}$ (possible since $N_{\mathrm{cls}} \geq 2$); by Step~2,
at every pixel
\[
    c^\top H_\tau^{(n)} c
    \;=\; \mathrm{Var}_{i \sim p^{(n)}}(c_i)
    \;\geq\; p^{(n)}_{\min} \norm{P_\perp c}^2
    \;=\; p^{(n)}_{\min}.
\]
Take the rank-one witness $V := c\, u^\top$, for which
$\norm{V}_F = \norm{c}\norm{u} = 1$ and $V h_n = (h_n^\top u)\, c$. Then
\begin{equation}
\label{eq:witness_bound}
\begin{aligned}
    \lambda_{\max}(G_{\phip\phip})
    &\;\geq\; \frac{1}{N} \sum_{n=1}^{N}
        (V h_n)^\top H_\tau^{(n)} (V h_n)
    \;=\; \frac{1}{N} \sum_{n=1}^{N}
        (h_n^\top u)^2\, c^\top H_\tau^{(n)} c\\
    &\;\geq\; u^\top S_h^{p} u
    \;=\; \lambda_{\max}(S_h^{p}).
\end{aligned}
\end{equation}
Note what~\eqref{eq:witness_bound} is: a deterministic, hypothesis-free
inequality---one identity and one application of Step~2---valid at every
parameter value. All $N$ pixels contribute, and no minimum over pixels is
taken, so no union bound and no factor of $N$ enters.

The head hypothesis is exactly $\lambda_{\max}(S_h^{p}) \geq q_h M_h$ with
probability at least $1 - \delta_h$. On that event
$\lambda_{\max}(G_{\phip\phip}) \geq q_h M_h = \Omega(M)$; since
$\lambda_{\max} \geq 0$ always,
\[
    \E\!\left[\lambda_{\max}(G_{\phip\phip})\right]
    \;\geq\; (1 - \delta_h)\, q_h M_h
    \;=\; \Omega(M),
\]
which is the claim.

\emph{The other two conventions.}
The bound just proved is for the Gauss--Newton block
$G_{\phip\phip} = J_\phip^\top H_\tau J_\phip$ of~\eqref{eq:ggn}, the object
appearing in Definition~\ref{def:dcurv}. Two neighbouring matrices inherit
it. For the \emph{logit} Gram, $H_\tau \preceq I$ gives
$J_\phip^\top J_\phip \succeq G_{\phip\phip}$, so
$\lambda_{\max}(J_\phip^\top J_\phip) \geq \lambda_{\max}(G_{\phip\phip}) =
\Omega(M)$ immediately. For the \emph{residual} Gram
$J_\phip^{\mathrm{res}\top} J_\phip^{\mathrm{res}} = J_\phip^\top H_\tau^2
J_\phip$, note that $H_\tau$ is symmetric with $H_\tau \mathbf{1} = 0$, so
$\mathbf{1}^\perp$ is invariant and, by Step~2,
$H_\tau^{(n)} \succeq p^{(n)}_{\min} I$ there; squaring the eigenvalues on
that subspace gives $(H_\tau^{(n)})^2 \succeq (p^{(n)}_{\min})^2 P_\perp$,
so the identical witness argument runs with the weights
$(p^{(n)}_{\min})^2$ in place of $p^{(n)}_{\min}$ and yields
$\lambda_{\max}(J_\phip^{\mathrm{res}\top} J_\phip^{\mathrm{res}}) \geq
\lambda_{\max}(S_h^{p^2}) = \Omega(M)$ under the corresponding head
hypothesis.
\end{proof}

\paragraph{On the head hypothesis, and its counterpart in Karakida et al.}
Write $S_h := N^{-1} \sum_n h_n h_n^\top$ for the unweighted mean feature
Gram. The hypothesis $\lambda_{\max}(S_h^{p}) = \Omega(M_h)$ asks two things
at once: that the penultimate features carry a common direction of
$\Theta(M_h)$ energy across inputs ($\lambda_{\max}(S_h) = \Omega(M_h)$),
and that the initial predictions are non-degenerate on the pixels supplying
that energy. The second is mild at initialization---the logits are
$\mathcal{O}(1)$ uniformly in $M$, since the head weights have variance
$\Theta(1/M_h)$ against $\norm{h_n}^2 = \mathcal{O}(M_h)$, so
$p^{(n)}_{\min}$ is bounded below by a constant of order
$1/N_{\mathrm{cls}}$ on all but a small fraction of pixels, and $S_h^{p}$
dominates that constant times the unweighted Gram restricted to those
pixels. (Proposition~\ref{prop:verified_instance} below turns this sketch
into a proof for a concrete head.) The first is the substantive condition.
It is implied by $\norm{\bar{h}}^2 \geq q M_h$ on the \emph{mean} feature
vector $\bar{h} := N^{-1}\sum_n h_n$: for $\bar{u} =
\bar{h}/\norm{\bar{h}}$,
\[
    \lambda_{\max}(S_h)
    \;\geq\; \bar{u}^\top S_h \bar{u}
    \;=\; \frac{1}{\norm{\bar{h}}^2} \cdot
        \frac{1}{N}\sum_n (h_n^\top \bar{h})^2
    \;\geq\; \frac{1}{\norm{\bar{h}}^2}
        \left( \frac{1}{N}\sum_n h_n^\top \bar{h} \right)^{\!2}
    \;=\; \frac{\norm{\bar{h}}^4}{\norm{\bar{h}}^2}
    \;=\; \norm{\bar{h}}^2,
\]
the middle step being Jensen's inequality (the mean of squares dominates the
square of the mean). At variance-scaled initialization each entry of $h_n$
has $\Theta(1)$ second moment (the signal-propagation recursions of
Karakida et al.~\citep{karakida2019universal}, Eq.~(9)), so the condition
holds whenever the activations are correlated across inputs---in particular
for any nonnegative or nonzero-mean activation (ReLU and its variants), and
for odd activations whenever the bias variance $\sigma_b^2 > 0$ or the
inputs are themselves correlated.

Per-entry second moments alone are \emph{not} enough:
$\norm{h_n}^2 = \Theta(M_h)$ for every $n$ gives only
$\mathrm{tr}(S_h) = \Theta(M_h)$, hence
$\lambda_{\max}(S_h) \geq \Theta(M_h)/\min(N, M_h)$, which is weaker by a
factor of $N$. Cross-input correlation, not per-vector norm, is what makes
the top eigenvalue grow. The same quantity carries the width-linear term in
Karakida et al.'s Theorem~4: their $\kappa_2$ is built from the cross-input
activation correlations $\hat{q}_{st}$. Their footnote notes that odd
activations with $\sigma_b = 0$ give $\hat{q}_{st} = 0$ and $\kappa_2 = 0$,
and is explicit that this exceptional case falls outside their analysis:
``In such exceptional cases, we need to evaluate the lower order terms of
$s_\lambda$ and $\lambda_{\max}$ (outside the scope of this study).'' We
therefore attribute no finite-$T$ claim to their theorem in the degenerate
case. Our hypothesis has the parallel structure, and its degenerate
behaviour \emph{can} be evaluated directly: for odd activations with
$\sigma_b = 0$ and uncorrelated inputs the features lose their common
direction, and the per-vector norms alone give only
$\lambda_{\max}(S_h) \geq \mathrm{tr}(S_h)/\min(N, M_h)$. In our regime
$N \gg M_h$, so this floor is $\Theta(1)$---\emph{no} width scaling.
Width-linear growth in the degenerate case would require a separate
argument, which we do not attempt.

\begin{remark}[Scope, and the role of the mean-field law]
\label{rem:meanfield_scope}
The witness argument uses only: (a)~a dense classifier head over penultimate
features of dimension $\Theta(M)$; (b)~the head hypothesis
$\lambda_{\max}(S_h^{p}) \geq q_h M_h$ at initialization---a
\emph{cross-input correlation} condition, which variance-scaled
initialization does not by itself deliver (see the per-vector-norm
counterexample above) but which holds for the standard activation choices
discussed there; and (c)~softmax cross-entropy with
$N_{\mathrm{cls}} \geq 2$. Given (b), the inequality
$\lambda_{\max}(G_{\phip\phip}) \geq \lambda_{\max}(S_h^{p})$ is
deterministic and architecture-agnostic, so the argument applies to
per-pixel dense heads on MLP, CNN (channel count $M$), and ViT (embedding
dimension $M$) spatial modules alike. For fully connected modules (b)
follows from standard signal propagation at Kaiming initialization; for CNN
and ViT heads it is an explicit assumption, to be checked. No i.i.d.\
Gaussian input assumption is required. Karakida et al.'s Theorem~4 plays the
role of the sharp asymptotic benchmark: within its model class (fully
connected, input and hidden widths proportional to $M$, output fixed at
$C = \mathcal{O}(1)$, squared loss), it upgrades the witness lower bound to
an exact width-linear law with explicit constant, and their Theorem~1
supplies the complementary fact that the \emph{mean} eigenvalue shrinks as
$\mathcal{O}(1/M)$, so curvature concentrates in a vanishing fraction of
directions as the module widens. The loss transfer from squared loss to
cross-entropy is handled exactly at initialization by Step~2, at the cost of
the constant $p_{\min}$.
\end{remark}

\begin{remark}[Empirical status, and width versus parameter count]
\label{rem:phi_scaling_empirical}
Karakida et al.'s own numerical validation (their Figure~1, right column)
confirms $\lambda_{\max}$ growing linearly in $M$, matching the exact
constant of their Theorem~4, for tanh, ReLU, and linear fully connected
networks. The theorem's native variable is width, not parameter count: for
architecture families with $\Cg = \Theta(M^2)$ the equivalent
parameter-count statement is $\lambda_{\max} = \Omega(\sqrt{\Cg})$
(log--log slope $0.5$ against $\Cg$), whereas the shallow family measured in
Section~4 has $\Cg = \Theta(M)$, for which the width-linear law predicts
slope $1.0$ against $\Cg$ as well as against $M$. The two parameter-count
forms must not be conflated. The measured slopes, their dependence on the
initialization scheme---which controls whether either variance term of the
head hypothesis is width-independent---and the levels that accompany them
are reported in Section~4 and Appendix~\ref{app:synthetic}; we quote slope
and level together, and read finite-range fits as finite-range fits.
\end{remark}

\subsection{The spectral block: an operator-norm cap}
\label{app:cap}

The proof of Lemma~\ref{lem:opcap} factors the spectral Jacobian via the
chain rule and bounds each factor's operator norm using
Assumptions~\ref{ass:linear}, \ref{ass:subg}, and~\ref{ass:inputlip}. We
bound the \emph{logit} Jacobian $J_\thetap = \partial \hat{y}/\partial
\thetap$ of~\eqref{eq:jacobians}; the Gauss--Newton block and the residual
Jacobian then inherit the bound, since both insert factors of $H_\tau$ with
$0 \preceq H_\tau \preceq I$.

\begin{proof}[Proof of Lemma~\ref{lem:opcap}]
\emph{Step 1: Chain-rule factorisation.}
With $Z = f_\thetap(X)$ and $\hat{y} = g_\phip(Z)$, the chain rule gives, in
the normalization of~\eqref{eq:jacobians} (the $N^{-1/2}$ riding on the
right-hand factor),
\[
    J_\thetap
    \;=\; J_Z \cdot \frac{\partial Z}{\partial \thetap},
    \qquad J_Z := \partial g_\phip / \partial Z ,
\]
where $J_Z$ is unnormalized---the object Assumption~\ref{ass:inputlip}
bounds. Since $g_\phip$ acts on one image at a time, the dataset-stacked
$J_Z$ is block diagonal across images, so its operator norm is the maximum
over per-image blocks and Assumption~\ref{ass:inputlip} applies to the stack
with the same constant $\widetilde{L}$. (For the residual map
$r(\thetap,\phip) = \tau(g_\phip(Z))$ with $\tau(\hat{y}) =
\mathrm{softmax}(\hat{y}) - y_{\mathrm{onehot}}$, the same computation
carries the extra left factor $J_\tau := \partial \tau / \partial \hat{y} =
\mathrm{diag}(p) - p p^\top = H_\tau$, the softmax Jacobian, which satisfies
$\|H_\tau\|_{\mathrm{op}} \leq 1$; see Botev, Ritter, and
Barber~\citep{botev2017practical}.)

\emph{Step 2: Operator-norm inequality.}
Sub-multiplicativity of the operator norm gives
\[
    \|J_\thetap\|_{\mathrm{op}}
    \;\leq\; \|J_Z\|_{\mathrm{op}}
             \cdot \left\| \frac{\partial Z}{\partial \thetap}
             \right\|_{\mathrm{op}}
    \;\leq\; \widetilde{L} \cdot
             \left\| \frac{\partial Z}{\partial \thetap}
             \right\|_{\mathrm{op}},
\]
using Assumption~\ref{ass:inputlip} on the first factor.

\emph{Step 3: The input-Jacobian factor, exactly.}
Under Assumption~\ref{ass:linear}, $Z_n = W^\top x_n$ with
$\thetap = \mathrm{vec}(W)$, $W \in \mathbb{R}^{S \times K}$. The per-pixel
Jacobian is $\partial Z_n / \partial \mathrm{vec}(W) = I_K \otimes x_n^\top$
(Magnus and Neudecker~\citep{magnus1988matrix}, Ch.~9). Stacking over the
$N$ pixel positions with the $N^{-1/2}$ normalization
of~\eqref{eq:jacobians} gives
$\partial Z / \partial \thetap \in \mathbb{R}^{NK \times SK}$ with
\[
    \left( \frac{\partial Z}{\partial \thetap} \right)^{\!\top}
    \frac{\partial Z}{\partial \thetap}
    \;=\; \frac{1}{N}\sum_{n=1}^{N} \left( I_K \otimes x_n \right)
          \left( I_K \otimes x_n^\top \right)
    \;=\; I_K \otimes \Sigma_X ,
\]
hence, \emph{exactly},
\[
    \left\| \frac{\partial Z}{\partial \thetap} \right\|_{\mathrm{op}}^{2}
    \;=\; \lambda_{\max}(I_K \otimes \Sigma_X)
    \;=\; \lambda_{\max}(\Sigma_X),
\]
with $\Sigma_X := N^{-1}\sum_{n=1}^{N} x_n x_n^\top$ the mean per-pixel input
Gram of the main text;
Assumption~\ref{ass:subg} bounds $\lambda_{\max}(\Sigma_X) \leq \Lambda_X$
uniformly in the dataset size $N$. No factor of $K$ appears: the Kronecker
structure replicates the same spectrum $K$ times rather than accumulating
it.

\emph{Step 4: Assembling.}
Combining Steps~2 and~3,
\[
    \|J_\thetap\|_{\mathrm{op}}^{2}
    \;\leq\; \widetilde{L}^{2} \cdot \lambda_{\max}(\Sigma_X)
    \;=:\; C_\thetap .
\]
For the Gauss--Newton block, $G_{\thetap\thetap} = J_\thetap^\top H_\tau
J_\thetap \preceq J_\thetap^\top J_\thetap$ since $H_\tau \preceq I$, so
$\lambda_{\max}(G_{\thetap\thetap}) \leq C_\thetap$; for the residual
Jacobian, $\|H_\tau J_\thetap\|_{\mathrm{op}} \leq
\|J_\thetap\|_{\mathrm{op}}$ gives the same cap. The constant $C_\thetap$
depends on the classifier's input Lipschitz constant $\widetilde{L}$ (which
may grow with the depth of $g_\phip$ but not with its width $M$) and the top
eigenvalue of the mean input Gram $\lambda_{\max}(\Sigma_X)$. It does not
depend on $M$, on the spatial parameter count $\Cg$, or on the dataset size
$N$.
\end{proof}

\begin{remark}[Why an operator-norm cap suffices]
\label{rem:opcap_route}
An alternative route via the smallest positive eigenvalue
$\lambda_{\min}^{+}(J_\thetap^\top J_\thetap)$ requires either an exact
Kronecker factorisation (valid only for per-pixel classifiers) or a
trace-over-rank argument via the softmax residual dimension---both of which
need additional structural assumptions on $g_\phip$. The operator-norm bound
above is significantly simpler, applies to general nonlinear spatial
classifiers (CNN, ViT), and is exactly what Theorem~\ref{thm:hessian} needs:
a $\Cg$-independent cap on the spectral block's \emph{top} eigenvalue, so
that $\Dcurv = \lambda_{\max}(G_{\phip\phip})/
\lambda_{\max}(G_{\thetap\thetap}) \geq \Omega(M)/C_\thetap$.
\end{remark}

\begin{remark}[Softmax Jacobian bound]
\label{rem:softmax_jac}
The bound $\|\mathrm{diag}(p) - p p^\top\|_{\mathrm{op}} \leq 1$ is
standard: $\mathrm{diag}(p) - p p^\top$ is the covariance matrix of the
categorical distribution with parameter $p$ (equivalently, the Jacobian of
the softmax map), hence positive semi-definite with
$\lambda_{\max} \leq \max_i p_i \leq 1$. See Botev, Ritter, and
Barber~\citep{botev2017practical} for its role in Gauss--Newton
computations.
\end{remark}

\subsection{A fully proved instance}
\label{app:instance}

For general heads, the head hypothesis of Lemma~\ref{lem:phi_scaling} and
the uniform input-Jacobian cap of Assumption~\ref{ass:inputlip} are
hypotheses. For one concrete architecture they are both theorems, and the
disparity law follows outright.

\begin{proposition}[A fully proved instance]
\label{prop:verified_instance}
Fix bottleneck features $z_1, \dots, z_N \in \mathbb{R}^K$ with
$\norm{z_n} \leq R$, and let the spatial module's head be the two-layer
ReLU network
\[
    h_n = \mathrm{ReLU}(W_1 z_n + b_1),
    \qquad
    \hat{y}_n = W_2 h_n + b_2,
\]
with Kaiming Gaussian initialization: entries of
$W_1 \in \mathbb{R}^{M_h \times K}$ i.i.d.\
$\mathcal{N}(0, \sigma_w^2/K)$, entries of
$W_2 \in \mathbb{R}^{N_{\mathrm{cls}} \times M_h}$ i.i.d.\
$\mathcal{N}(0, \sigma_w^2/M_h)$, biases i.i.d.\
$\mathcal{N}(0, \sigma_b^2)$ with $\sigma_b > 0$, and
$N_{\mathrm{cls}} \geq 2$. Let the spectral module be the linear per-pixel
$f_\thetap$ of Assumption~\ref{ass:linear} with inputs satisfying
Assumption~\ref{ass:subg}. Then there exist constants $q_h > 0$,
$M_0 < \infty$, and $C_G < \infty$, depending only on
$\sigma_w, \sigma_b, R, K, N_{\mathrm{cls}}$, and
$\mathrm{tr}(\Sigma_X)$---not on $M_h$ or $N$---such that for all
$M_h \geq M_0$, at initialization:
\begin{enumerate}
    \item[(i)] (head hypothesis)
    $\lambda_{\max}(S_h^{p}) \geq q_h M_h$ with probability at least
    $1/2$; consequently
    $\E[\lambda_{\max}(G_{\phip\phip})] \geq \tfrac{1}{2} q_h M_h$;
    \item[(ii)] (spectral-block cap, trace form)
    $\E[\lambda_{\max}(G_{\thetap\thetap})]
    \leq \E[\mathrm{tr}(G_{\thetap\thetap})] \leq C_G :=
    N_{\mathrm{cls}}\, \sigma_w^4\, \mathrm{tr}(\Sigma_X)$,
    uniformly in $M_h$ and $N$;
    \item[(iii)] (instance-level disparity law) with probability at least
    $3/8$, $\Dcurv \geq \left( q_h / (8 C_G) \right) M_h$; consequently
    $\E[\Dcurv] \geq \tfrac{3}{64} (q_h / C_G)\, M_h = \Omega(M)$.
\end{enumerate}
For this architecture, the disparity law of Theorem~\ref{thm:hessian}---not
merely its hypotheses---is proved outright at initialization. What remains
assumption-level in general is the uniform-over-inputs, during-training form
of Assumption~\ref{ass:inputlip}: a per-input Jacobian cap at initialization
follows from the computation in Step~D below, but a supremum over all inputs
by the natural union route would cost a $\sqrt{\log N}$ factor, and
training-time persistence is exactly what Assumption~\ref{ass:inputlip}
postulates. The trace-form cap~(ii) bypasses both issues for the instance.
\end{proposition}

\begin{proof}
\emph{Step A (mean feature vector).} Write
$g_{n,i} = w_i^\top z_n + b_{1,i}$ for the pre-activations, so
$g_{n,i} \sim \mathcal{N}(0, s_n^2)$ with
$s_n^2 = \sigma_w^2 \norm{z_n}^2 / K + \sigma_b^2 \in
[\sigma_b^2,\, s_{\max}^2]$, $s_{\max}^2 := \sigma_w^2 R^2/K + \sigma_b^2$.
Since $\E\,\mathrm{ReLU}(\mathcal{N}(0, s^2)) = s/\sqrt{2\pi}$, each
coordinate of the mean feature vector $\bar{h} = N^{-1} \sum_n h_n$ has
expectation
\[
    m \;:=\; \E[\bar{h}_i]
    \;=\; \frac{1}{N} \sum_n \frac{s_n}{\sqrt{2\pi}}
    \;\geq\; \frac{\sigma_b}{\sqrt{2\pi}} \;>\; 0 .
\]
To obtain constants free of the data and of $N$, we use throughout only the
uniform floor $m \geq m_0 := \sigma_b/\sqrt{2\pi}$. The rows
$(w_i, b_{1,i})$ are i.i.d.\ across $i$, so the coordinates $\bar{h}_i$ are
i.i.d.\ with $\E[\bar{h}_i^2] \leq N^{-1}\sum_n
\E[\mathrm{ReLU}(g_{n,i})^2] \leq s_{\max}^2$ and (being averages of
sub-Gaussian variables) finite fourth moment bounded by a constant
$C_4(s_{\max})$. By Chebyshev's inequality applied to the i.i.d.\ sum
$\norm{\bar{h}}^2 = \sum_i \bar{h}_i^2$, whose mean is at least
$m^2 M_h \geq m_0^2 M_h$,
\[
    \Pr\!\left[ \norm{\bar{h}}^2 < \tfrac{1}{2} m_0^2 M_h \right]
    \;\leq\; \frac{4\, C_4(s_{\max})}{m_0^4\, M_h}
    \;\xrightarrow[M_h \to \infty]{}\; 0 .
\]
On the complementary event $E_A$, the Jensen step above gives
$\lambda_{\max}(S_h) \geq \norm{\bar{h}}^2 \geq \tfrac{1}{2} m_0^2 M_h$, and
moreover the witness direction $\bar{u} = \bar{h}/\norm{\bar{h}}$ certifies
$A := N^{-1} \sum_n (h_n^\top \bar{u})^2 \geq \tfrac{1}{2} m_0^2 M_h$.

\emph{Step B (softmax weights).} Condition on
$\mathcal{F}_1 = \sigma(W_1, b_1)$, which fixes the features $h_n$, the
direction $\bar{u}$, and the energies $a_n := (h_n^\top \bar{u})^2$; the
head randomness $(W_2, b_2)$ is independent of $\mathcal{F}_1$. Call pixel
$n$ \emph{tame} if $\norm{h_n}^2 \leq 2 s_{\max}^2 M_h$. For a tame pixel,
the logits $\hat{y}_n \,|\, \mathcal{F}_1$ have i.i.d.\ Gaussian entries
with variance $\sigma_w^2 \norm{h_n}^2/M_h + \sigma_b^2 \leq
2\sigma_w^2 s_{\max}^2 + \sigma_b^2 =: s_*^2$, so by a Gaussian tail bound
on the logit range, there are constants $c_p(s_*, N_{\mathrm{cls}}) > 0$ and
$\delta_0 \leq 1/16$ with
$\Pr[\, p^{(n)}_{\min} < c_p \mid \mathcal{F}_1 \,] \leq \delta_0$ for every
tame pixel. Untame pixels are exponentially rare: $\norm{h_n}^2/M_h$ is an
average of $M_h$ i.i.d.\ sub-exponential variables with mean
$\leq s_{\max}^2/2$, so $\Pr[\text{$n$ untame}] \leq e^{-c M_h}$
(Bernstein's inequality; Vershynin~\citep{vershynin2018high},
Thm.~2.8.1), and by Cauchy--Schwarz with
$\E\norm{h_n}^4 \leq 3 s_{\max}^4 M_h^2$ (i.i.d.\ coordinates), the
$\mathcal{F}_1$-measurable untame energy
$U := N^{-1} \sum_n a_n \mathbf{1}[\text{$n$ untame}]$ satisfies
\[
    \E[U]
    \;\leq\; \max_n \sqrt{\E\norm{h_n}^4}
        \sqrt{\Pr[\text{$n$ untame}]}
    \;\leq\; \sqrt{3}\, s_{\max}^2 M_h\, e^{-c M_h / 2}
    \;=:\; \varepsilon(M_h)\, M_h ,
\]
with $\varepsilon(M_h) \to 0$ exponentially (we used
$a_n \leq \norm{h_n}^2$, valid since $\bar{u}$ is a unit vector).

\emph{Step C (assembly).} We control the two bad masses by separate
arguments, respecting which randomness each depends on. The untame energy
$U$ is $\mathcal{F}_1$-measurable; by (unconditional) Markov, the event
$E_U := \{ U \leq 16\, \varepsilon(M_h) M_h \}$ has $\Pr[E_U] \geq 15/16$,
and for $M_h \geq M_0$ large enough that
$16\, \varepsilon(M_h) \leq m_0^2/8$, on $E_A \cap E_U$ we have
$U \leq \tfrac{1}{8} m_0^2 M_h \leq \tfrac{1}{4} A$. The $p$-degenerate mass
over \emph{tame} pixels,
$P := N^{-1} \sum_{n\ \mathrm{tame}} a_n
\mathbf{1}[p^{(n)}_{\min} < c_p]$, is $(W_2, b_2)$-random given
$\mathcal{F}_1$, with $\E[P \mid \mathcal{F}_1] \leq \delta_0 A$ by Step~B;
conditional Markov with $\delta_0 \leq 1/16$ gives
$\Pr[\, P > \tfrac{1}{4} A \mid \mathcal{F}_1 \,] \leq 4\delta_0 \leq 1/4$.
On the intersection of the three events---probability at least
$1 - 4 C_4/(m_0^4 M_h) - 1/16 - 1/4 > 1/2$ for $M_h \geq M_0$---the total
bad mass is $P + U \leq \tfrac{1}{2} A$, and
\[
    \lambda_{\max}(S_h^{p})
    \;\geq\; \bar{u}^\top S_h^{p} \bar{u}
    \;\geq\; c_p \left( A - P - U \right)
    \;\geq\; \frac{c_p}{2}\, A
    \;\geq\; \frac{c_p m_0^2}{4}\, M_h .
\]
Take $q_h := c_p m_0^2 / 4 = c_p \sigma_b^2 / (8\pi)$, manifestly free of
the data and of $N$. The in-expectation clause follows from
$\lambda_{\max} \geq 0$ and~\eqref{eq:witness_bound}. This proves~(i).

\emph{Step D (spectral-block cap, trace form).} At any input $z$, the head's
input Jacobian is $\partial \hat{y}/\partial z = W_2 D_z W_1$ with
$D_z = \mathrm{diag}(\mathbf{1}[W_1 z + b_1 > 0])$. Conditioning on
$(W_1, b_1)$ (which fixes $D_z$) and using that the rows of $W_2$ are
i.i.d.\ $\mathcal{N}(0, (\sigma_w^2/M_h) I)$,
\[
    \E \norm{W_2 D_z W_1}_F^2
    \;=\; N_{\mathrm{cls}} \cdot \frac{\sigma_w^2}{M_h}
        \cdot \E \norm{D_z W_1}_F^2
    \;\leq\; N_{\mathrm{cls}} \cdot \frac{\sigma_w^2}{M_h}
        \cdot \E \norm{W_1}_F^2
    \;=\; N_{\mathrm{cls}}\, \sigma_w^4 ,
\]
using $\E\norm{W_1}_F^2 = M_h K \cdot \sigma_w^2/K = M_h \sigma_w^2$---%
uniformly over inputs $z$. (The naive product bound
$\norm{W_2}_{\mathrm{op}} \norm{W_1}_{\mathrm{op}}$ would grow like
$\sqrt{M_h/K}$; the composition is $O(1)$ because the top directions of the
two random factors are generically misaligned---which the Frobenius
computation captures without any alignment analysis.) Now bound the spectral
block. With $J_\thetap^{(n)} = (W_2 D_{z_n} W_1)(I_K \otimes x_n^\top)$ the
per-pixel spectral logit Jacobian (Appendix~\ref{app:cap}, Step~3) and
$H_\tau \preceq I$,
\[
\begin{aligned}
    \E\,\mathrm{tr}(G_{\thetap\thetap})
    &\;\leq\; \frac{1}{N} \sum_{n=1}^{N}
        \E \norm{J_\thetap^{(n)}}_F^2
    \;\leq\; \frac{1}{N} \sum_{n=1}^{N}
        \E \norm{W_2 D_{z_n} W_1}_F^2 \cdot \norm{x_n}^2\\
    &\;\leq\; N_{\mathrm{cls}}\, \sigma_w^4 \cdot
        \frac{1}{N} \sum_n \norm{x_n}^2
    \;=\; N_{\mathrm{cls}}\, \sigma_w^4\, \mathrm{tr}(\Sigma_X),
\end{aligned}
\]
using $\norm{AB}_F \leq \norm{A}_F \norm{B}_{\mathrm{op}}$ and
$\norm{I_K \otimes x_n^\top}_{\mathrm{op}} = \norm{x_n}$. Since
$\lambda_{\max} \leq \mathrm{tr}$ for positive semidefinite matrices,~(ii)
follows with $C_G = N_{\mathrm{cls}} \sigma_w^4 \mathrm{tr}(\Sigma_X)$, free
of $M_h$ and $N$.

\emph{Step E (combining).} If $C_G = 0$ then $\Sigma_X = 0$, the spectral
block vanishes almost surely, and $\Dcurv = +\infty$ by the convention of
Definition~\ref{def:dcurv}; assume $C_G > 0$ below. Let $E_1$ be the event
of~(i) ($\Pr \geq 1/2$) and
$E_2 := \{ \lambda_{\max}(G_{\thetap\thetap}) \leq 8 C_G \}$, which by~(ii)
and Markov has $\Pr[E_2] \geq 7/8$. On $E_1 \cap E_2$ (probability at least
$1 - 1/2 - 1/8 = 3/8$), the witness bound~\eqref{eq:witness_bound} gives
\[
    \Dcurv
    \;=\; \frac{\lambda_{\max}(G_{\phip\phip})}
               {\lambda_{\max}(G_{\thetap\thetap})}
    \;\geq\; \frac{q_h M_h}{8 C_G},
\]
and since $\Dcurv \geq 0$, $\E[\Dcurv] \geq \tfrac{3}{8} \cdot
q_h M_h / (8 C_G) = \tfrac{3}{64} (q_h/C_G) M_h$. This proves~(iii).
\end{proof}

\begin{remark}[Minimality]
\label{rem:instance_minimal}
Proposition~\ref{prop:verified_instance} is deliberately minimal: one hidden
layer, ReLU, bias variance $\sigma_b^2 > 0$ (which guarantees activation
correlation across inputs; cf.\ the $\kappa_2$ discussion above). Its role
is to certify that the two structural hypotheses of
Theorem~\ref{thm:hessian} are \emph{theorems}, not merely assumptions, for
at least one concrete member of each side of the pipeline---so the disparity
law has a fully verified instance. For deeper heads and attention-based
modules the hypotheses revert to assumption status, argued by signal
propagation and checked empirically.
\end{remark}

\begin{remark}[What the instance shows---and what it does not]
\label{rem:instance_interpretation}
The mechanism of Proposition~\ref{prop:verified_instance} deserves a candid
reading. The width-linear spike is generated by the \emph{mean feature
direction}: with $\sigma_b > 0$, every hidden unit has positive mean
activation, so the features of all inputs share a common direction of
$\Theta(M_h)$ energy. This requires no information in the features at all.
Take every bottleneck feature $z_n = 0$, so that all inputs are
indistinguishable: the features $h_n = \mathrm{ReLU}(b_1)$ are identical
across pixels, carry no sample-specific structure whatsoever, and the
proposition applies verbatim. Width-linear downstream curvature is therefore
\emph{not} evidence that the spatial module has found, or will find, useful
structure: $\Dcurv$ measures the downstream module's readiness to
move---the loss curvature its parameters command at initialization---not
the value of what it computes. It is the counterpart, on the spatial side,
of an ambiguity on the spectral side: small spectral curvature cannot by
itself distinguish a direction that is irrelevant to the task from one whose
learning signal is buried behind the bottleneck; here, large spatial
curvature cannot distinguish signal from capacity. It is also why we read
the geometry as a hazard rather than as a mechanism: the capacity advantage
is in place before the data has spoken, so whatever the spatial module can
most easily fit is what the geometry favours. Whether it in fact fits it
first is the separate dynamical question of Theorem~\ref{thm:twoscale} and
of the residual-excitation results of Appendix~\ref{app:math}.
\end{remark}

\begin{remark}[Smoothness scope]
\label{rem:relu_scope}
ReLU is not twice differentiable, but no smoothness enters
Proposition~\ref{prop:verified_instance}: the Gauss--Newton blocks involve
only first-order Jacobians, which at Gaussian initialization exist almost
surely (every preactivation is a.s.\ nonzero). Assumption~\ref{ass:reg}
plays no role here---its function is the well-posedness of the gradient flow
of Theorem~\ref{thm:twoscale}, whose statements are formulated for
absolutely continuous trajectories satisfying the flow almost everywhere and
therefore accommodate piecewise-smooth models.
\end{remark}

\subsection{Proof of Theorem~\ref{thm:hessian}}
\label{app:thm1}

\begin{proof}[Proof of Theorem~\ref{thm:hessian}]
By Lemma~\ref{lem:phi_scaling}, in expectation over random Kaiming
(variance-scaled Gaussian) initialization,
\[
    \lambda_{\max}(G^{(M)}_{\phip\phip}) \;\geq\; c_\phi \cdot M
\]
for a constant $c_\phi > 0$ bounded below independently of $M$: the witness
bound~\eqref{eq:witness_bound} gives
$\E[\lambda_{\max}] \geq (1 - \delta_h)\, q_h M_h \geq
(1 - \delta_h)\, q_h \beta M$, where $\beta > 0$ is the width-ratio floor
$M_h \geq \beta M$ supplied by the head hypothesis, so one may take
$c_\phi := (1 - \delta_h) q_h \beta$. The witness route delivers this with
no additive correction; within Karakida et
al.'s~\citep{karakida2019universal} proportional-width model class, their
Theorem~4 sharpens the width-linear rate to an exact asymptotic law. By
Lemma~\ref{lem:opcap},
\[
    \lambda_{\max}(G_{\thetap\thetap})
    \;\leq\; \|J_\thetap\|_{\mathrm{op}}^2 \;\leq\; C_\thetap
\]
with $C_\thetap$ independent of $M$ and $\Cg$. (If $C_\thetap = 0$ the
spectral block vanishes identically and the bound holds trivially under the
$+\infty$ convention of Definition~\ref{def:dcurv}; assume $C_\thetap > 0$
below.) The cap is deterministic---valid for every realization of the
initialization---so it divides through inside the expectation:
\begin{equation}
    \E\!\left[\Dcurv^{(M)}\right]
    \;=\; \E\!\left[\frac{\lambda_{\max}(G^{(M)}_{\phip\phip})}
                {\lambda_{\max}(G_{\thetap\thetap})}\right]
    \;\geq\; \frac{\E\!\left[\lambda_{\max}(G^{(M)}_{\phip\phip})\right]}
                  {C_\thetap}
    \;\geq\; \frac{c_\phi \cdot M}{C_\thetap}
    \;=\; c_1 \cdot M,
    \label{eq:disparity_bound}
\end{equation}
where $c_1 := c_\phi / C_\thetap$ depends only on the initialization scheme,
the architecture class of $g_\phip$, the classifier's input Lipschitz
constant $\widetilde{L}$, the bottleneck dimension $K$, the input spectral
dimension $S$, and the data distribution---exactly the dependencies claimed
in the theorem statement. The witness route needs no additive correction
($c_2 = 0$ suffices); the subtracted constant is retained in the statement
so that sharper asymptotic constants, which hold only for $M$ large, may
also be substituted. Under $\Cg = \Theta(M^2)$, $M = \Theta(\sqrt{\Cg})$ and
the ratio form $\Dcurv \geq c_1' \sqrt{\Cg / \Cf} - c_2$ follows on
multiplying and dividing by $\sqrt{\Cf}$.
\end{proof}

\begin{remark}[What Theorem~\ref{thm:hessian} does and does not claim]
\label{rem:thm1_scope}
Theorem~\ref{thm:hessian} is a scaling law for the curvature ratio between
two Gauss--Newton blocks; it does not claim that the full joint
Gauss--Newton matrix has a large condition number in the strict
$\lambda_{\max}/\lambda_{\min}$ sense (which is trivially infinite under
rank-deficiency), nor that joint gradient descent produces disparate
effective learning rates (the separate content of
Theorem~\ref{thm:twoscale}, under its own hypotheses). The scaling
$\Dcurv = \Omega(M)$ (equivalently $\Omega(\sqrt{\Cg / \Cf})$ for
$\Theta(M^2)$-parameter families) is a purely geometric statement about how
curvature concentrates in the downstream module as the spatial network
widens. In particular $\Dcurv$ concerns the \emph{top} of the spatial
block's spectrum, while the fitting rate $\mu$ of
Assumption~\ref{ass:residual} is governed by how fast the residual can be
driven down---a property of other parts of the spatial spectrum. The two are
linked only under spectral-shape assumptions we do not make.
\end{remark}

\subsection{Relation to classical ill-conditioning}
\label{app:classical}

Ill-conditioned optimization has classical causes with classical remedies,
and it is natural to ask whether they remove the width scaling of
Theorem~\ref{thm:hessian}. Each acts on a different lever, and none acts on
the one the theorem uses.

\paragraph{Feature normalization.}
Scale asymmetry or collinearity among input features degrades the
conditioning of $\Sigma_X$ and is classically cured by standardizing (e.g.\
mean-centering) the inputs. Here the cure has a sharper, quantitative role:
Lemma~\ref{lem:opcap} places $\lambda_{\max}(\Sigma_X)$ in the spectral cap,
so input transformations that deflate the top of the input spectrum
\emph{raise} the bound on the disparity at every width---normalization moves
the constants. What no input transformation can touch is the width scaling
of the spatial block (Lemma~\ref{lem:phi_scaling}): capacity moves the rate.
On raw nonnegative spectra, whose second moment is dominated by the shared
mean shape, the cap is large and the disparity the bound permits
correspondingly small. Whether an input transformation moves the
\emph{realized} disparity, however, depends on the architecture between
$f_\thetap$ and the loss: internal normalization layers can absorb input
shifts entirely (their Jacobians annihilate constant directions), in which
case the effective regressor is already centered and the lever acts through
the normalization choice, not the preprocessing. Section~\ref{sec:diagnostics}
measures exactly this on the production architecture.

\paragraph{Weight decay.}
Adding $\alpha \norm{w}^2$ to the loss shifts the Hessian by $2\alpha I$,
lifting every eigenvalue by $2\alpha$ and famously repairing classical
condition numbers. But the shift acts on \emph{both} blocks equally: the
disparity becomes $(\lambda_{\max}(G_{\phip\phip}) + 2\alpha) /
(\lambda_{\max}(G_{\thetap\thetap}) + 2\alpha)$, which remains $\Theta(M)$
for any fixed $\alpha$ as $M$ grows. Removing the disparity this way would
require $2\alpha \sim \lambda_{\max}(G_{\phip\phip})$, a regularization
strength that would dominate the data term and prevent fitting. The same
arithmetic applies to any penalty that shifts the two blocks by a common
multiple of the identity, including logit penalties of the Spectral
Decoupling type~\citep{pezeshki2021gradient}.

\paragraph{Adaptive optimizers.}
Adam and RMSProp precondition per-coordinate gradient magnitudes with a
diagonal rescaling. The raw block spectra are optimizer-invariant, but the
\emph{dynamics} under a diagonal preconditioner $P$ are governed by the
preconditioned spectra of $P^{-1/2} G P^{-1/2}$ and need not respect the
disparity: under Adam, per-coordinate steps are near-constant in gradient
magnitude, so small spectral gradients do not imply small spectral movement.
The gradient-flow statements around Theorem~\ref{thm:twoscale} are
accordingly scoped to (stochastic) gradient descent, with the discrete
pathwise versions of Appendix~\ref{app:math}; joint-versus-frozen gaps
observed under Adam are evidence about the phenomenon, not applications of
the theorem. The synthetic dynamics experiment reported in
Section~\ref{sec:diagnostics} and Appendix~\ref{app:synthetic} is run under
Adam, and at three seeds its arm ordering is not resolved; we treat the
stronger claim---that no shared-learning-rate adaptive method can remove the
disparity's dynamical consequence---as an open question rather than a
theorem. Optimizers with explicitly different learning rates per module
(two-timescale update rules; Heusel et al.~\citep{heusel2017gans};
Borkar~\citep{borkar2008stochastic}, Ch.~6) are the natural remedy at the
optimizer level; we do not evaluate them here, and
Section~\ref{sec:matched} is our reason for not recommending untuned
comparisons.

\begin{remark}[Parameterization scope]
\label{rem:param_scope}
$\Dcurv$ is not invariant to per-module reparameterization, and
Theorem~\ref{thm:hessian} is a statement about the standard variance-scaled
(Kaiming) parameterization---the one used in the pipelines we analyze and in
all our experiments. Under alternative parameterizations that rescale layers
with width, such as $\mu$P~\citep{yang2021tensor}, both the initialization
magnitudes and the effective per-layer learning rates change, and the width
scaling of the blocks would need to be re-derived; we do not claim the
disparity law transfers. Our uses of $\mu$P elsewhere are limited to a
different role: as one route to keeping the input-Jacobian bound
(Assumption~\ref{ass:inputlip}) uniform in width during training. Whether
principled reparameterization can neutralize the disparity---rather than
merely re-expressing it---is open. A closely related caution applies to
normalization layers, whose gain can move measured curvature with width at
fixed function (Appendix~\ref{app:math} and Section~\ref{sec:diagnostics};
cf.\ van Laarhoven~\citep{vanlaarhoven2017l2}).
\end{remark}

\begin{remark}[Defense against linear triviality]
\label{rem:not_pca}
A natural objection is that, since $f_\thetap$ is linear, the analysis must
reduce to principal component analysis (PCA) of the input. This is not the
case. PCA finds directions of maximum input variance \emph{without reference
to labels}. At initialization, however, the distinction is limited, and we
state this candidly: the Gauss--Newton blocks carry no label information
there (the softmax weights are near-uniform), and $G_{\thetap\thetap}$ is,
by Lemma~\ref{lem:opcap}, an input-second-moment object propagated through a
label-free random network. The defense against PCA-triviality is therefore a
statement about \emph{training}: along the trajectory the supervised
residual shapes $\partial \hat{y}/\partial Z$ and hence the curvature, and
the conditional pathology we identify is precisely that this supervised
signal may fail to reach $\thetap$---not that $f_\thetap$ is incapable of
representing useful features.
\end{remark}

\subsection{Scope and open extensions for Theorem~\ref{thm:twoscale}}
\label{app:thm2_open}

\begin{remark}[Open extensions]
\label{rem:thm2_open}
The proof of Theorem~\ref{thm:twoscale} is complete for the gradient-flow
model: the displacement bound follows by direct integration of the residual
envelope, with the operator-norm cap supplied by Lemma~\ref{lem:opcap}. Two
extensions are deliberately left open rather than sketched, and one is
supplied elsewhere. (i)~\emph{Discrete updates}: pathwise discrete
counterparts---gradient descent, momentum with clipping in either order,
Adam with a positive stabilizer, and decoupled weight decay---are proved in
Appendix~\ref{app:math} with $\mu_s = \eta \mu_f$, under the same
containment and gradient-bound hypotheses imposed at \emph{each} iterate; a
stochastic-approximation route that instead tracks the flow under gradient
noise (Borkar~\citep{borkar2008stochastic}, Ch.~2) would relax that
requirement, and we have not carried it out. (ii)~\emph{Deriving the
envelope}: obtaining Assumption~\ref{ass:residual}---in particular the
growth of $\mu$ with $\Cg$---from the architecture, for instance by a
kernel analysis of the spatial pathway at large width adapting Pezeshki et
al.~\citep{pezeshki2021gradient}, is the missing link that would chain
Theorems~\ref{thm:hessian} and~\ref{thm:twoscale} into a single causal
statement; we test its content empirically rather than claiming it, and the
residual-excitation and phase-restricted results of
Appendix~\ref{app:math} record why curvature anisotropy alone does not
supply it. (iii)~\emph{Envelope verification}: fitting $C/(1+\mu t)$
envelopes to measured residual trajectories, with the fitted $\hat{\mu}$
increasing in capacity, is an empirical programme, reported for the shallow
model in Appendix~\ref{app:synthetic}.
\end{remark}

\section{Additional theory}
\label{app:math}

This appendix collects the mathematical results that the main text cites but
does not prove. In order: pathwise discrete-time counterparts of the
displacement bound (gradient descent, momentum under either clipping order,
velocity clipping, Adam with a positive stabilizer, decoupled decay);
monotone and affine joint-versus-frozen comparisons, the finite-horizon
attribution bound that Corollary~\ref{cor:attribution} states in the main
text, the counterexample showing that a large \emph{top} kernel eigenvalue
is not a rate certificate, and a finite-data initialization certificate;
interior curvature witnesses and the normalization propositions, including
the conditional gain law under which a batch-normalized block's measured
curvature moves with width at fixed function; the residual-excitation
theorem and the phase-restricted attribution bound with its width
corollary; the solvable serial cross-entropy instance with its isotropic
extension, its displacement and margin-gap bounds and its normalized-readout
control; and a scoped nonlinear existence result for the shallow
biased-ReLU model under squared loss. It is reproduced verbatim from a
self-contained mathematical note, so its internal labels are namespaced and
its normalization is restated in its first paragraph; in particular
$K_b = J_b J_b^\top$ there is an output-space kernel and is \emph{not} the
parameter-space block $G_{bb}$, and the two must not be substituted for one
another. Its section headings are typeset as subsections of this appendix.

% appendix_math.tex is reproduced verbatim and is NOT edited here; its own
% \section headings are demoted locally so that the whole file sits inside
% Appendix~\ref{app:math}. The grouping keeps the redefinition local.
\begingroup
\let\section\subsection
\let\subsection\subsubsection
% Mathematics appendix, Astra, 2026-09-10.
% Integration: input AFTER \appendix. Requires the macros/environments and
% amsmath,amssymb,amsthm,mathtools,natbib of paper/main.tex.
% All local labels use astra_ to avoid collisions. Main-text cor:attribution
% should state the sharp result proved in Corollary cor:astra_attribution.
% Do not input the old and new proofs as duplicate theorem statements.
% BibTeX records for every citation used here are in the final comment block.
% Review-note numbers (1.1, 3.2, N, etc.) are mnemonic names, not hard-coded
% manuscript counters. The existing main-text labels remain unchanged.

\section{Mathematical scope and normalization}
\label{sec:astra_math}
We use the parameter blocks $\thetap,\phip$ of the main text and
$z=N^{-1/2}\operatorname{stack}(\hat y_n)$, $r=\nabla_z\ell(z)$,
$J_b=\partial z/\partial b$, $\loss=\ell\circ z$. Write
\[
 K_b=J_bJ_b^\top,\qquad G_{bb}=J_b^\top H_\tau J_b,
 \qquad \nabla_b\loss=J_b^\top r.
\]
Thus $K_b$ is an output-space logit kernel; it is not the parameter-space
GGN block $G_{bb}$. For averaged softmax CE,
$r=N^{-1/2}\operatorname{stack}(p_n-y_n)$ and $H_\tau$ is the block-diagonal
softmax Hessian. For normalized squared loss, $r=z-y$ and $H_\tau=I$.
All parameter norms are Euclidean (Frobenius for matrix parameters);
matrix-map bounds are operator norms unless explicitly marked otherwise.
Here $C=N_{\mathrm{cls}}$ denotes output dimension; envelope amplitudes
are written $C_r$ to avoid a collision. A normalized convention removes
extraneous powers of $N$, but does not make data-dependent constants
such as kernel minimum eigenvalues independent of $N$.

Theorem~\ref{thm:hessian} and Lemma~\ref{lem:phi_scaling} are geometric
statements. Theorem~\ref{thm:twoscale} additionally requires a trajectory
Jacobian cap and a residual envelope. Neither a large
$\lambda_{\max}(G_{\phip\phip})$ nor parameter count alone supplies
residual-subspace coercivity or a comparison with a separately trained
frozen model. For unit gradient flow, $\dot z=-(K_\thetap+K_\phip)r$;
under CE, $\dot r=-H_\tau(K_\thetap+K_\phip)r$, not the squared-loss
constant-kernel residual equation.

\section{Discrete displacement, clipping, and optimizer conventions}
\label{sec:astra_discrete}
The following seven statements collect the bounds numbered (4.1)--(4.7)
in the review notes and distinguish the two clipping algorithms. They are
pathwise: no convergence, stochastic independence, or smoothness assumption
is needed beyond the stated gradient bound at every iterate. There is no
encoder weight decay unless explicitly included. A uniform bound
$\|g_k\|\le B_kR_k\le BR_k$ derived from the main theorem requires its
containment/Jacobian hypothesis at \emph{each discrete iterate}. A full-data
or smoothed residual envelope does not establish a sampled-gradient envelope.
An observed residual floor limits a fitted envelope; a positive floor is
not a universal property of mini-batch training.

\begin{lemma}[Pathwise GD and clipping before a GD update]
\label{lem:astra_gd}
Let $K\ge0$, $\eta_k\ge0$, $a_k\in[0,1]$, and
$\|g_k\|\le B_kR_k$, where $B_k,R_k\ge0$. If
$\thetap_{k+1}=\thetap_k-\eta_ka_kg_k$, then
\begin{equation}\label{eq:astra_gd}
 \|\thetap_K-\thetap_0\|
 \le\sum_{k=0}^{K-1}\eta_ka_kB_kR_k
 \le\sum_{k=0}^{K-1}\eta_kB_kR_k.
\end{equation}
The clipping scalar may depend on all parameter groups and the history.
\end{lemma}
\begin{proof}
Telescope the updates and apply the triangle inequality.
\end{proof}

\begin{corollary}[Polynomial envelope in step units]
\label{cor:astra_steps}
In Lemma~\ref{lem:astra_gd}, suppose $\eta_k=\eta\ge0$, $B_k\le B$,
and $R_k\le C_r/(1+\mu_s k)$, with $C_r,\mu_s>0$.
For $K\ge1$,
\begin{equation}\label{eq:astra_steps}
 \|\thetap_K-\thetap_0\|
 \le\eta BC_r\left[1+\frac{\log(1+\mu_s(K-1))}{\mu_s}\right].
\end{equation}
For $0<\delta<C_r$, suppose the envelope holds through
$K_* =\lceil(C_r/\delta-1)/\mu_s\rceil$. Then the first index
$K_\delta=\min\{k:R_k\le\delta\}$ satisfies $K_\delta\le K_*$ and
\begin{equation}\label{eq:astra_hit}
 \|\thetap_{K_\delta}-\thetap_0\|
 \le\eta BC_r[1+\mu_s^{-1}\log(C_r/\delta)].
\end{equation}
The displacement is zero if $K_\delta=0$.
\end{corollary}
\begin{proof}
For decreasing $f(x)=(1+\mu_sx)^{-1}$,
$\sum_{k=1}^{K-1}f(k)\le\int_0^{K-1}f(x)\,dx$.
The envelope gives $R_{K_*}\le\delta$. If $K_\delta\ge1$,
$K_\delta-1\le K_*-1<(C_r/\delta-1)/\mu_s$; substitute in
\eqref{eq:astra_steps}.
\end{proof}
\begin{remark}[Flow-to-step dictionary]
For constant positive $\eta$, $t_k=\eta k$ and $\mu_s=\eta\mu_f$.
Consequently $\eta BC_r/\mu_s=BC_r/\mu_f$. Omitting the leading
$\eta$ in \eqref{eq:astra_steps} is a units error. The extra
$\eta BC_r$ is a left-sum allowance, not an additional mechanism.
\end{remark}

\begin{lemma}[Variable-step envelope]
\label{lem:astra_variable}
Under Lemma~\ref{lem:astra_gd}, let $t_k=\sum_{i<k}\eta_i$,
$B_k\le B$, $R_k\le C_r/(1+\mu_ft_k)$ with $C_r,\mu_f>0$.
For $K\ge1$ and $\eta_{\max}=\max_{k<K}\eta_k$,
\begin{equation}\label{eq:astra_variable}
 \|\thetap_K-\thetap_0\|\le BC_r\left[
 \frac{\log(1+\mu_ft_K)}{\mu_f}
 +\eta_{\max}\left(1-\frac1{1+\mu_ft_K}\right)\right].
\end{equation}
\end{lemma}
\begin{proof}
The excess of the left sum for $f(t)=(1+\mu_ft)^{-1}$ over its
integral on $[t_k,t_{k+1}]$ is at most
$\eta_k(f(t_k)-f(t_{k+1}))$. Bound $\eta_k$ by $\eta_{\max}$ and
telescope. Zero step sizes cause no difficulty.
\end{proof}

\begin{lemma}[Momentum, including gradient clipping before accumulation]
\label{lem:astra_momentum}
Let $0\le\beta<1$, $\eta_k\ge0$, $a_k\in[0,1]$, and
$\|g_k\|\le B_kR_k$. For
$v_{k+1}=\beta v_k+a_kg_k$ and
$\thetap_{k+1}=\thetap_k-\eta_kv_{k+1}$,
\begin{align}
 \thetap_K-\thetap_0
 &=-v_0\sum_{k<K}\eta_k\beta^{k+1}
   -\sum_{i<K}a_ig_i\sum_{k=i}^{K-1}\eta_k\beta^{k-i}.
 \label{eq:astra_momentum_exact}
\end{align}
For constant $\eta$ this implies
\begin{equation}\label{eq:astra_momentum}
 \|\thetap_K-\thetap_0\|\le
 \frac{\eta\beta(1-\beta^K)}{1-\beta}\|v_0\|
 +\eta\sum_{i<K}\frac{1-\beta^{K-i}}{1-\beta}a_iB_iR_i.
\end{equation}
For $v_0=0$, \eqref{eq:astra_steps}--\eqref{eq:astra_hit} therefore
remain valid with an additional factor $1/(1-\beta)$.
\end{lemma}
\begin{proof}
Induction gives $v_{k+1}=\beta^{k+1}v_0+
\sum_{i=0}^k\beta^{k-i}a_ig_i$. Substitute into the parameter update,
interchange the finite sums, and sum the geometric series. For variable
steps the inner coefficient is at most $\eta_{\max}/(1-\beta)$.
\end{proof}
\begin{remark}
This is zero-dampening, non-Nesterov momentum. It includes global
\emph{gradient} clipping through its actual $a_k$; the gradients and
later iterates belong to the clipped run itself. It does not compare
a clipped run's displacement with a separately evolved unclipped run.
The factor $1/(1-\beta)$ is approached by long sequences of aligned gradients.
\end{remark}

\begin{lemma}[Velocity clipping: sharper norm induction]
\label{lem:astra_velocity}
Let $c_k\ge0$, $0\le\beta<1$, and let
$\operatorname{clip}_{c}(u)=u\min\{1,c/\|u\|\}$, with
$\operatorname{clip}_{c}(0)=0$. Suppose
\[
 v_{k+1}=\operatorname{clip}_{c_k}(\beta v_k+g_k),\qquad
 \thetap_{k+1}=\thetap_k-\eta_kv_{k+1},\quad \eta_k\ge0.
\]
Define the \emph{comparison} coefficient
$a_k=\min\{1,c_k/\|g_k\|\}$ for $g_k\ne0$ and $a_k=1$ otherwise.
Then
\[
 \|v_{k+1}\|\le\beta^{k+1}\|v_0\|
       +\sum_{i=0}^k\beta^{k-i}a_i\|g_i\|.
\]
Consequently \eqref{eq:astra_momentum} holds when
$\|g_i\|\le B_iR_i$. For variable steps the corresponding norm bound
uses the same finite weights as \eqref{eq:astra_momentum_exact}.
The vector identity \eqref{eq:astra_momentum_exact} is not asserted.
\end{lemma}
\begin{proof}
Write $x=\beta\|v_k\|$, $y=\|g_k\|$. Then
\[
 \|v_{k+1}\|\le\min(c_k,x+y)\le x+\min(c_k,y)
 =\beta\|v_k\|+a_k\|g_k\|.
\]
If $y\le c_k$ use the second argument of the minimum; if $y>c_k$ use
its first. Induct and sum the parameter-step norms.
\end{proof}
\begin{remark}[Two valid coefficient conventions]
The comparison $a_k$ above differs from the actual multiplier
$t_k=\min\{1,c_k/\|\beta v_k+g_k\|\}$, defined as one at a zero
argument. The same norm bound is also valid with $t_i$ in place of
$a_i$, by the separate induction
$\|v_{k+1}\|\le\beta\|v_k\|+t_k\|g_k\|$.
In fact the exact unrolling is
\[
 v_{k+1}=\beta^{k+1}\left(\prod_{j=0}^kt_j\right)v_0
 +\sum_{i=0}^k\beta^{k-i}\left(\prod_{j=i}^kt_j\right)g_i.
\]
Dropping subsequent multipliers gives the claimed $t_i$ bound.
The two coefficient conventions need not be ordered, owing to gradient
cancellation. State which one is used. For radial implementations with
$c/(\|u\|+\varepsilon_{\rm clip})$ and clipping multiplier in $[0,1]$,
this actual-multiplier argument still applies; the ideal-clipping
comparison $a_i$ is not obtained by silently equating the two formulas.
\end{remark}

\begin{lemma}[Positive-stabilizer Adam]
\label{lem:astra_adam}
Suppose the first moment is initialized at zero, $0\le\beta_1<1$,
$0\le\eta_k\le\eta_{\max}$, and
\[
 \widehat m_k=\sum_{i=0}^k
 \frac{(1-\beta_1)\beta_1^{k-i}}{1-\beta_1^{k+1}}g_i,
 \qquad \thetap_{k+1}-\thetap_k=-\eta_kD_k\widehat m_k,
 \qquad \|D_k\|\le\zeta^{-1},\quad\zeta>0.
\]
Then
\begin{equation}\label{eq:astra_adam}
 \|\thetap_K-\thetap_0\|
 \le\frac{\eta_{\max}}{\zeta(1-\beta_1)}\sum_{i<K}\|g_i\|.
\end{equation}
Standard Adam denominators $\sqrt{\widehat v_k}+\zeta$ satisfy the
preconditioner bound when $\widehat v_k\ge0$ coordinatewise. If
clipping occurs before moment accumulation, $g_i$ here is the gradient
\emph{after} that clipping.
\end{lemma}
\begin{proof}
Bound the update by $\eta_{\max}\zeta^{-1}\|\widehat m_k\|$ and
interchange sums. For each $i$, the sum of its nonnegative weights is
at most $\sum_{k\ge i}\beta_1^{k-i}\le(1-\beta_1)^{-1}$, since
$(1-\beta_1)/(1-\beta_1^{k+1})\le1$.
\end{proof}
\begin{remark}[Numerically weak optimizer constant]
The factor $1/\zeta$ makes this certificate numerically vacuous at
standard small stabilizers unless a much smaller preconditioner bound
is independently certified. At $\zeta=0$, the first scalar Adam step
has magnitude $\eta$ for every nonzero $g_0$, so no constant uniform
in residual scale bounds it by a multiple of $\eta\|g_0\|$.
\end{remark}

\begin{lemma}[Decoupled encoder decay must be included]
\label{lem:astra_decay}
If an update is $\thetap_{k+1}-\thetap_k=-u_k-\eta_k\lambda_k\thetap_k$
with $\eta_k,\lambda_k\ge0$, then
$\|\thetap_K-\thetap_0\|\le\sum_{k<K}\|u_k\|+
\sum_{k<K}\eta_k\lambda_k\|\thetap_k\|$.
Even with $u_k=0$, constant decay $\lambda$ gives
$\thetap_K=\prod_{k<K}(1-\eta_k\lambda)\thetap_0$.
\end{lemma}
\begin{proof}
Telescope for the bound and iterate the scalar multiplication for the
identity. A nonzero decay displacement can occur at identically zero
residual, so a residual-only bound cannot omit this term.
\end{proof}

\section{Joint and frozen trajectory comparisons}
\label{sec:astra_bridges}
\begin{theorem}[Monotone frozen flow; review Theorem 1.1]
\label{thm:astra_monotone}
On $[0,T]$ let the joint and frozen paths be absolutely continuous,
with common spatial initialization and
$\dot\phip_J=-\nabla_\phip\loss(\thetap_J,\phip_J)$,
$\dot\phip_F=-\nabla_\phip\loss(\thetap_0,\phip_F)$ a.e.
Let $d\ge0$ be locally integrable with
$\|\thetap_J(t)-\thetap_0\|\le d(t)$.
Assume the differentiable frozen objective has an $\alpha$-monotone
gradient on a region containing both spatial paths, $\alpha\ge0$:
\[
 \langle v-w,\nabla_\phip\loss(\thetap_0,v)
 -\nabla_\phip\loss(\thetap_0,w)\rangle\ge\alpha\|v-w\|^2.
\]
Assume finite nonnegative constants $K_{\theta\phi},L_\theta,L_\phi$
satisfy, on the parameter pairs used below,
\begin{align*}
 \|\nabla_\phip\loss(\thetap_J,\phip_J)
       -\nabla_\phip\loss(\thetap_0,\phip_J)\|&\le K_{\theta\phi}d(t),\\
 \|z(\thetap_J,\phip_J)-z(\thetap_0,\phip_J)\|&\le L_\theta d(t),\\
 \|z(\thetap_0,\phip_J)-z(\thetap_0,\phip_F)\|&\le L_\phi\|\phip_J-\phip_F\|.
\end{align*}
Then
\begin{align}
 \|\phip_J(t)-\phip_F(t)\|&\le K_{\theta\phi}\int_0^t e^{-\alpha(t-s)}d(s)\,ds,
 \label{eq:astra_monotone}\\
 \|z_J(t)-z_F(t)\|&\le L_\theta d(t)+
 L_\phi K_{\theta\phi}\int_0^t e^{-\alpha(t-s)}d(s)\,ds.
\end{align}
If $d$ is nondecreasing, the last bound is
$[L_\theta+L_\phi K_{\theta\phi}\psi_\alpha(t)]d(t)$, where
$\psi_\alpha(t)=(1-e^{-\alpha t})/\alpha$ for $\alpha>0$ and
$\psi_0(t)=t$.
\end{theorem}
\begin{proof}
For $\Delta=\phip_J-\phip_F$, split the derivative into a frozen-gradient
difference and an encoder-induced forcing. Then
$\frac12(d/dt)\|\Delta\|^2\le-\alpha\|\Delta\|^2+
K_{\theta\phi}d\|\Delta\|$. The scalar norm comparison (or the
regularization $\sqrt{\|\Delta\|^2+\epsilon^2}$ followed by
$\epsilon\downarrow0$) proves \eqref{eq:astra_monotone}. Add and subtract
$z(\thetap_0,\phip_J)$ for the output bound; monotonicity of $d$ gives
the final expression.
\end{proof}
\begin{remark}
Under Theorem~\ref{thm:twoscale}, one may use
$d(t)=BC_r\log(1+\mu t)/\mu$ on its valid horizon. If the first fitting
time lies in that horizon, replace this logarithm by $\log(C_r/\delta)$
only up to that fitting time. The constants $L_\theta,L_\phi,K_{\theta\phi}$
and $\alpha$ are additional path-dependent hypotheses; all must be
controlled across widths before inferring a width rate. A PL inequality
alone supplies none of this monotonicity.
\end{remark}

\begin{proposition}[Counterexample 1.2: PL does not imply contraction]
\label{prop:astra_pl_counter}
For $a,b>0$, $f(x,y)=\frac12(x+ay^2)^2$ satisfies
$\frac12\|\nabla f\|^2\ge f-\inf f$ globally, but nearby gradient-flow
trajectories can expand by an arbitrarily large factor at time one.
\end{proposition}
\begin{proof}
$\inf f=0$ and $\|\nabla f\|^2=(x+ay^2)^2(1+4a^2y^2)\ge2f$.
Along $(x(t),y(t))=(-be^{-t},0)$, transverse linearized variation satisfies
$\dot\xi=2abe^{-t}\xi$, hence
$\xi(t)=\xi(0)\exp(2ab(1-e^{-t}))$. Vary $a$ with the PL constant fixed.
This uses the PL convention of \citet{astra_karimi2016}.
\end{proof}

\begin{theorem}[Affine squared-loss bridge; review Theorem 1.3]
\label{thm:astra_affine}
Let $z=z_0+A(\thetap-\thetap_0)+B(\phip-\phip_0)$ and
$\loss=\frac12\|z-y\|^2$. Train the joint model and the model with
$\thetap=\thetap_0$ by unit Euclidean gradient flow from the same
initialization. Let $K_\thetap=AA^\top$, $K_\phip=BB^\top$.
Suppose a common invariant linear subspace $\mathcal V$ contains
$e_0=z_0-y$ and, on $\mathcal V$,
$K_\phip\succeq\kappa I$, $0\preceq K_\thetap\preceq aI$ with
$\kappa>0$, $a\ge0$. For $t\ge0$,
\begin{align}
 \|z_J(t)-z_F(t)\|&\le at e^{-\kappa t}\|e_0\|,
 &\sup_{t\ge0}\|z_J(t)-z_F(t)\|&\le\frac{a\|e_0\|}{e\kappa},
 \label{eq:astra_affine_gap}\\
 \|\thetap_J(t)-\thetap_0\|&\le
 \frac{\sqrt a}{\kappa}(1-e^{-\kappa t})\|e_0\|,
 &\int_0^\infty\|\nabla_\thetap\loss\|^2dt&\le
 \frac{a\|e_0\|^2}{2\kappa}.\label{eq:astra_affine_disp}
\end{align}
At every finite horizon with positive loss decrease the encoder share
is at most $a/(a+\kappa)$, by Corollary~\ref{cor:astra_attribution}.
\end{theorem}
\begin{proof}
On $\mathcal V$, both semigroups have norm at most $e^{-\kappa t}$.
Their difference obeys the noncommutative Duhamel identity
\[
 e_J(t)-e_F(t)=-\int_0^t e^{-K_\phip(t-s)}K_\thetap
                e^{-(K_\phip+K_\thetap)s}e_0\,ds.
\]
This gives $at e^{-\kappa t}\|e_0\|$, maximized at $t=1/\kappa$.
Since $\|A^\top v\|\le\sqrt a\|v\|$ for $v\in\mathcal V$,
integrate $\dot\thetap=-A^\top e_J$ and its squared norm to get
\eqref{eq:astra_affine_disp}. Both residuals converge to zero.
\end{proof}

\begin{corollary}[Sharp finite-horizon attribution; main Corollary~\ref{cor:attribution}]
\label{cor:astra_attribution}
Explicitly: $\loss=\ell\circ z$ is differentiable, the parameter path
is absolutely continuous, the residual is $r=\nabla_z\ell$ with
$K_\thetap=J_\thetap J_\thetap^\top$,
$K_\phip=J_\phip J_\phip^\top$, both blocks follow unit Euclidean
gradient flow, and the chain rule holds a.e. Suppose for constants
$a\ge0$, $\kappa>0$ that the Rayleigh inequalities
\[
 r^\top K_\thetap r\le a\|r\|^2,\qquad
 r^\top K_\phip r\ge\kappa\|r\|^2
\]
hold at almost every time along the realized trajectory on $[0,T]$,
and $\loss(0)-\loss(T)>0$. Then
\begin{equation}\label{eq:astra_attribution}
 \frac{\int_0^T\|\nabla_\thetap\loss\|^2dt}{\loss(0)-\loss(T)}
 \le\frac{a}{a+\kappa}.
\end{equation}
No common invariant subspace is required: that hypothesis of
Theorem~\ref{thm:astra_affine} is a device for guaranteeing the
Rayleigh bounds along the path, and here they are assumed directly.
\end{corollary}
\begin{proof}
Put $q_b=r^\top K_b r=\|\nabla_b\loss\|^2$.
Pointwise $\kappa q_\thetap\le a q_\phip$. Integrate and use
$\loss(0)-\loss(T)=\int_0^T(q_\thetap+q_\phip)dt$.
Scalar kernels $K_\thetap=a$, $K_\phip=\kappa$ attain equality.
\end{proof}
\begin{remark}
The same result holds for the nonsmooth flow convention defined below
whenever its a.e. energy identity and Rayleigh inequalities hold.
The unit rate is not cosmetic: with per-block rates $\eta_\theta\ne
\eta_\phi$, the numerator $\int\|\nabla_\thetap\loss\|^2dt$ is no
longer the encoder share. For positive constant rates, the actual share
uses $\eta_\theta$ in the numerator and is bounded by
$\eta_\theta a/(\eta_\theta a+\eta_\phi\kappa)$.
\end{remark}

\begin{proposition}[Counterexample 1.4: top eigenvalues are insufficient]
\label{prop:astra_top_counter}
Let $K_\phip=\operatorname{diag}(M,0)$,
$K_\thetap=\operatorname{diag}(0,1)$, $e_0=(0,1)^\top$, $M>0$.
In the affine squared-loss model, the joint/frozen logit gap is
$1-e^{-t}$ although $\lambda_{\max}(K_\phip)=M$.
\end{proposition}
\begin{proof}
$e_F(t)=(0,1)^\top$ and $e_J(t)=(0,e^{-t})^\top$. Each kernel is
realized by its PSD square root as an affine Jacobian. The fast spatial
direction contains none of the residual.
\end{proof}

\begin{proposition}[Robustness of the affine comparison]
\label{prop:astra_robust}
Under the kernel/subspace assumptions of Theorem~\ref{thm:astra_affine},
suppose instead the residuals solve
$\dot e_J=-(K_\phip+K_\thetap)e_J+\xi_J$ and
$\dot e_F=-K_\phip e_F+\xi_F$, with common $e_0$ and locally
integrable $\xi_J,\xi_F\in\mathcal V$. Then add
\[
 \int_0^t e^{-\kappa(t-s)}(\|\xi_J(s)\|+\|\xi_F(s)\|)\,ds
\]
to the first bound in \eqref{eq:astra_affine_gap}. If
$\|\xi_J\|\le\nu_J$ and $\|\xi_F\|\le\nu_F$, the additional
uniform term is at most $(\nu_J+\nu_F)/\kappa$.
\end{proposition}
\begin{proof}
Compare each perturbed solution with its own constant-kernel unperturbed
solution by variation of constants, then use the triangle inequality.
\end{proof}

\begin{proposition}[Logit margins imply a disagreement bound]
\label{prop:astra_margins}
If normalized logits satisfy $\|z_J-z_F\|\le h$ and at most a fraction
$b_\gamma$ of the N samples have frozen top-two unnormalized logit
margin at most $\gamma>0$, the disagreement fraction is at most
$b_\gamma+2h^2/\gamma^2$ (and at most one).
\end{proposition}
\begin{proof}
On a remaining disagreement, for the frozen winner $c_*$ and the joint
winner $c$, the logit perturbation $d_n$ satisfies
$d_{n,c}-d_{n,c_*}>\gamma$. Thus
$\|d_n\|^2\ge(d_{n,c}-d_{n,c_*})^2/2>\gamma^2/2$.
Average and apply Markov. Macro-F1 needs further class-specific
confusion and denominator control.
\end{proof}
\subsection{A finite-data initialization certificate}
\begin{corollary}[Biased-ReLU initialization; review Corollary 1.5]
\label{cor:astra_init}
Fix $N<\infty$, $x_n\in\R^S$, and a deterministic encoder matrix
$\Theta_0$ for which $u_n=\Theta_0x_n\in\R^K$ are pairwise distinct
and $\|u_n\|\le R$. Let
\[
 h_i(u)=\operatorname{ReLU}(w_i^\top u+b_i),\qquad \hat y=Vh+b_2,
\]
with independent $w_i\sim\mathcal N(0,s_w^2I_K/K)$,
$b_i\sim\mathcal N(0,s_b^2)$, $V_{ci}\sim\mathcal N(0,s_w^2/M)$,
$b_{2,c}\sim\mathcal N(0,s_b^2)$, and $s_w,s_b>0$, $1\le c\le C$.
Let $A,B$ be the normalized encoder/spatial logit Jacobians at this
initialization, including all hidden and readout parameters in the
spatial block. Train the resulting \emph{exactly affine} model with
squared loss as in Theorem~\ref{thm:astra_affine}.
Define
\[
 \psi=(h_1(u_n))_{n=1}^N,\quad K_*=N^{-1}\E\psi\psi^\top,
 \quad\kappa_*=\lambda_{\min}(K_*),\quad
 s_{\max}^2=s_w^2R^2/K+s_b^2.
\]
Then $\kappa_*>0$. For $0<\rho<1$, put
\begin{align*}
 M_0&=\left\lceil\frac{18s_{\max}^4}{\rho\kappa_*^2}\right\rceil,
 & A_\rho&=\frac{3C_G}{\rho},
 & R_\rho&=\sqrt{3C_r/\rho},\\
 C_G&=Cs_w^4\operatorname{tr}(\Sigma_X),
 & C_r&=2C(s_w^2s_{\max}^2/2+s_b^2)+2\|y\|^2,
 &\Sigma_X&=N^{-1}\sum_nx_nx_n^\top.
\end{align*}
For each $M\ge M_0$, with probability at least $1-\rho$,
\[
 K_\phip\succeq\kappa_*M I/2,\qquad \|K_\thetap\|\le A_\rho,
 \qquad\|e_0\|\le R_\rho.
\]
Consequently
\[
 \sup_t\|z_J-z_F\|\le\frac{2A_\rho R_\rho}{e\kappa_*M},
 \qquad\sup_t\|\thetap_J-\thetap_0\|
 \le\frac{2\sqrt{A_\rho}R_\rho}{\kappa_*M}.
\]
The displayed sufficient threshold necessarily satisfies
$M_0\ge72N^2/\rho$; it is an existence certificate, not a practical
onset width or a lower bound on widths at which coercivity actually holds.
\end{corollary}
\begin{proof}
If $c^\top K_*c=0$, the continuous function
$\sum_nc_n\operatorname{ReLU}(w^\top u_n+b)$ vanishes almost surely
under a Gaussian with full support, hence everywhere. The hyperplanes
$w^\top u_n+b=0$ have pairwise nonproportional normals $(u_n,1)$.
At a point of the nth hyperplane on none of the others, crossing the
hyperplane produces the gradient jump $c_n(u_n,1)$; all other terms
are affine locally. The jump must vanish, so each $c_n=0$.

Let $\mathsf H_{ni}=h_i(u_n)$ be the unnormalized feature matrix.
The readout kernel is $(\mathsf H\mathsf H^\top/N)\otimes I_C$.
Gaussian ReLU fourth moments and Cauchy--Schwarz give
$\E\|\psi\|^4\le3N^2s_{\max}^4/2$. Independence of columns yields
\[
 \E\left\|\frac{\mathsf H\mathsf H^\top}{NM}-K_*\right\|_F^2
 \le\frac{3s_{\max}^4}{2M},\qquad
 \Pr\left(\left\|\frac{\mathsf H\mathsf H^\top}{NM}-K_*\right\|>
                  \kappa_*/2\right)
 \le\frac{6s_{\max}^4}{M\kappa_*^2}.
\]
For the encoder Jacobian, conditioning on the hidden weights kills the
cross terms in the independent centered readouts. Bounding each gate by
one gives $\E\|J_\thetap\|_F^2\le C_G$. Also
$\E\|e_0\|^2\le C_r$ by the activation second moment and
$\|z_0-y\|^2\le2\|z_0\|^2+2\|y\|^2$.
Markov and a three-event union bound prove the asserted event.
Apply Theorem~\ref{thm:astra_affine} on the full output space.
Finally $\operatorname{tr}(K_*)\le s_{\max}^2/2$ implies
$\kappa_*\le s_{\max}^2/(2N)$ and the threshold inequality.
\end{proof}
\begin{remark}
This proves a frozen-linearization instance only. Constants depend on
the fixed dataset and encoder; nearly duplicate bottleneck inputs can
make $\kappa_*$ arbitrarily small. Exact duplicates require a compatible
restricted-space formulation. The displayed constants are numerically
vacuous for the reported large datasets and are not empirical width thresholds.
\end{remark}

\begin{proposition}[Truncated-Chernoff alternative]
\label{prop:astra_chernoff}
For the same Gaussian/ReLU feature law, let $X_i=\psi_i\psi_i^\top/N$.
Choose $R_c>0$ such that
\[
 2N(R_c+2s_{\max}^2)e^{-R_c/(2s_{\max}^2)}\le\kappa_*/2.
\]
Then
\begin{equation}\label{eq:astra_chernoff}
 \Pr\left\{\lambda_{\min}\left(\sum_{i=1}^M X_i\right)
               <M\kappa_*/4\right\}
 \le N\exp[-M\kappa_*/(16R_c)].
\end{equation}
In particular $M\ge(16R_c/\kappa_*)\log(N/\eta)$ suffices for failure
probability $\eta\in(0,1)$. If used in the previous corollary or the
nonlinear theorem below, replace the head floor $\kappa_*M/2$ by
$\kappa_*M/4$ and propagate this changed constant.
\end{proposition}
\begin{proof}
Let $Z$ be the maximum of the N squared Gaussian preactivations.
Then $\|X_i\|\le Z$ and $\Pr(Z>u)\le2N e^{-u/(2s_{\max}^2)}$.
Integrating the tail gives
$\E[\|X_i\|1_{\{\|X_i\|>R_c\}}]
\le2N(R_c+2s_{\max}^2)e^{-R_c/(2s_{\max}^2)}$.
Thus $Y_i=X_i1_{\{\|X_i\|\le R_c\}}$ satisfies
$0\preceq Y_i\preceq R_cI$ and $\E Y_i\succeq\kappa_*I/2$.
Since each $X_i$ is positive semidefinite,
$\sum_iX_i\succeq\sum_iX_i1_{\{\|X_i\|\le R_c\}}$ and therefore
$\lambda_{\min}(\sum_iX_i)\ge\lambda_{\min}(\sum_iY_i)$; the
lower-tail bound for the truncated sum applies verbatim to the
untruncated sum. Apply matrix Chernoff \citep{astra_tropp2012} with
$\delta=1/2$ and $\mu_{\min}\ge M\kappa_*/2$.
\end{proof}
\begin{remark}
The constant is Tropp's $\delta=1/2$ lower tail after the standard
simplification
$[e^{-\delta}(1-\delta)^{-(1-\delta)}]^{\mu_{\min}/R_c}
\le e^{-\delta^2\mu_{\min}/(2R_c)}$.
Keeping the exact constant $-\log(e^{-1/2}\sqrt2)=0.15343\ldots$
replaces $1/16$ by approximately $1/13.04$; the difference is immaterial
at the widths involved. The one-sided ReLU tail also permits $N$ in
place of $2N$ in the truncation estimate. This certificate removes a
power of $\kappa_*^{-1}$ at the cost of logarithmic truncation factors;
it can still be numerically vacuous. Neither its threshold nor the
Frobenius threshold is a necessity theorem.
\end{remark}

\section{Interior curvature witnesses and normalization}
\label{sec:astra_interior}
\begin{lemma}[Pointwise interior witness; review Lemma 3.1]
\label{lem:astra_interior}
At N sites let $h_n\in\R^d$, $a_n=Wh_n+b\in\R^q$, and
$\hat y_n=\Phi_n(a_n)\in\R^C$, $C\ge2$.
Perturbing $a_n$ must not change logits at other sites in this lemma.
Fix a parameter point with finite logits and existing downstream
Jacobians $D_n=\partial\Phi_n/\partial a_n$. Let
$\Pi=I-\boldsymbol1\boldsymbol1^\top/C$,
$p_{\min,n}=\min_c p_{n,c}$. For a unit $a\in\R^q$ define
\[
 S_{h,a}=N^{-1}\sum_n p_{\min,n}\|\Pi D_na\|^2h_nh_n^\top,
 \qquad S_h^p=N^{-1}\sum_n p_{\min,n}h_nh_n^\top.
\]
If W belongs to $\phip$, then
\[
 \lambda_{\max}(G_{\phip\phip})\ge\lambda_{\max}(G_{WW})
 \ge\lambda_{\max}(S_{h,a}).
\]
If additionally $\|\Pi D_na\|\ge s$ for every n with $h_n\ne0$,
then the final lower bound is at least $s^2\lambda_{\max}(S_h^p)$.
\end{lemma}
\begin{proof}
For each unit $u\in\R^d$, $\Delta W=au^\top$ has unit Frobenius norm.
The CE Hessian satisfies $H_n\succeq p_{\min,n}\Pi$: indeed
$v^\top H_nv=\min_c\sum_jp_{n,j}(v_j-c)^2
\ge p_{\min,n}\min_c\sum_j(v_j-c)^2$.
Consequently the witness Rayleigh quotient is at least
$N^{-1}\sum_n(h_n^\top u)^2p_{\min,n}\|\Pi D_na\|^2
=u^\top S_{h,a}u$. Maximize over u. Zero-padding the W perturbation
proves the full-block inequality. The additional hypothesis gives
$S_{h,a}\succeq s^2S_h^p$.
\end{proof}
\begin{remark}[Mixing, masks, and nonsmooth points]
For a general interior perturbation, let $\mathcal R_h\operatorname{vec}
(\Delta W)$ be its full unnormalized preactivation stack, including all
sites that can affect the loss. If $D_\Phi$ is the full downstream
Jacobian and $P_{\rm out}$ selects the supervised logits, the valid
replacement is
\begin{equation}\label{eq:astra_full_witness}
 \lambda_{\max}(G_{\phip\phip})\ge
 \frac1{N_{\rm out}}\|H_\tau^{1/2}P_{\rm out}D_\Phi
       \mathcal R_h\operatorname{vec}(\Delta W)\|^2,
 \qquad\|\Delta W\|_F=1.
\end{equation}
This follows by the chain rule and the Rayleigh principle and includes
BN, convolutions, and ignored sites. A product lower bound requires a
restricted downstream gain on this same witness, with matching site
normalization; a smallest nonzero singular value does not protect a
witness in a nullspace. At ReLU kinks one must specify a gate selection;
for nondegenerate Gaussian preactivations the ordinary Jacobian exists
a.s. Sites with zero feature and zero bias contribute zero to the
incoming lower bound, but a Jacobian convention is still necessary.
\end{remark}

\begin{theorem}[Fixed bottleneck, growing interior width; review Theorem 3.2]
\label{thm:astra_interior_relu}
Fix integers $q\ge1$, $C\ge2$, and constants $R,a_0,s_w,s_v>0$,
$s_b\ge0$. For each M let deterministic features $h_n\in\R^M$ satisfy
$\|h_n\|^2\le R^2M$ and
$\lambda_{\max}(S_h)\ge a_0M$, where
$S_h=N^{-1}\sum_nh_nh_n^\top$.
Consider $\hat y_n=V\operatorname{ReLU}(Wh_n+b)$ with independent
$W_{ij}\sim\mathcal N(0,s_w^2/M)$, $b_i\sim\mathcal N(0,s_b^2)$,
$V_{ci}\sim\mathcal N(0,s_v^2/q)$, independent of the features.
At any identically zero preactivation fix a gate in $[0,1]$ consistently
in all Jacobian blocks; such sites have $h_n=0$ and cannot contribute
to the W witness below. Under averaged CE in these Euclidean
coordinates there exist $c_0,p_0>0$
independent of M,N with
\[
 \Pr\{\lambda_{\max}(G_{WW})\ge c_0a_0M\}\ge p_0,
 \qquad \E\lambda_{\max}(G_{\phip\phip})\ge p_0c_0a_0M
\]
whenever W belongs to $\phip$.
One explicit choice is
\[
 s_{\max}^2=s_w^2R^2+s_b^2,\quad L_0=\sqrt{48q s_{\max}^2},
 \quad c_0=e^{-4L_0}/(32C),\quad p_0=p_V/4,
\]
where $p_V=\Pr\{\|V\|_F\le2,\ \|\Pi Ve_1\|\ge1/2\}>0$.
These constants are \emph{existence-only and numerically vacuous}; this
is a positive-probability/expectation statement, not a probability
approaching one as M grows.
\end{theorem}
\begin{proof}
Fix a top unit eigenvector u of $S_h$, write
$a_n=(h_n^\top u)^2$, $A=N^{-1}\sum_na_n\ge a_0M$, and
$\chi_n=1_{\{(Wh_n+b)_1>0\}}$.
At sites with $a_n>0$ this preactivation is a nondegenerate centered
Gaussian. For $Q=N^{-1}\sum_na_n\chi_n$, $\E Q=A/2$ and $0\le Q\le A$,
hence $\Pr(Q\ge A/4)\ge1/3$ by splitting the expectation.
The exact moment bound
$\E\|\operatorname{ReLU}(Wh_n+b)\|^2\le qs_{\max}^2/2$
gives $\Pr(\|\operatorname{ReLU}(Wh_n+b)\|>L_0)\le1/96$.
Thus the weighted untame energy U has $\E U\le A/96$ and
$\Pr(U>A/8)\le1/12$. With probability at least $1/4$, active tame
sites therefore carry energy at least $A/8$. No independence across
sites was used.

Independently, the V event has positive probability: a small open ball
around a matrix with first column a unit class contrast and other
columns zero is contained in it. At tame sites on that event,
$\|\hat y_n\|_\infty\le2L_0$ and $p_{\min,n}\ge e^{-4L_0}/C$.
The unit witness $\Delta W=e_1u^\top$ has logit perturbation
$(h_n^\top u)\chi_nVe_1$, hence curvature at least
$(e^{-4L_0}/C)(1/4)(A/8)=c_0A$. Multiply the independent event
probabilities. Nonnegativity gives the expectation bound.
\end{proof}
\begin{remark}[Scope and a squared-loss counterpart]
A fixed channel direction may be inactive on every site, making its
bound zero on some draws. Adaptive directions can improve the witness,
but are not required for validity. For random upstream features, condition
on an event where both feature bounds hold; its probability multiplies
$p_0$. The theorem covers a pointwise ReLU bottleneck followed by a
linear readout, not an arbitrary convolutional BN decoder. It supplies
no residual-subspace coercivity for Theorem~\ref{thm:astra_affine}.
For scalar squared loss, independent readout entries
$v_i\sim\mathcal N(0,s_v^2/q)$ give the simpler bound
\[
 \E\lambda_{\max}(G_{WW})\ge\frac{s_v^2}{2q}\lambda_{\max}(S_h).
\]
Indeed, the same witness has curvature $v_1^2Q$, and independence gives
$\E(v_1^2Q)=(s_v^2/q)(A/2)$.
\end{remark}

\begin{proposition}[Counterexamples 3.3 and 3.4]
\label{prop:astra_bottleneck_counter}
A fixed ReLU bottleneck does not in general admit a positive interior
floor with probability $1-o_M(1)$. Moreover, train-mode BN can erase
an entire width-linear incoming witness even with a positive stabilizer.
\end{proposition}
\begin{proof}
For the first claim take all $h_n=h\ne0$ with $\|h\|^2=M$ and q
independent centered nondegenerate Gaussian preactivations. With
probability $2^{-q}$ all q gates are inactive at every site, and
$J_W=0$. This probability is independent of M; centered Gaussian bias
does not remove it.

For the second take $h_n=h$ throughout the \emph{whole normalization
group}, $\|h\|^2=M$, and place train-mode BN immediately after $Wh_n$.
For every positive stabilizer $\varepsilon>0$, its output is the BN
affine offset, independent of W. Thus every weight perturbation is
annihilated and $G_{WW}=0$, whereas
$S_h^p=p_{\min}hh^\top$ has eigenvalue $p_{\min}M$ if the downstream
logits have fixed finite nondegenerate probabilities. At
$\varepsilon=0$ this zero-variance forward pass is undefined; it is
not that counterexample. Constant inputs with a $1\times1$ convolution
or periodic boundaries give literal convolutional realizations.
\end{proof}

\begin{proposition}[BN derivative: an upper bound, not a lower witness]
\label{prop:astra_bn_derivative}
For one normalization group of size L, let
$P=I-\boldsymbol1\boldsymbol1^\top/L$, $t=Pu$,
$s^2=\|t\|^2/L$, and $\widetilde\sigma=\sqrt{s^2+\varepsilon}>0$,
with $\varepsilon\ge0$. With fixed affine scale $\gamma$ and offset,
\begin{equation}\label{eq:astra_bn_derivative}
 D_{\rm BN}=\frac{\gamma}{\widetilde\sigma}
 \left(P-\frac{tt^\top}{L(s^2+\varepsilon)}\right),\qquad
 \|D_{\rm BN}d\|^2\le\frac{\gamma^2}{s^2+\varepsilon}\|Pd\|^2.
\end{equation}
If $s^2>0$, decompose $Pd=d_\parallel+d_\perp$ parallel and orthogonal
to t. Then exactly
\[
 \|D_{\rm BN}d\|^2=\frac{\gamma^2}{s^2+\varepsilon}
 \left[\|d_\perp\|^2+
 \left(\frac{\varepsilon}{s^2+\varepsilon}\right)^2\|d_\parallel\|^2\right].
\]
\end{proposition}
\begin{proof}
Differentiate $\gamma Pu/\sqrt{\|Pu\|^2/L+\varepsilon}$.
The middle matrix annihilates constants, has eigenvalue one on
centered vectors orthogonal to t, and eigenvalue
$\varepsilon/(s^2+\varepsilon)$ in the t direction. This proves both
identities and the contraction upper bound. If $s^2=0$ and
$\varepsilon>0$, the derivative is simply $\gamma P/\sqrt\varepsilon$.
\end{proof}

\begin{proposition}[Function-preserving BN scale law]
\label{prop:astra_bn_scale}
Suppose a linear or convolutional block W is used downstream only
through its immediately following train-mode BN. Scale W and its bias
by c>0, keep the other parameters fixed, and co-scale the stabilizer
$\varepsilon\mapsto c^2\varepsilon$. On batches with defined forward
passes, logits are unchanged and
\begin{equation}\label{eq:astra_bn_scale}
 J_W(cW)=c^{-1}J_W(W),\qquad G_{WW}(cW)=c^{-2}G_{WW}(W),
\end{equation}
where the left side uses the co-scaled bias and stabilizer.
For $\varepsilon=0$ the same identities hold without a stabilizer change
provided every channel has positive variance.
\end{proposition}
\begin{proof}
Centering is linear and variance scales by $c^2$, so
$\operatorname{BN}_{c^2\varepsilon}(cu)=\operatorname{BN}_{\varepsilon}(u)$.
This identity holds as a function of W in a neighborhood where defined;
differentiate it and use unchanged output probabilities in the GGN.
The incoming bias Jacobian also scales by $c^{-1}$. Jacobians with
respect to parameters that are not co-scaled remain unchanged, by
differentiating the function identity with respect to those parameters.
\end{proof}
\begin{remark}[Known parameterization sensitivity]
The effective-learning-rate effect of normalization is established
\citep{vanlaarhoven2017l2}. The scalar
$\|W\|_F^2\lambda_{\max}(G_{WW})$ is invariant under this uniform
rescaling, not under arbitrary channelwise rescaling. For positive
diagonal D, channelwise scaling W to DW with channelwise stabilizers
co-scaled is also function preserving, but the GGN transforms by a
nonuniform inverse congruence; the scalar generally changes.
The audited four-channel example with scales $(0.2,1,3,8)$ changed it
by a factor 235 at identical logits. This is an example, not a universal
factor. At fixed positive $\varepsilon$, report each $s_i^2/\varepsilon$
and the scale-control departure: neither a universal accuracy threshold
such as $s_i^2\ge10^3\varepsilon$ nor a final-logit error bound follows
without bounds on c and on downstream sensitivity. Empirical tolerances
from a particular block are diagnostics only.
\end{remark}

\begin{proposition}[A constant centered Gram can coexist with growing curvature]
\label{prop:astra_center_counter}
For $M\ge2$, take three feature rows $(1,0)$, $(0,1)$, $(-1,-1)$ padded
with zeros to dimension M, and $w_M=e_1/\sqrt M$. Put a single train-mode
BN with $\gamma=1$, $\varepsilon=0$ after $w_M^\top h$, followed by
binary logits $(\operatorname{BN}(w_M^\top h),0)$. Then the centered
incoming Gram has top eigenvalue one, while
$\lambda_{\max}(G_{ww})=M c_*$ for a constant $c_*>0$.
\end{proposition}
\begin{proof}
The nonzero Gram block is
$\left(\begin{smallmatrix}2/3&1/3\\1/3&2/3\end{smallmatrix}\right)$.
At $w=e_1$ the centered preactivation is $(1,0,-1)$, with positive
variance, and the perturbation $e_2$ induces $(0,1,-1)$ with a
nonzero component transverse to it. The BN derivative and finite-logit
CE curvature give a positive Rayleigh quotient. Apply
Proposition~\ref{prop:astra_bn_scale} with $c=M^{-1/2}$.
This growth is precisely that function-preserving rescaling: the
normalization $\|w_M\|=M^{-1/2}$ is the whole mechanism. It refutes
an \emph{unqualified} necessity claim for a width-linear centered Gram,
not a claim that assumes width-uniform downstream gain. It is not a
normalization-free mechanism.
\end{proof}

\begin{proposition}[Conditional normalization-gain width law]
\label{prop:astra_bn_gain}
Let input dimension be d (for a fixed-size convolution, d is a fixed
multiple of $M+K$). On a fixed BN group let the \emph{centered} feature
covariance be $S_c$, and initialize a row w independently with
$\E w=0$, $\E ww^\top=\sigma_w^2 I_d/d$. Its pre-BN variance satisfies
\begin{equation}\label{eq:astra_bn_variance}
 \E_w s^2=\sigma_w^2\operatorname{tr}(S_c)/d.
\end{equation}
For a realized row assume $c_-/d\le s^2\le c_+/d$, with
$0<c_-\le c_+<\infty$, and
$0<\gamma_-\le|\gamma|\le\gamma_+$. Then
\begin{equation}\label{eq:astra_bn_gain}
 \frac{\gamma_-^2d}{c_++\varepsilon d}
 \le\frac{\gamma^2}{s^2+\varepsilon}
 \le\frac{\gamma_+^2d}{c_-+\varepsilon d}.
\end{equation}
For a given unit parameter witness whose pre-BN perturbation d' has
$\|d'_\perp\|^2\ge\tau\|d'\|^2$ with $\tau>0$, its post-BN energy
is at least
$\tau\gamma_-^2d\|d'\|^2/(c_++\varepsilon d)$.
A final GGN lower bound further requires a positive lower gain of the
actual downstream softmax-weighted, masked map on that post-BN perturbation.
\end{proposition}
\begin{proof}
$s^2=w^\top S_cw$, so taking a trace proves
\eqref{eq:astra_bn_variance}. Invert the two realized variance bounds
to obtain \eqref{eq:astra_bn_gain}. The transverse energy identity in
Proposition~\ref{prop:astra_bn_derivative} proves the witness assertion;
then apply \eqref{eq:astra_full_witness} with the stated downstream gain.
\end{proof}
\begin{remark}[Diagnostic model and its premises]
If $s^2=c_0/(M+K)$, its squared BN gain is exactly
\[
 \gamma^2(M+K)/(c_0+\varepsilon(M+K)).
\]
The premise is fixed positive \emph{total centered} incoming energy as
channels are added, together with control of the realized row variance.
Standard fan-in initialization with per-channel centered energy fixed
instead gives width-independent \emph{expected} pre-BN variance.
A fixed total uncentered second moment is insufficient (constant
features have zero centered variance). An expectation of variance cannot
be inverted into a concentration statement about gain. Even with fixed
pre-BN witness energy, a radial witness is annihilated at
$\varepsilon=0$ and need not inherit the gain. These qualifications
make this a conditional diagnostic identity, not a new general BN theorem
or an unconditional replacement for Lemma~\ref{lem:phi_scaling}.
\end{remark}
\section{Residual excitation and phase-restricted attribution}
\label{sec:astra_excitation}
\begin{lemma}[Directional gradient gain is not directional GGN curvature]
\label{lem:astra_direction_gain}
Let $R_u$ isometrically embed a directional parameter subspace and
$J_u=J_\thetap R_u$. Then
\[
 \|J_u^\top r\|^2\le\lambda_{\max}(J_u^\top J_u)\|r\|^2.
\]
If $H\succeq0$ and $r\in\operatorname{range}(H)$, also
\begin{equation}\label{eq:astra_pseudoinverse}
 \|J_u^\top r\|^2\le\lambda_{\max}(J_u^\top HJ_u)r^\top H^\dagger r.
\end{equation}
For finite-logit softmax CE the range condition holds. A bound
$H\succeq h_0\Pi_{\rm cls}$ on the class-contrast space, $h_0>0$,
permits $r^\top H^\dagger r\le\|r\|^2/h_0$.
\end{lemma}
\begin{proof}
The first claim is the operator-norm inequality. For the second, write
\[J_u^\top r=(H^{1/2}J_u)^\top H^{\dagger/2}r.\]
Apply the operator-norm inequality again.
The last assertion follows by diagonalizing H on its range.
Here the Hessian floor is understood for the softmax Hessian, whose
range is exactly the class-contrast space at finite logits.
\end{proof}
\begin{remark}[CE counterexample and saturation]
For one example with logits $(\theta,0)$, $\theta=0$, and label one, the residual and Jacobian are
\[r=(-1/2,1/2),\qquad J=(1,0)^\top.\]
Then
$|\nabla_\theta\loss|^2=1/4$, while
$G_{\theta\theta}\|r\|^2=(1/4)(1/2)=1/8$.
Thus replacing the logit Gram by the GGN without the factor in
\eqref{eq:astra_pseudoinverse} is false. At general first-class probability
p the ratio is $1/(2p(1-p))$, unbounded near zero or one.
A uniform positive h0 is a further assumption, not supplied by
saturation. For a linear encoder and unit input direction u,
$R_u b=\operatorname{vec}(bu^\top)$, so $J_u^\top r=(\nabla_\Theta\loss)u$.
\end{remark}

\begin{theorem}[Cumulative energy and excitation; review Theorem 2.1]
\label{thm:astra_excitation}
Let $A_u,A_v$ be constant directional logit Jacobians along a path,
$r(t)=\nabla_z\ell(z(t))$ square integrable on $[0,T]$, and
$C_u=A_uA_u^\top$, $\lambda_u=\|C_u\|$, with analogous v notation.
Define $Q_T=\int_0^T r(t)r(t)^\top dt$ and
$E_u=\int_0^T\|A_u^\top r(t)\|^2dt$. Then
\[
 E_u=\operatorname{tr}(C_uQ_T)\le\lambda_u\operatorname{tr}(Q_T).
\]
Suppose $\lambda_v>0$, $\lambda_u>0$, $\operatorname{tr}(Q_T)>0$,
and a top unit eigenvector $q_v$ of $C_v$ satisfies
$q_v^\top Q_Tq_v\ge\rho\operatorname{tr}(Q_T)$ with $\rho>0$.
Then $E_v\ge\rho(\lambda_v/\lambda_u)E_u$; division by $E_u$ requires
$E_u>0$. More generally, if $C_v\succeq c\lambda_vP$ and
$\operatorname{tr}(PQ_T)\ge\rho\operatorname{tr}(Q_T)$ for a projector P,
replace $\rho$ by $c\rho$.
\end{theorem}
\begin{proof}
Integrate $\|A_u^\top r\|^2=r^\top C_ur$ and use
$C_u\preceq\lambda_uI$. Since $C_v\succeq\lambda_vq_vq_v^\top$,
$E_v\ge\lambda_v q_v^\top Q_Tq_v\ge\rho\lambda_v\operatorname{tr}(Q_T)$.
The projector version is identical.
\end{proof}
\begin{remark}
This theorem is loss-agnostic under its constant-Jacobian and integrability
hypotheses; it includes affine CE. An envelope
$\|r(t)\|\le C_r/(1+\mu t)$ gives $E_u\le\lambda_u C_r^2/\mu$.
Here $\lambda_u$ is a logit-kernel eigenvalue. The exponential residual
formulas in the next counterexamples specifically use squared loss.
Excitation is an output-space condition on the visited residual; an
input principal component alone does not establish it.
\end{remark}

\begin{proposition}[Counterexamples 2.2 and 2.3 and the decoupled ratio]
\label{prop:astra_energy_counter}
For affine squared loss take
$C_v=L e_1e_1^\top$, $C_u=\ell e_2e_2^\top$, $L,\ell>0$, with
total kernel $\operatorname{diag}(L,\ell)$.
If $r(0)=e_2$, then $E_v(T)=0<E_u(T)$ for $T>0$.
If $r(0)=e_1+e_2$, then
\[
 E_v(T)=(1-e^{-2LT})/2,\qquad E_u(T)=(1-e^{-2\ell T})/2,
\]
so their ratio tends to one even for arbitrarily large $L/\ell$.
More generally assume e1,e2 are orthonormal eigenvectors of the total
constant kernel with positive eigenvalues $\gamma_v,\gamma_u$,
$C_v=\lambda_v e_1e_1^\top$, $C_u=\lambda_u e_2e_2^\top$, and the
initial residual components are $a_v,a_u$ with $a_u\ne0$,
$\lambda_u>0$, $T>0$. Then
\begin{equation}\label{eq:astra_decoupled}
 \frac{E_v(T)}{E_u(T)}=
 \frac{\lambda_v}{\lambda_u}\frac{a_v^2}{a_u^2}
 \frac{\gamma_u}{\gamma_v}
 \frac{1-e^{-2\gamma_vT}}{1-e^{-2\gamma_uT}}.
\end{equation}
\end{proposition}
\begin{proof}
Each residual component is $a_je^{-\gamma_jt}$; integrate
$\lambda_ja_j^2e^{-2\gamma_jt}$. The first two constructions are
realized by the affine map $(\sqrt L\theta_1,\sqrt\ell\theta_2)$.
Without the decoupling hypothesis, these scalar exponentials do not
represent mixed residual components.
\end{proof}

\begin{proposition}[Attribution during a residual-excited phase]
\label{prop:astra_phase}
Assume the differentiability, absolute continuity, unit-flow and
energy-identity conventions of Corollary~\ref{cor:astra_attribution},
but not its global Rayleigh bounds. Let
$q_b(t)=r^\top K_br$, $D(t)=q_\thetap(t)+q_\phip(t)$, and
$\Delta L=\int_0^TD(t)dt=\loss(0)-\loss(T)>0$.
For a measurable phase $E\subset[0,T]$, put
\[
 \tau_E=|E|,\qquad w_E=\frac{\int_E D(t)dt}{\Delta L}.
\]
On E suppose a measurable top-k projector $P_k(t)$ of $K_\phip(t)$
(with a fixed measurable tie convention) satisfies
\[
 \|P_k(t)r(t)\|^2\ge\alpha\|r(t)\|^2,\quad
 \lambda_k(K_\phip(t))\ge\ell_*>0,\quad
 q_\thetap(t)\le a\|r(t)\|^2,
\]
where $0<\alpha\le1$, $a\ge0$. Then, when the phase denominator is positive,
\begin{align}
 \frac{\int_E q_\thetap}{\int_E D}&\le\frac{a}{a+\alpha\ell_*},
 \label{eq:astra_phase}\\
 \frac{\int_0^Tq_\thetap}{\Delta L}
 &\le w_E\frac{a}{a+\alpha\ell_*}+(1-w_E).
 \label{eq:astra_phase_total}
\end{align}
A verified complementary share bound b replaces the last term by
$(1-w_E)b$. Time-varying pointwise constants give the sharper
$D(t)dt/\Delta L$-weighted average of their share caps.
\end{proposition}
\begin{proof}
PSD spectral decomposition gives
$q_\phip\ge\lambda_k\|P_kr\|^2\ge\alpha\ell_*\|r\|^2$.
Thus $q_\thetap\le[a/(a+\alpha\ell_*)]D$ on E.
Integrate on E and its complement, where $q_\thetap\le D$.
The phase duration is explicit but imposes no value on $w_E$.
For $E=[0,\tau]$, $\int_E D=\loss(0)-\loss(\tau)$; under squared
loss additionally $\|r(t)\|\le e^{-\alpha\ell_*t}\|r(0)\|$
through this interval by differentiating its squared norm.
\end{proof}
\begin{corollary}[What a width-dependent cumulative share additionally requires]
\label{cor:astra_phase_width}
For a family satisfying Proposition~\ref{prop:astra_phase}, suppose
its early-phase constants obey $a\le a_*$ and
$\alpha\ell_*\ge cM$, with fixed $a_*\ge0,c>0$. On the complement
assume $0<b_-\le q_\thetap/D\le b_+\le1$ at almost every time with
$D>0$, for constants independent of M. Writing the cumulative share
as $S_M$, one has
\[
 b_-(1-w_E)\le S_M
 \le w_E\frac{a_*}{a_*+cM}+b_+(1-w_E).
\]
Consequently $S_M=O(M^{-1})$ holds if and only if
$1-w_E=O(M^{-1})$. The same statement holds for nonnegative weighted
sums on one discrete sequence, interpreting $w_E$ and $S_M$ as
\emph{gradient-energy} fractions.
\end{corollary}
\begin{proof}
The upper bound is Proposition~\ref{prop:astra_phase}. For the lower
bound retain only the complementary numerator and use its share floor.
Both implications follow since $b_->0$ and $a_*/(a_*+cM)=O(M^{-1})$.
This characterizes what the claimed cumulative scaling needs under the
stated late-phase floor; early excitation or a short phase alone is insufficient.
\end{proof}

\begin{remark}[Discrete diagnostic and admissible reference space]
Replacing integrals by nonnegative weighted sums of gradient squares
proves the same \emph{gradient-energy} inequalities on any one sampled
sequence, regardless of optimizer. This is not Adam loss-decrease
attribution. Probe-defined phases and training-minibatch energies are
different sequences unless explicitly matched. No phase-share bound
alone replaces the stability assumptions needed to compare two trajectories.
For CE let $\Pi_{\rm cls}=I_N\otimes(I_C-\boldsymbol1\boldsymbol1^\top/C)$.
The residual belongs to its rank-$N(C-1)$ range. For a uniformly random
unit vector in that range, expected energy in a projector P is
$\operatorname{tr}(P\Pi_{\rm cls})/[N(C-1)]$, since its second moment is
$\Pi_{\rm cls}/[N(C-1)]$. The full-space baseline
$\operatorname{rank}(P)/(NC)$ is generally inappropriate.
\end{remark}

\section{A solvable serial cross-entropy instance}
\label{sec:astra_serial}
\begin{proposition}[Information retained by a fitting serial encoder]
\label{prop:astra_information}
If the predictions depend on X only through $Z=f_\thetap(X)$ and a
classifier $q_\phip(Y\mid Z)$ has population CE at most $\eta<\infty$,
then $H(Y\mid Z)\le\eta$ and $I(Y;Z)\ge H(Y)-\eta$.
The same statement holds for the empirical distribution with a fixed model.
\end{proposition}
\begin{proof}
CE is conditional entropy plus the nonnegative expected conditional KL
divergence. Suppression of a selected coordinate or failure under a
specified shift is therefore different from absence of all training-label
information in the serial representation.
\end{proof}

\begin{theorem}[Serial CE instance; review Theorem 5.1]
\label{thm:astra_serial}
Let $Y$ be uniform on $\{-1,1\}$ and let training inputs be
$x_1=(S,0)$, $x_2=(0,C_{\rm ctx})$ with $S=C_{\rm ctx}=Y$.
The shared linear encoder is $W=(a,v)$, giving $z_1=aS,z_2=vC_{\rm ctx}$.
For integer $M\ge1$ define
\[
 F=z_1+\sum_{j=1}^M\beta_jz_2=aS+bvC_{\rm ctx},\quad
 b=\sum_j\beta_j,\quad\thetap=(a,v),\quad\phip=\beta.
\]
All coordinates follow unit Euclidean gradient flow on the averaged
binary CE $\loss=\E\log(1+e^{-YF})=\log(1+e^{-q})$, $q=a+bv$,
from $a_0=0,v_0=1,\beta_0=0$, with no clipping, decay or momentum.
Fix $0<\delta<1/2$, $m=\log((1-\delta)/\delta)$, and let $T_m$ be
the first time q reaches m. Write $\varrho=(1+e^q)^{-1}$ and
$\Psi(q)=q+e^q-1$. Then:
\begin{align}
 \dot a&=\varrho,&\dot v&=b\varrho,&\dot b&=Mv\varrho,\label{eq:astra_serial_ode}\\
 v&=\cosh(\sqrt M a),&b&=\sqrt M\sinh(\sqrt M a),\label{eq:astra_serial_phase}\\
 q&=a+\frac{\sqrt M}{2}\sinh(2\sqrt M a),&
 \dot q&=(M+1+2b^2)\varrho.\label{eq:astra_serial_q}
\end{align}
The solution is global and reaches m in finite time, with
\begin{equation}\label{eq:astra_serial_time}
 \frac{\Psi(m)}{M+1+2m^2}\le T_m\le\frac{\Psi(m)}{M+1}.
\end{equation}
At that time, writing $a_M=a(T_m)$,
\begin{gather}
 0<a_M\le\frac{m}{M+1},\qquad b(T_m)v(T_m)\ge\frac{Mm}{M+1},
 \qquad Ma_M\longrightarrow m,\label{eq:astra_serial_update}\\
 0\le v(T_m)-1\le\frac{m^2}{2M},\qquad
 \|W(T_m)-W_0\|\le
 \sqrt{\frac{m^2}{(M+1)^2}+\frac{m^4}{4M^2}}.\label{eq:astra_serial_disp}
\end{gather}
For fixed m, $a_M$ decreases strictly with M; removing the contextual
path entirely requires $a=m$ at the same fitting margin. Moreover,
\begin{equation}\label{eq:astra_serial_envelope}
 \varrho(t)\le\frac{1/2}{1+(M+1)t/4},\qquad t\ge0.
\end{equation}
In the stated parameter blocks and loss convention,
\[
 \lambda_{\max}(G_{\thetap\thetap})=\varrho(1-\varrho)(1+b^2),\quad
 \lambda_{\max}(G_{\phip\phip})=M\varrho(1-\varrho)v^2,
\]
so at initialization they equal $1/4,M/4$ and the disparity is M.
For the separately trained frozen encoder $W=(0,1)$, $\beta_F(0)=0$,
its margin obeys $\Psi(q_F(t))=Mt$ and
\begin{equation}\label{eq:astra_serial_gap}
 0\le q(t)-q_F(t)\le\frac{(1+2m^2)t}{2+Mt/2}
 \le\frac{2(1+2m^2)}M,\qquad0\le t\le T_m.
\end{equation}
Under the specified reversal $S=Y,C_{\rm ctx}=-Y$, the joint model at
$T_m$ has error one; the fitted spectral-only comparator has error zero.
\end{theorem}
\begin{proof}
The training margin is the same on each example, so direct differentiation
gives \eqref{eq:astra_serial_ode}. Divide by $\dot a=\varrho>0$ to get
$dv/da=b$, $db/da=Mv$, which proves \eqref{eq:astra_serial_phase}.
It implies $b^2=M(v^2-1)$ and \eqref{eq:astra_serial_q}.
Since $0<\dot a\le1/2$, a stays finite on every finite time interval;
the phase formula then gives global existence. The strictly increasing
map $Q_M(a)=a+(\sqrt M/2)\sinh(2\sqrt M a)$ maps $[0,\infty)$
onto itself and has a finite inverse at m. On this compact a interval,
$\dot a\ge\delta>0$, so the hitting time is finite.

$\sinh x\ge x$ gives $Q_M(a)\ge(M+1)a$, strictly for $a>0$.
Before fitting, $v\ge1$, $0\le bv=q-a\le m$, hence $b\le m$.
The invariant gives $v-1=b^2/[M(v+1)]\le m^2/(2M)$.
For fixed positive a the contextual term increases strictly with M;
invert to get strict decrease of $a_M$. Its upper bound implies
$\sqrt M a_M\to0$, and the expansion
$m=(M+1)a_M+O(M^2a_M^3)$ gives $Ma_M\to m$.

$\dot\varrho=-(M+1+2b^2)\varrho^2(1-\varrho)$ and
$\varrho\le1/2$ imply $(d/dt)(1/\varrho)\ge(M+1)/2$.
Integrate from $\varrho(0)=1/2$ for the envelope.
Since $\Psi'(q)=1/\varrho$, before fitting
$\frac d{dt}\Psi(q)=M+1+2b^2\in[M+1,M+1+2m^2]$;
integration proves \eqref{eq:astra_serial_time}.
The margin gradients are $Y(1,b)$ and $Yv\boldsymbol1_M$;
the scalar logistic Hessian is $\varrho(1-\varrho)$, giving the GGN.

For the frozen system $\dot q_F=M/(1+e^{q_F})$. Thus
$0\le\Psi(q)-\Psi(q_F)=\int_0^t(1+2b^2)ds\le(1+2m^2)t$.
Convexity gives a lower bound $\Psi'(q_F)(q-q_F)$ for this difference.
As $q_F\le e^{q_F}-1$,
$Mt=\Psi(q_F)\le2(e^{q_F}-1)$, so
$\Psi'(q_F)\ge2+Mt/2$, proving \eqref{eq:astra_serial_gap}.
The shifted margin is $a-bv=2a_M-m<0$: for $M>1$ use
\eqref{eq:astra_serial_update}, and for $M=1$ use the strict
inequality $Q_1(a)>2a$. The spectral-only shifted margin is m.
\end{proof}

\begin{corollary}[Analytic instantiation of the displacement bound]
\label{cor:astra_serial_bound}
For Theorem~\ref{thm:astra_serial}, embed the scalar prediction in
symmetric logits $(F/2,-F/2)$. The packet residual norm is
$R=\sqrt2\varrho$, and on $[0,T_m]$
\[
 \|J_\thetap\|=\sqrt{(1+b^2)/2}\le B_m:=\sqrt{(1+m^2)/2}.
\]
With envelope amplitude $C_r=1$ and $\mu=(M+1)/4$,
Theorem~\ref{thm:twoscale} gives
\begin{equation}\label{eq:astra_serial_bound}
 \|W(T_m)-W_0\|\le\frac{4B_m}{M+1}\log\frac1{\sqrt2\delta}.
\end{equation}
Also $\|W(T_m)-W_0\|\le\sqrt{1+m^2}\,a_M
\le\sqrt{1+m^2}m/(M+1)$.
\end{corollary}
\begin{proof}
The symmetric logit derivative has squared norm $(1+b^2)/2$, while
the full residual norm is $\sqrt2\varrho$, so their product is the
exact encoder gradient norm. The envelope has true initial amplitude
$1/\sqrt2$; $C_r=1$ is a valid larger amplitude matching the main
text's $C_r\ge1$ convention. Its target is $\sqrt2\delta<1$.
All constants apply on the fitting region $|b|\le m$.
Alternatively integrate $\|\dot W\|=\varrho\sqrt{1+b^2}
\le\sqrt{1+m^2}\dot a$.
\end{proof}
\begin{remark}[Embedding, head, and initialization conventions]
For logits $(F,0)$, loss, residual norm and GGN agree, but the unprojected
logit-Jacobian norm is $\sqrt{1+b^2}$; the symmetric embedding makes
the norm-product bound exact. A freely trained dense two-row head is
not the parameter block $\phip=\beta$: its initial head eigenvalue is
$M/2$, not $M/4$, while the encoder eigenvalue remains $1/4$ at the
same initial function. Labels need not be balanced for the training
ODE; balance is used in the information/noisy-accuracy statement below.
The initial encoder orientation in Theorem~\ref{thm:astra_serial} is
explicitly asymmetric; the next corollary removes that choice for W,
not the routing/readout asymmetry.
\end{remark}

\begin{proposition}[Readout normalization and representation interpretation]
\label{prop:astra_serial_control}
Replace the contextual readout by $\alpha_M\sum_j\gamma_jz_2$,
$\gamma_j(0)=0$, at unit rates. Its reduced equation for
$b=\alpha_M\sum_j\gamma_j$ is $\dot b=M\alpha_M^2v\varrho$.
In particular $\alpha_M=M^{-1/2}$ gives exactly the $M=1$ dynamics
for every M. At every finite M in the original theorem, the noiseless
spectral coordinate $a_MS$ still determines Y. If it is instead
observed as $Z_S=a_MY+\sigma\xi$, with $\xi\sim\mathcal N(0,1)$
independent, $\sigma>0$ fixed, then Bayes accuracy is
$\Phi(a_M/\sigma)$ and $I(Y;Z_S)\le a_M^2/(2\sigma^2)$.
\end{proposition}
\begin{proof}
Differentiate each readout and sum. The Gaussian likelihood-ratio test
is thresholding at zero, giving the accuracy. For
$Q=\mathcal N(0,\sigma^2)$,
$\E_Y\operatorname{KL}(P_{Z_S|Y}\|Q)=a_M^2/(2\sigma^2)
=I(Y;Z_S)+\operatorname{KL}(P_{Z_S}\|Q)$.
This noise is a specified observation perturbation, not training noise.
\end{proof}

\begin{corollary}[General and isotropic encoder initialization]
\label{cor:astra_isotropic}
In the serial model let $m>0$, $a_0<m$, $v_0\ne0$, and allow any
readout initialization with $\sum_j\beta_{j0}=0$. Set
$u=a-a_0$, $\Delta=m-a_0>0$. Then
\begin{align}
 v&=v_0\cosh(\sqrt M u),\quad b=\sqrt Mv_0\sinh(\sqrt M u),
 \label{eq:astra_iso_phase}\\
 q&=a_0+u+\frac{\sqrt Mv_0^2}{2}\sinh(2\sqrt M u),\notag\\
 0<u_M&\le\frac{\Delta}{1+Mv_0^2},\qquad
 Mu_M\longrightarrow\frac{\Delta}{v_0^2},\qquad
 \frac{a(T_m)-a_0}{m-a_0}\longrightarrow0.
 \label{eq:astra_iso_update}
\end{align}
The fitting root decreases strictly with M and
\begin{equation}\label{eq:astra_iso_disp}
 |v(T_m)-v_0|\le\frac{\Delta^2}{2M|v_0|^3},\qquad
 \|W(T_m)-W_0\|\le
 \sqrt{\frac{\Delta^2}{(1+Mv_0^2)^2}+
       \frac{\Delta^4}{4M^2|v_0|^6}}.
\end{equation}
With $\Psi_{a_0}(q)=q-a_0+e^q-e^{a_0}$,
\[
 \frac{\Psi_{a_0}(m)}{1+Mv_0^2+2\Delta^2/v_0^2}
 \le T_m\le\frac{\Psi_{a_0}(m)}{1+Mv_0^2}.
\]
For a frozen encoder at $(a_0,v_0)$ and the same initial readouts,
$\Psi_{a_0}(q_F)=Mv_0^2t$. If
$p_0=e^{a_0}/(1+e^{a_0})$ and $C_0=1+2\Delta^2/v_0^2$, then
\begin{equation}\label{eq:astra_iso_gap}
 0\le q_J-q_F\le\frac{C_0t}{1+e^{a_0}+p_0Mv_0^2t}
 \le\frac{C_0}{p_0Mv_0^2},\qquad0\le t\le T_m.
\end{equation}
Under contextual reversal the fitted margin is $2a(T_m)-m$.
For $a_0<m/2$, error is one for all sufficiently large M, with sufficient
condition $Mv_0^2>m/(m-2a_0)$. For $m/2\le a_0<m$, the reversed
classifier is correct at every width. If $a_0\ge m$, define fitting to
stop immediately, in which case it is also correct.

If $(a_0,v_0)\sim\mathcal N(0,\sigma^2 I_2)$, $\sigma>0$, then
$v_0\ne0$ almost surely; use the preceding stopping convention.
The expected reversal error satisfies
\begin{equation}\label{eq:astra_iso_error}
 \lim_{M\to\infty}\E_{W_0}\operatorname{Err}_{\rm reversal}
 =\Phi\left(\frac{m}{2\sigma}\right).
\end{equation}
The same-threshold spectral-only comparator has zero reversal error.
\end{corollary}
\begin{proof}
Dividing the ODE by $\dot a>0$ gives
\eqref{eq:astra_iso_phase}; only the initial readout sum enters.
Each $\beta_i-\beta_j$ is conserved, so equal initial readouts are not
needed for closure. Positivity of the contextual margin
$b v=(\sqrt Mv_0^2/2)\sinh(2\sqrt M u)$ and the sinh inequality
prove \eqref{eq:astra_iso_update}; Taylor expansion at
$\sqrt M u_M\to0$ gives the limit. The invariant is
$b^2=M(v^2-v_0^2)$; v keeps its sign and $|v|\ge|v_0|$.
Since $bv\le\Delta$, $|b|\le\Delta/|v_0|$ and
\[
 |v-v_0|=\frac{b^2}{M(|v|+|v_0|)}
 \le\frac{\Delta^2}{2M|v_0|^3}.
\]
This finite-M inequality proves \eqref{eq:astra_iso_disp}.
Now $\dot q=(1+Mv_0^2+2b^2)\varrho$ and
$b^2\le\Delta^2/v_0^2$ on the fitting interval; integrate the
transformed margin for the time bounds. Global existence follows as
before from $\dot u\le1$ and the phase formula; every finite m is hit.

The transformed joint/frozen difference lies between zero and $C_0t$.
Also $q_F-a_0\le(e^{q_F}-e^{a_0})/e^{a_0}$ gives
$\Psi'_{a_0}(q_F)\ge1+e^{a_0}+p_0Mv_0^2t$.
Convexity proves \eqref{eq:astra_iso_gap}. For $a_0<m/2$, use
$a(T_m)\le a_0+\Delta/(1+Mv_0^2)$ to get the finite-M reversal
condition. For $a_0>m/2$ the classifier is correct at every fitted
width, since $u_M>0$ gives $2a(T_m)-m>2a_0-m>0$; equality
$a_0=m/2$ also gives a positive margin. Immediate stopping covers
$a_0\ge m$. Almost surely the limiting error is
$1_{\{a_0<m/2\}}$, apart from the zero-probability boundary and
$v_0=0$ events. Dominated convergence gives \eqref{eq:astra_iso_error}.
\end{proof}

\begin{remark}[Confidence dependence and limiting quantity]
For $a_0\ne0$, $a(T_m)/a_0\to1$ but $a(T_m)/m\to a_0/m$;
the vanishing quantity is the \emph{update} fraction in
\eqref{eq:astra_iso_update}. For $0<\rho<1$, a Gaussian tail and its
maximum density give, with probability at least $1-\rho$,
\[
 |a_0|\le\sigma\sqrt{2\log(4/\rho)},\qquad
 |v_0|\ge\frac{\rho\sigma}{2}\sqrt{\pi/2}.
\]
Substitute these bounds in the finite-M inequalities for uniform
confidence bounds. The dependence is severe: for fixed
$(\Delta,v_0)\in(0,\infty)\times(\R\setminus\{0\})$, expansion
of the root yields
\[
 Mu_M=\frac{\Delta}{v_0^2}\left[
 1-\frac1{Mv_0^2}-\frac{2\Delta^2}{3Mv_0^4}+O(M^{-2})\right].
\]
Indeed substitute $u=A/M+B/M^2$ in the sinh equation and match its
constant and $M^{-1}$ coefficients. The expansion's remainder is
pointwise, not already uniform in rho. Its small-$|v_0|$ crossover
requires control of $\Delta^2/(Mv_0^4)$; at the confidence cutoff this
has scale $\Delta^2\rho^{-4}\sigma^{-4}/M$, a fourth power of
$1/\rho$, not a second. A rigorous uniform remainder can instead be
bounded from the sinh series on this controlled argument. Do not
present this asymptotic crossover as a universal necessary width.
\end{remark}

\paragraph{Provenance and interpretation.}
Architecture- and initialization-dependent implicit bias is established
in linear-network analyses \citep{astra_yun2021,astra_moroshko2020};
trajectory-level incremental learning is studied by
\citet{astra_berthier2023}. Exact deep-linear dynamics
\citep{astra_saxe2014} and conservation of layer-norm differences
\citep{astra_du2018balance} supply the underlying tools. CE feature
competition is central to \citet{pezeshki2021gradient}.
The serial example instantiates a known optimization-metric mechanism
with finite-margin formulas, an encoder-update comparison and a specified
reversal distribution. Replication is equivalent to a rate M on b,
and normalization removes it; the function class does not grow.
These particular calculations are not a claim of a new general implicit-bias
mechanism, and finite-time or serial composition alone is not a novelty claim.

\section{A scoped nonlinear existence result}
\label{sec:astra_nonlinear}

\paragraph{Flow convention and normalization.}
Here the loss is squared loss, the encoder is linear, the decoder has
one hidden ReLU layer, and every parameter uses unit Euclidean gradient
flow. There is no BatchNorm, convolution, weight decay, clipping,
preconditioning, or stochastic sampling. Write
\[
 f_n=V\sigma(W\Theta x_n+b)+\beta,\qquad
 \loss=\frac1{2N}\sum_{n=1}^N\norm{f_n-y_n}^2,
 \qquad \thetap=\Theta,\quad\phip=(W,b,V,\beta).
\]
Parameter norms are Euclidean, with Frobenius norms on matrices.
Set $\bar f=N^{-1/2}(f_1,\ldots,f_N)$ and
$e=\bar f-\bar y$, so that $\loss=\norm e^2/2$.
At a ReLU kink choose a measurable gate in $[0,1]$, using the same gate
in every parameter block for that neuron and datum. A solution is an
absolutely continuous parameter path satisfying the resulting gradient
flow almost everywhere. The theorem concerns every such solution;
uniqueness is not asserted. Local existence and the energy identity
needed below are proved in the proof. This formulation avoids identifying
arbitrary backpropagation selections with the Clarke subdifferential.

\begin{theorem}[Theorem N: nonlinear, unnormalized-readout existence result]
\label{thm:astra_nonlinear}
Fix positive integers $N,D,K,C$, fixed data $x_n\in\R^D,y_n\in\R^C$,
and a deterministic $\Theta_0\in\R^{K\times D}$ such that the
$\Theta_0x_n$ are pairwise distinct. Fix $s_w,s_v,s_b>0$ and initialize
independently across neurons, entries, and blocks:
\[
 w_i\sim\mathcal N(0,s_w^2 I_K/K),\quad
 b_i\sim\mathcal N(0,s_b^2),\quad
 V_{ci}\sim\mathcal N(0,s_v^2/M).
\]
The output bias $\beta_0$ is fixed, or independent of the other blocks
with an $M$-independent distribution of finite second moment. In
particular $\beta_0$ need not be centered. Compare joint training from
this initialization with training of all decoder parameters from the
same initialization while fixing $\Theta=\Theta_0$.
For every $\rho\in(0,1)$ there are finite positive constants
$\kappa,A,B,R_0,C_{\rm gap},M_*$, depending on the fixed data,
initialization scales and $\rho$, but not on $M$, with the following
property. For each $M\ge M_*$, with probability at least $1-\rho$,
both flows have global solutions, and every pair of solutions obeys,
for every $t\ge0$,
\begin{align*}
 \norm{e^J(t)},\norm{e^F(t)}
   &\le R_0e^{-\kappa Mt/2}, \tag{N1}\label{eq:astra_n1}\\
 \sup_t\norm{\Theta^J(t)-\Theta_0}
   &\le\frac{2AR_0}{\kappa M},&
 \sup_t\norm{\phip^{J,F}(t)-\phip_0}
   &\le\frac{2BR_0}{\kappa\sqrt M}, \tag{N2}\label{eq:astra_n2}\\
 \sup_t\norm{\bar f^J(t)-\bar f^F(t)}
   &\le\frac{C_{\rm gap}R_0}{\kappa M}. \tag{N3}\label{eq:astra_n3}
\end{align*}
For every finite $T>0$ with $\loss(0)>\loss(T)$ on the joint path,
\begin{equation*}
 \frac{\int_0^T\norm{\nabla_{\thetap}\loss}^2dt}
 {\loss(0)-\loss(T)}
 \le\frac{A^2}{A^2+\kappa M/2}. \tag{N4}\label{eq:astra_n4}
\end{equation*}
For $0<\delta<R_0$, residual norm at most $\delta$ is reached by
$2(\kappa M)^{-1}\log(R_0/\delta)$. If the initial residual is already
at most $\delta$, the first hitting time can be zero.
The constants and threshold constructed below are existence-only and
can be numerically vacuous. This is not a uniform theorem for learned
encoders, growing datasets, general deep decoders, or production Adam.
\end{theorem}

\begin{proof}
We give explicit finite constants to make all quantifiers and
probability events checkable.

\emph{Population kernel and initialization events.}
Let $X=\max_n\norm{x_n}$, $z_n=\Theta_0x_n$, and
\[
 H_{ni}=N^{-1/2}\sigma(w_i^Tz_n+b_i),\quad
 K_* =\E[H_{:i}H_{:i}^T],\quad
 \kappa_* =\lambda_{\min}(K_*),\quad \kappa=\kappa_*/2.
\]
Here $K_*$ uses one random neuron with the displayed normalization,
so that $\E[HH^T]=MK_*$. The distinct-input kink argument in the proof
of Corollary~\ref{cor:astra_init} gives $\kappa_*>0$: if a linear
combination of the ReLU features vanishes almost surely, it vanishes
everywhere by continuity and full Gaussian support in $(w,b)$;
one can cross each distinct affine hyperplane away from the others,
forcing its coefficient to vanish.
Put
\[
 s_{\max}^2=\frac{s_w^2}{K}\max_n\norm{z_n}^2+s_b^2,
 \quad Z_0=\sqrt{1+(\norm{\Theta_0}X)^2},\quad Z=Z_0+X.
\]
Operator norm may be used for $\norm{\Theta_0}$ here; its Frobenius
norm is also a valid upper bound. For $\norm{\Theta-\Theta_0}\le1$,
\begin{equation}\label{eq:astra_zbound}
 \sqrt{1+\norm{\Theta x_n}^2}\le Z.
\end{equation}
Indeed, with $a=\norm{\Theta_0}X$,
$(\sqrt{1+a^2}+X)^2-[1+(a+X)^2]
=2X(\sqrt{1+a^2}-a)\ge0$.

Choose positive $h_0,C_w,C_v,A_0,R_0$ such that each of
\begin{equation}\label{eq:astra_fiveevents}
 \norm{H_0}_F\le h_0\sqrt M,\quad
 \norm{W_0}_F\le C_w\sqrt M,\quad
 \norm{V_0}_F\le C_v,\quad
 \norm{J_{\Theta,0}}_{\rm op}\le A_0,\quad
 \norm{e_0}\le R_0
\end{equation}
has failure probability at most $\rho/7$. Here and below $J$ is
the Jacobian of $\bar f$. For example, Markov's inequality applies
with respective second-moment upper bounds
\[
 \frac{s_{\max}^2}{2},\quad s_w^2,\quad Cs_v^2,\quad
 Cs_v^2s_w^2X^2,\quad
 Cs_v^2s_{\max}^2+2\E\norm{\beta_0}^2+2\norm{\bar y}^2
\]
after dividing the first two squared norms by $M$. Take each squared
constant at least $7/\rho$ times the corresponding bound, enlarging
to a positive value if necessary. For $J_{\Theta,0}$, conditional
centering and independence of the entries of $V_0$ eliminate cross
terms in $V_0\Gamma_{n,0}W_0$; its expected squared Frobenius norm
is at most $Cs_v^2s_w^2$. The factors $x_n$ and $N^{-1/2}$ give the
stated bound. The residual bound follows by conditioning on $(W_0,b_0)$
and using the independent centered $V_0$; the output bias itself need
not be centered.
The variance calculation used in Corollary~\ref{cor:astra_init} gives
\[
 \Pr\{\lambda_{\min}(H_0H_0^T)<\kappa M\}
 \le\frac{6s_{\max}^4}{M\kappa_*^2}.
\]
Thus $M\ge42s_{\max}^4/(\rho\kappa_*^2)$ makes this sixth failure
probability at most $\rho/7$.

\emph{Boundary event and constants.}
Set
\[
 h=h_0+1,\quad \Omega=Z(C_v+1),\quad
 B=\sqrt{h^2+\Omega^2+1},\quad R_\phi=4BR_0/\kappa,
 \quad \tau_i=\frac{X\norm{w_{i,0}}+ZR_\phi}{\sqrt M}.
\]
For the seventh event define
\[
 \mathcal B=\sum_{i=1}^M\norm{V_{:i,0}}\norm{w_{i,0}}
 \mathbf1\{\min_n|w_{i,0}^Tz_n+b_{i,0}|\le\tau_i\}.
\]
Conditioning on $w_{i,0}$, the Gaussian density of $b_{i,0}$ is bounded
by $(s_b\sqrt{2\pi})^{-1}$. A union bound over $n$, independence
of $V_{:i,0}$, and $\E\norm{V_{:i,0}}\le s_v\sqrt{C/M}$ give
\[
 \E\mathcal B\le C_{\rm bd}
 :=\frac{2Ns_v\sqrt C}{s_b\sqrt{2\pi}}
 \left(X\E\norm w^2+ZR_\phi\E\norm w\right).
\]
Here $\E\norm w^2=s_w^2$ and $\E\norm w\le s_w$; either the actual
moment or its bound can define the constant. Put
$\mathcal B_*=7C_{\rm bd}/\rho$. Then
$\Pr\{\mathcal B>\mathcal B_*\}\le\rho/7$.
Define
\begin{align*}
 A&=A_0+X\{R_\phi C_w+(C_v+1)+\mathcal B_*\},&
 R_\theta&=4AR_0/\kappa,\\
 D_H&=XC_wR_\theta+ZR_\phi,&
 D_K&=(2h_0+1)D_H,\\
 C_{\rm gap}&=2D_K+2\Omega^2+A^2.
\end{align*}
It suffices to take the integer ceiling of
\begin{equation}\label{eq:astra_widththreshold}
 M_*\ge\max\left\{1,R_\theta^2,R_\phi^2,D_H,
              2D_K/\kappa,42s_{\max}^4/(\rho\kappa_*^2)\right\}.
\end{equation}
All constants are determined before $M$ is chosen. The intersection
of the five events in \eqref{eq:astra_fiveevents}, the head-kernel
event, and the boundary event has probability at least $1-\rho$.
Initialization at a kink has probability zero and may be excluded
without another probability charge.

\emph{Deterministic estimates in the bootstrap tube.}
Work on this intersection and in the closed tube
\[
 \norm{\Theta-\Theta_0}\le R_\theta/M,\qquad
 \norm{\phip-\phip_0}\le R_\phi/\sqrt M.
\]
The width conditions imply $\norm{\Theta-\Theta_0}\le1$ and
$\norm{\phip-\phip_0}\le1$, as well as
$R_\theta/M\le M^{-1/2}$. Write $\omega=(W,b)$.
ReLU is 1-Lipschitz, so
\begin{align*}
 \norm{H-H_0}_F
 &\le X\norm{W_0}_F\norm{\Theta-\Theta_0}
        +Z\norm{\omega-\omega_0}
 \le D_H/\sqrt M,\\
 \norm H_F&\le h\sqrt M,\qquad
 \norm{HH^T-H_0H_0^T}_{\rm op}\le D_K.
\end{align*}
For the second line use $D_H\le M$ and expand the product; in
particular $2h_0D_H+D_H^2/M\le D_K$.
Consequently
\begin{equation}\label{eq:astra_headfloor}
 HH^T\succeq(\kappa M/2)I_N.
\end{equation}
The preactivation change for neuron $i$ on any datum is bounded by
\[X\norm{w_{i,0}}R_\theta/M+Z\norm{\omega_i-\omega_{i,0}}\le\tau_i.\]
Thus a changed gate belongs to the initial boundary event above.
To control the encoder Jacobian use the exact three-term decomposition
\begin{equation}\label{eq:astra_three_terms}
 V\Gamma_nW-V_0\Gamma_{n,0}W_0
 =(V-V_0)\Gamma_nW_0+V\Gamma_n(W-W_0)
       +V_0(\Gamma_n-\Gamma_{n,0})W_0.
\end{equation}
The last term has operator norm at most $\mathcal B$. The first
uses the current gates and initial $W_0$; the second uses the current
$V$; both distinctions matter. Averaging over $n$ and multiplying by
$x_n$ yields
\[
 \norm{J_\Theta}_{\rm op}
 \le A_0+X\{\norm{V-V_0}_F\norm{W_0}_F
       +\norm V_F\norm{W-W_0}_F+\mathcal B_*\}\le A.
\]
The separate block bounds are
\[
 \norm{J_\omega}_{\rm op}\le Z\norm V_F\le\Omega,\qquad
 \norm{J_V}_{\rm op}\le h\sqrt M,\qquad
 \norm{J_\beta}_{\rm op}=1,
 \quad \norm{J_\phip}_{\rm op}\le B\sqrt M.
\]
Furthermore $J_\phip J_\phip^T\succeq
(HH^T)\otimes I_C\succeq(\kappa M/2)I_{NC}$.
These estimates also hold for the frozen path in its decoder tube.

\emph{Flow existence, chain rule and continuation.}
For completeness, local solutions in the convention above can be
obtained by replacing ReLU by
$\sigma_\epsilon(u)=(u+\sqrt{u^2+\epsilon^2})/2$.
For fixed $M$, on a compact parameter neighborhood its values and
first derivatives are uniformly bounded for $0<\epsilon\le1$.
The smooth gradient vector fields are therefore uniformly bounded
there, giving a common positive interval of existence and uniformly
Lipschitz paths on a smaller neighborhood. A subsequence converges
uniformly. The finitely many gates have a weak-* convergent subsequence
in $L^\infty$ with limits in $[0,1]$ and the correct values away from
zero preactivations. Each block gradient is linear in the corresponding
gate, with all remaining coefficients converging uniformly. Passing
to the integral equations gives one common limiting gate selection
and a local solution.
Along any such absolutely continuous solution, each preactivation
has derivative zero almost everywhere on its zero level set.
Consequently $d\sigma(g(t))/dt=\gamma(t)\dot g(t)$ almost everywhere,
for every admissible gate selection. The ordinary product rule then
gives $\dot e=-JJ^Te$ and
$\dot\loss=-\norm{J^Te}^2$ almost everywhere. These facts follow
directly here; related nonsmooth chain-rule frameworks are discussed
by \citet{astra_davis2020,astra_bolte2021}. We do not infer a Clarke
inclusion solely from the computational graph.

Until a first exit from the tube,
$\frac d{dt}\norm e^2\le-\kappa M\norm e^2$, proving (N1).
The integrated speeds are bounded by
\[
 \int_0^\infty\norm{\dot\Theta}\,dt\le\frac{2AR_0}{\kappa M}
     =\frac{R_\theta}{2M},\qquad
 \int_0^\infty\norm{\dot\phip}\,dt
     \le\frac{2BR_0}{\kappa\sqrt M}
     =\frac{R_\phi}{2\sqrt M}.
\]
At a first exit these bounds contradict the tube boundary. Thus no
exit occurs. Bounded velocities in this compact tube preclude finite
escape, give a limit at any finite endpoint, and the local existence
argument continues the path from that endpoint. This proves global
existence and (N1)--(N2) for every selection, without uniqueness.

\emph{Common linear reference and attribution.}
Let
\[
 K_0=(H_0H_0^T)\otimes I_C+
             (\mathbf1\mathbf1^T/N)\otimes I_C,
 \qquad e^{\rm lin}(t)=e^{-K_0t}e_0.
\]
The second summand is the constant output-bias kernel.
For the full output kernels $K^J(t)$ and $K^F(t)$,
\[
 \norm{K^J(t)-K_0}_{\rm op}\le D_K+\Omega^2+A^2,
 \qquad
 \norm{K^F(t)-K_0}_{\rm op}\le D_K+\Omega^2.
\]
Since $K_0\succeq\kappa MI$, variation of constants gives
\[
 e^a(t)-e^{\rm lin}(t)
 =-\int_0^t e^{-K_0(t-s)}(K^a(s)-K_0)e^a(s)\,ds,
 \qquad a\in\{J,F\}.
\]
Use $\norm{e^a(s)}\le R_0$ and
$\int_0^te^{-\kappa M(t-s)}ds\le1/(\kappa M)$ and add the two
bounds to obtain (N3). Both paths track the same linear reference;
(N3) is not a general stability theorem for arbitrary nonlinear
joint and frozen flows. Applying
Corollary~\ref{cor:astra_attribution} with $a=A^2$ and decoder floor
$\kappa M/2$ gives (N4), completing the proof.
\end{proof}

\begin{remark}[Threshold and scope]
The threshold \eqref{eq:astra_widththreshold} is deliberately
conservative. The boundary term contains an explicit factor $N/\rho$
and $R_0$ can scale as $\rho^{-1/2}$, before its further dependence
on $R_\phi$, the data, and $\kappa_*$ is substituted. When these
factors are nonzero and the remaining scales are held fixed, already
$R_\theta^2$ can contain a contribution of order
$N^2/(\rho^3\kappa_*^2)$; the full displayed construction can be
substantially worse. This is an illustration of vacuity, not a
universal necessary lower bound (degenerate inputs can remove some
terms), and it is not an empirical width prediction.
\end{remark}

\paragraph{Relation to prior work and a normalization control.}
Positive limiting kernels, small parameter movement, and control of
ReLU gate changes are established tools in wide-network convergence
proofs \citep{astra_du2019,astra_arora2019}. Linearized dynamics are
studied by \citet{astra_lee2019}; their dependence on parameterization
and the distinction between lazy and feature-learning limits are
central to \citet{chizat2019lazy,yang2021tensor}.
The divergent kernel scale of the naive standard parameterization and
its relation to layerwise training rates are also explicit in
\citet{astra_sohldickstein2020}.
The result above is a scoped instantiation with a coupled finite-width
encoder and a joint/frozen comparison. It does not establish a new
general lazy-training mechanism. Its gain disparity vanishes under
an explicit $1/\sqrt M$ readout in the following precise convention:
write $f=M^{-1/2}U\sigma(W\Theta x+b)+\beta$, initialize the entries
of $U$ with variance $s_v^2$, and train $U$ at unit Euclidean rate.
This preserves the initial function distribution but changes the
optimization metric. The readout kernel is then $HH^T/M=O(1)$,
and the encoder kernel remains $O(1)$; there is no forced factor-$M$
readout/encoder separation. This does not assert equality of the two
kernels or of their directional gains. Merely inserting the factor
while retaining variance $s_v^2/M$ for the trained readout would be
a different, shrinking-output initialization.

% BEGIN ASTRA BIBTEX: copy the entries below, without the leading percent, to references.bib.
% @inproceedings{astra_karimi2016,
%   author = {Karimi, Hamed and Nutini, Julie and Schmidt, Mark},
%   title = {Linear Convergence of Gradient and Proximal-Gradient Methods Under the {Polyak--Lojasiewicz} Condition},
%   booktitle = {European Conference on Machine Learning and Knowledge Discovery in Databases},
%   year = {2016}, url = {https://arxiv.org/abs/1608.04636}
% }
% @article{astra_tropp2012,
%   author = {Tropp, Joel A.},
%   title = {User-Friendly Tail Bounds for Sums of Random Matrices},
%   journal = {Foundations of Computational Mathematics}, volume = {12},
%   pages = {389--434}, year = {2012}, doi = {10.1007/s10208-011-9099-z}
% }
% @article{vanlaarhoven2017l2,
%   author = {van Laarhoven, Twan},
%   title = {{L2} Regularization versus Batch and Weight Normalization},
%   journal = {arXiv preprint arXiv:1706.05350}, year = {2017},
%   url = {https://arxiv.org/abs/1706.05350}
% }
% @inproceedings{astra_yun2021,
%   author = {Yun, Chulhee and Krishnan, Shankar and Mobahi, Hossein},
%   title = {A Unifying View on Implicit Bias in Training Linear Neural Networks},
%   booktitle = {International Conference on Learning Representations}, year = {2021},
%   url = {https://arxiv.org/abs/2010.02501}
% }
% @inproceedings{astra_moroshko2020,
%   author = {Moroshko, Edward and Woodworth, Blake and Gunasekar, Suriya and Lee, Jason D. and Srebro, Nathan and Soudry, Daniel},
%   title = {Implicit Bias in Deep Linear Classification: Initialization Scale vs Training Accuracy},
%   booktitle = {Advances in Neural Information Processing Systems}, volume = {33}, year = {2020},
%   url = {https://arxiv.org/abs/2007.06738}
% }
% @article{astra_berthier2023,
%   author = {Berthier, Rapha{\"e}l},
%   title = {Incremental Learning in Diagonal Linear Networks},
%   journal = {Journal of Machine Learning Research}, volume = {24}, number = {171},
%   pages = {1--26}, year = {2023}, url = {https://jmlr.org/papers/v24/22-1395.html}
% }
% @inproceedings{astra_saxe2014,
%   author = {Saxe, Andrew M. and McClelland, James L. and Ganguli, Surya},
%   title = {Exact Solutions to the Nonlinear Dynamics of Learning in Deep Linear Neural Networks},
%   booktitle = {International Conference on Learning Representations}, year = {2014},
%   url = {https://arxiv.org/abs/1312.6120}
% }
% @inproceedings{astra_du2018balance,
%   author = {Du, Simon S. and Hu, Wei and Lee, Jason D.},
%   title = {Algorithmic Regularization in Learning Deep Homogeneous Models: Layers are Automatically Balanced},
%   booktitle = {Advances in Neural Information Processing Systems}, volume = {31}, year = {2018},
%   url = {https://arxiv.org/abs/1806.00900}
% }
% @inproceedings{pezeshki2021gradient,
%   author = {Pezeshki, Mohammad and Kaba, Sekou-Oumar and Bengio, Yoshua and Courville, Aaron and Precup, Doina and Lajoie, Guillaume},
%   title = {Gradient Starvation: A Learning Proclivity in Neural Networks},
%   booktitle = {Advances in Neural Information Processing Systems}, volume = {34}, year = {2021},
%   url = {https://arxiv.org/abs/2011.09468}
% }
% @inproceedings{astra_du2019,
%   author = {Du, Simon S. and Zhai, Xiyu and Poczos, Barnabas and Singh, Aarti},
%   title = {Gradient Descent Provably Optimizes Over-parameterized Neural Networks},
%   booktitle = {International Conference on Learning Representations}, year = {2019},
%   url = {https://openreview.net/forum?id=S1eK3i09YQ}
% }
% @article{astra_davis2020,
%   author = {Davis, Damek and Drusvyatskiy, Dmitriy and Kakade, Sham and Lee, Jason D.},
%   title = {Stochastic Subgradient Method Converges on Tame Functions},
%   journal = {Foundations of Computational Mathematics}, volume = {20}, pages = {119--154},
%   year = {2020}, url = {https://arxiv.org/abs/1804.07795}
% }
% @article{astra_bolte2021,
%   author = {Bolte, J{\'e}r{\^o}me and Pauwels, Edouard},
%   title = {Conservative Set Valued Fields, Automatic Differentiation, Stochastic Gradient Methods and Deep Learning},
%   journal = {Mathematical Programming}, volume = {188}, pages = {19--51},
%   year = {2021}, doi = {10.1007/s10107-020-01501-5},
%   url = {https://arxiv.org/abs/1909.10300}
% }
% @inproceedings{astra_lee2019,
%   author = {Lee, Jaehoon and Xiao, Lechao and Schoenholz, Samuel S. and Bahri, Yasaman and Novak, Roman and Sohl-Dickstein, Jascha and Pennington, Jeffrey},
%   title = {Wide Neural Networks of Any Depth Evolve as Linear Models Under Gradient Descent},
%   booktitle = {Advances in Neural Information Processing Systems}, volume = {32}, year = {2019},
%   url = {https://arxiv.org/abs/1902.06720}
% }
% @inproceedings{astra_arora2019,
%   author = {Arora, Sanjeev and Du, Simon S. and Hu, Wei and Li, Zhiyuan and Wang, Ruosong},
%   title = {Fine-Grained Analysis of Optimization and Generalization for Overparameterized Two-Layer Neural Networks},
%   booktitle = {Proceedings of the 36th International Conference on Machine Learning},
%   series = {Proceedings of Machine Learning Research}, volume = {97}, pages = {322--332}, year = {2019},
%   url = {https://proceedings.mlr.press/v97/arora19a.html}
% }
% @inproceedings{chizat2019lazy,
%   author = {Chizat, Lenaic and Oyallon, Edouard and Bach, Francis},
%   title = {On Lazy Training in Differentiable Programming},
%   booktitle = {Advances in Neural Information Processing Systems}, volume = {32}, year = {2019},
%   url = {https://arxiv.org/abs/1812.07956}
% }
% @inproceedings{yang2021tensor,
%   author = {Yang, Greg and Hu, Edward J.},
%   title = {Tensor Programs {IV}: Feature Learning in Infinite-Width Neural Networks},
%   booktitle = {Proceedings of the 38th International Conference on Machine Learning},
%   series = {Proceedings of Machine Learning Research}, volume = {139}, pages = {11727--11737}, year = {2021},
%   url = {https://proceedings.mlr.press/v139/yang21c.html}
% }
% @article{astra_sohldickstein2020,
%   author = {Sohl-Dickstein, Jascha and Novak, Roman and Schoenholz, Samuel S. and Lee, Jaehoon},
%   title = {On the Infinite Width Limit of Neural Networks with a Standard Parameterization},
%   journal = {arXiv preprint arXiv:2001.07301}, year = {2020},
%   url = {https://arxiv.org/abs/2001.07301}
% }
% END ASTRA BIBTEX

\endgroup

\section{The effective gradient ratio}
\label{app:egr}

The theorems of the main text bound the spectral module's learning signal
but say nothing about how a practitioner would notice its decay inside a
single run. The \emph{effective gradient ratio} (EGR) is the cheapest
observable that tracks it. We define it, prove the one statement the theory
supports---decay of the numerator---and report what it does and does not
predict empirically. It appears here, and not in the main text, because the
evidence supports it as a complementary measurement rather than as a
decision rule.

\begin{definition}[Effective gradient ratio]
\label{def:egr}
At training step $t$,
\[
    \EGR(t) \;=\;
    \frac{\norm{\nabla_\thetap \loss(\thetap_t, \phip_t)}}
         {\norm{\nabla_\phip \loss(\thetap_t, \phip_t)}} .
\]
Both norms are available from the standard backward pass, so the cost is one
extra reduction per step. $\EGR$ is unitless and invariant to the shared
learning rate.
\end{definition}

\begin{proposition}[Spectral gradient decay]
\label{prop:egr_monotonic}
Under the hypotheses of Theorem~\ref{thm:twoscale}---%
Assumptions~\ref{ass:linear}, \ref{ass:subg}, \ref{ass:inputlip}
and~\ref{ass:residual}, with the trajectory contained in the regular set up
to a horizon $T$---write
$B_T := \sup_{t \leq T} \norm{J_\thetap(t)}_{\mathrm{op}}$. Then for every
$t \leq T$,
\[
    \norm{\nabla_\thetap \loss(t)} \;\leq\; \frac{B_T\,C}{1 + \mu t},
    \qquad
    B_T \;\leq\; \sqrt{C_\thetap}
    \;=\; \widetilde{L}\sqrt{\lambda_{\max}(\Sigma_X)},
\]
with $C_\thetap$ the cap of Lemma~\ref{lem:opcap} and $C, \mu$ the constants
of Assumption~\ref{ass:residual}.
\end{proposition}

\begin{proof}
$\norm{\nabla_\thetap \loss} = \norm{J_\thetap^\top r} \leq
\norm{J_\thetap}_{\mathrm{op}} \norm{r} \leq B_T \norm{r}$ for $t \leq T$,
and Assumption~\ref{ass:residual} bounds $\norm{r(t)} \leq C/(1 + \mu t)$.
The cap on $B_T$ is Lemma~\ref{lem:opcap}. The horizon is not a
technicality: in the unregularized separable regime that
Assumption~\ref{ass:residual} models, a cap uniform over
$t \in [0, \infty)$ cannot be assumed; in practice $T$ is the length of the
run, which is the horizon the diagnostic operates over.
\end{proof}

\begin{remark}[From gradient decay to the ratio]
\label{rem:egr_ratio_caveat}
Proposition~\ref{prop:egr_monotonic} bounds the \emph{numerator} only. The
denominator $\norm{\nabla_\phip \loss} = \norm{J_\phip^\top r}$ is also
proportional to the residual, so both quantities vanish together as the loss
is fitted and the behaviour of their ratio is not determined by residual
decay alone: it depends on the relative conditioning of the two Jacobians
along the trajectory. The EGR trajectory is therefore an empirical object,
and we use the proposition only for the absolute claim that the spectral
module's learning signal decays polynomially while the loss is being
fitted.
\end{remark}

\paragraph{Descriptive statistics.}
On the synthetic width sweep of Experiment~1.2 ($N = 24$ runs, widths
$\{16, 64, 256, 1024\}$, six seeds each), EGR depth correlates with final
test accuracy at Pearson $r = -0.723$ and Spearman $\rho = -0.808$, and with
the train--test gap at $r = +0.741$. The pooled correlation is confounded
with capacity: EGR depth correlates with $\log \mathrm{width}$ at
$r = +0.903$, and the partial correlation with test accuracy controlling for
$\log \mathrm{width}$ collapses to $-0.169$. A fixed-capacity companion
sweep isolates the remainder: width fixed at $D = 256$, varying data noise
$\{0.05, 0.10, 0.15, 0.20\}$, learning rate
$\{3{\times}10^{-4}, 10^{-3}, 3{\times}10^{-3}\}$, weight decay
$\{0, 10^{-3}\}$ and six seeds, $N = 144$ runs. There
$r(\EGR_{\min}, \mathrm{acc}) = +0.344$; $\EGR_{\min}$ is nearly orthogonal
to the swept hyperparameters (marginal correlations $+0.029$, $-0.073$,
$+0.007$; $R^2 = 0.015$ for the full set), and the partial correlation given
all three, $+0.395$, is \emph{larger} than the pooled value---the opposite
signature of a hidden hyperparameter confound. Controlling additionally for
final training loss (itself a predictor, $r = +0.47$ for accuracy) leaves it
at $+0.335$.

\paragraph{Why we report no aggregate test on the cell analysis.}
Within the 24 $(\mathrm{noise}, \mathrm{lr}, \mathrm{wd})$ cells in which
only the seed varies, 23 of 24 give a positive
$r(\EGR_{\min}, \mathrm{acc})$, with mean $+0.584$. We report this
descriptively and attach no $p$-value to it. The same six seeds populate
every cell, so the cells are not independent replicates; the between-cell
standard deviation of the Fisher-transformed correlations, $0.233$, sits far
below the $0.577$ that independent cells of this size would produce, and a
one-sample test over cells would therefore overstate the evidence by an
unknown factor. The $N = 144$ partial correlation above, computed once over
runs rather than over cells, is the inferential statement we stand behind.

\paragraph{Effect size, honestly.}
As a binary alarm, $\EGR_{\min}$ reaches ROC-AUC $0.612$ for flagging
below-median test accuracy and $\EGR_{\mathrm{depth}}$ only $0.538$ for a
high overfit gap; the best operating point ($\EGR_{\min} < 0.10$) catches
$90\%$ of poor runs while also flagging $57\%$ of healthy ones. At fixed
capacity, EGR is a real, capacity-independent signal correlated with
generalization, and it carries information that train-loss monitoring does
not; the within-capacity effect size is small enough that it belongs
alongside loss monitoring as a complementary measurement and not as a
standalone run-quality indicator.

\IfFileExists{sections_iclr/appendix_experiments.tex}{\section{Synthetic experiment details}
\label{app:synthetic}

This appendix records the settings, artifacts and complete tables behind the
synthetic measurements quoted in the main text. Every number below is taken
from a named file under \texttt{results/}; where a value is a mean, a median
or a fitted slope that we computed from the raw rows of such a file rather
than being printed in it, we say so at that place.

\subsection{Generator, architectures and common protocol}
\label{app:synthetic_setup}

The generator produces hyperspectral images $X \in \R^{H \times W \times S}$
with $H = W = 16$ and $S = 64$ and a two-class per-pixel label. A known
spectral direction $u \in \R^S$ is placed in each pixel, scaled by a per-pixel
latent $z \sim \mathcal{N}(0,1)$ and modulated by a scalar strength $\alpha$;
smooth per-class spatial templates, modulated by a scalar strength $\beta$,
enter the label-generating logits. At $\alpha = \beta = 1$ and noise level
$0.1$ the Bayes-optimal accuracy is approximately $0.85$--$0.95$.

The spectral module $f_\thetap$ is a linear per-pixel map $\R^{64} \to \R^{16}$
without bias ($\Cf = 1024$ throughout). Four spatial modules are used across
the experiments: \emph{SpatialMLP}, a per-pixel one-hidden-layer network of
width $M$ (no spatial receptive field, $\Cg = \Theta(M)$); \emph{SpatialDeepMLP},
three hidden layers all $M \times M$ ($\Cg = \Theta(M^2)$); \emph{SpatialCNN},
two $3 \times 3$ convolutional layers; and \emph{SpatialViT}, a two-layer
pre-LN pixel-token transformer. The per-pixel loss is softmax cross-entropy
averaged over sites.

Two initializations of the spatial module are distinguished throughout, and
the spectral module's initialization is never altered. \emph{Default} is
PyTorch's \texttt{nn.Linear} initialization, whose bias variance is
$\Theta(1/\mathrm{fan\text{-}in})$. \emph{Theorem} initialization draws weights
$\mathcal{N}(0, 2/\mathrm{fan\text{-}in})$ and biases $\mathcal{N}(0,
\sigma_b^2)$ with $\sigma_b = 0.5$ held fixed in width, which is what the
curvature-disparity theorem's head hypothesis requires. Three further
controls decouple the weight law from the bias law:
\texttt{default\_fixedbias} (default weights, fixed-variance biases),
\texttt{he\_scaledbias} (He weights, Gaussian biases at the default variance
$1/(3\,\mathrm{fan\text{-}in})$) and \texttt{he\_zerobias} (He weights, zero
biases).

\subsection{Experiment 1.1: curvature blocks at initialization}
\label{app:exp11}

\paragraph{What is measured.} For each (architecture, width, seed) we compute
$\lambda_{\max}$ of the two diagonal Gauss--Newton blocks $G_{\thetap\thetap}$
and $G_{\phip\phip}$ of the averaged cross-entropy loss at random
initialization, in float64, on one fixed batch of 64 images per seed, by
Lanczos with full reorthogonalization applied to a matrix-free GGN block
operator ($80$ Lanczos steps, top-1). Before any sweep the runner executes a
validation gate that cross-checks the operator's $\lambda_{\max}$ against a
densely materialized $J^\top H_\tau J$ eigendecomposition on tiny
configurations for both architectures and both blocks, and aborts unless the
relative error is below $10^{-6}$. This closes the loss-Hessian versus
Gauss--Newton measurement gap present in the first version of this experiment:
the quantity measured here is the Gauss--Newton block the theory is stated for.
Five seeds per width. Artifacts:
\texttt{results/exp1\_1\_v3\_ggn\_aggregated.csv},
\texttt{results/exp1\_1\_v3\_ggn\_theoreminit\_aggregated.csv} and
\texttt{results/exp1\_1\_v3\_ggn\_sigmab\_isolation\_aggregated.csv}, each with
per-width seed means and standard deviations.

\paragraph{Results.} Table~\ref{tab:app_exp11} lists the two shallow-head arms
and the five deep-head arms. Eigenvalue entries are the per-width means over
five seeds printed in those files; the log--log slopes are ordinary
least-squares fits of $\log \lambda$ on $\log M$ over the full width range,
computed by us from those per-width means.

\begin{table}[htbp]
\centering
\small
\setlength{\tabcolsep}{4pt}
\caption{Experiment 1.1 (v3): top Gauss--Newton block eigenvalues at
initialization, five seeds per width. Each eigenvalue is given at the smallest
and the largest width of the arm's range. Slopes are least-squares log--log
fits over that range, computed by us from the per-width seed means in the three
aggregated CSVs; they use three orders of magnitude of width for the shallow
head and two for the deep head, and are finite-range descriptions rather than
asymptotic exponents. The five deep-head rows are the five arms of the
initialization isolation.}
\label{tab:app_exp11}
\begin{tabular}{lrrrrrrr}
\toprule
& & \multicolumn{3}{c}{$\lambda_\thetap$} & \multicolumn{3}{c}{$\lambda_\phip$} \\
\cmidrule(lr){3-5}\cmidrule(lr){6-8}
Initialization & $M$ range & first & last & slope & first & last & slope \\
\midrule
\multicolumn{8}{l}{\emph{Shallow head} (one hidden layer, $\Cg = \Theta(M)$)} \\
default & $16$--$8192$ & $0.0184$ & $0.0220$ & $-0.019$ & $0.62$ & $46.0$ & $0.705$ \\
theorem & $16$--$8192$ & $0.691$ & $0.538$ & $-0.040$ & $2.14$ & $441.6$ & $0.895$ \\
\midrule
\multicolumn{8}{l}{\emph{Deep head} (three hidden $M\times M$ layers, $\Cg = \Theta(M^2)$)} \\
default & $16$--$1024$ & $0.00032$ & $0.00037$ & $+0.012$ & $0.68$ & $1.11$ & $0.105$ \\
theorem & $16$--$1024$ & $0.371$ & $0.398$ & $-0.006$ & $3.62$ & $243.1$ & $0.989$ \\
default, fixed bias & $16$--$1024$ & $0.00055$ & $0.00063$ & $-0.041$ & $1.67$ & $73.8$ & $0.924$ \\
He, scaled bias & $16$--$1024$ & $0.310$ & $0.343$ & $+0.030$ & $1.50$ & $27.6$ & $0.696$ \\
He, zero bias & $16$--$1024$ & $0.336$ & $0.372$ & $-0.050$ & $1.45$ & $17.6$ & $0.606$ \\
\bottomrule
\end{tabular}
\end{table}

\paragraph{Reading.} The spectral cap is flat in width in every arm: the seven
$\lambda_\thetap$ slopes lie in $[-0.050, +0.030]$, and the level shifts by
three orders of magnitude between initializations without acquiring a width
trend. On the shallow head under the theorem's initialization the spatial
block grows with slope $0.895$ from $2.14$ to $441.6$; under default fan-in
initialization the same architecture gives slope $0.705$. The deep head
separates the two mechanisms most sharply: slope $0.105$ under default
initialization against $0.989$ under theorem initialization. The isolation
arms show that the width-linear behaviour survives if \emph{either} the
weight-variance term or the bias-variance term of the head hypothesis is
width-independent --- fixed-variance biases with default weights reach slope
$0.924$, He weights with default-variance biases $0.696$ and He weights with
zero biases $0.606$ --- while default fan-in initialization, which makes both
terms vanish with width, does not produce it. The corresponding curvature
ratios $\Dcurv$, first and last width, are $39.8 \to 2405$ and $3.4 \to 869$
for the two shallow arms and $3167 \to 3060$, $18.1 \to 668$, $5552 \to
1.17\times10^{5}$, $5.7 \to 84.5$ and $4.6 \to 46.7$ for the five deep arms, in
table order. They illustrate why the slope and the level must be quoted
together, and why we do not use $\Dcurv$ to attribute differences across arms:
the deep default arm has the largest ratio of any arm and the smallest
$\lambda_\phip$ slope, because its $\lambda_\thetap$ is $10^{-4}$ rather than
$10^{-1}$.

\subsection{Experiment 1.2 (v4): joint and frozen trajectories}
\label{app:exp12v4}

Data as in Appendix~\ref{app:synthetic_setup} with $\alpha = \beta = 1.0$ and
noise $0.10$; $512$ training and $128$ test images; Adam at learning rate
$10^{-3}$, batch size $32$, $150$ epochs ($2400$ steps); three seeds
($42$,$43$,$44$). The frozen arm sets \texttt{requires\_grad} to false on the
spectral parameters, so the spectral gradient is identically zero there by
construction. Widths are architecture-specific: $D \in \{16, 64, 256, 1024\}$
for the CNN head and $D \in \{64, 128, 256\}$ for the ViT head.
Table~\ref{tab:app_exp12v4} gives the full grid; means and standard deviations
over the three seeds were computed by us from the per-run rows of
\texttt{results/exp1\_2v4\_summary.csv}.

\begin{table}[htbp]
\centering
\small
\setlength{\tabcolsep}{4pt}
\caption{Experiment 1.2 (v4): endpoint and peak test accuracy and endpoint
losses, mean $\pm$ sd over three seeds, computed from the per-run rows of
\texttt{results/exp1\_2v4\_summary.csv}. \emph{Peak} is the maximum over the
run's per-epoch test evaluations and is therefore selection-optimistic.
$\Cg$ is the spatial module's parameter count; $\Cf = 1024$ in every row.}
\label{tab:app_exp12v4}
\begin{tabular}{llrrcccc}
\toprule
Head & Arm & $D$ & $\Cg$ & Final test acc. & Peak test acc. & Final test loss & Final train loss \\
\midrule
CNN & joint  & $16$   & $2{,}610$     & $0.618 \pm 0.061$ & $0.682 \pm 0.055$ & $0.650 \pm 0.037$ & $0.453 \pm 0.002$ \\
CNN & frozen & $16$   & $2{,}610$     & $0.644 \pm 0.044$ & $0.688 \pm 0.054$ & $0.628 \pm 0.038$ & $0.465 \pm 0.004$ \\
CNN & joint  & $64$   & $10{,}434$    & $0.574 \pm 0.038$ & $0.703 \pm 0.021$ & $0.698 \pm 0.035$ & $0.421 \pm 0.004$ \\
CNN & frozen & $64$   & $10{,}434$    & $0.611 \pm 0.042$ & $0.688 \pm 0.054$ & $0.658 \pm 0.036$ & $0.437 \pm 0.004$ \\
CNN & joint  & $256$  & $41{,}730$    & $0.540 \pm 0.023$ & $0.712 \pm 0.006$ & $1.070 \pm 0.112$ & $0.290 \pm 0.009$ \\
CNN & frozen & $256$  & $41{,}730$    & $0.574 \pm 0.026$ & $0.690 \pm 0.053$ & $0.787 \pm 0.029$ & $0.345 \pm 0.008$ \\
CNN & joint  & $1024$ & $166{,}914$   & $0.534 \pm 0.013$ & $0.711 \pm 0.007$ & $2.334 \pm 0.169$ & $0.029 \pm 0.004$ \\
CNN & frozen & $1024$ & $166{,}914$   & $0.560 \pm 0.024$ & $0.690 \pm 0.051$ & $1.119 \pm 0.054$ & $0.153 \pm 0.019$ \\
\midrule
ViT & joint  & $64$   & $84{,}546$    & $0.533 \pm 0.021$ & $0.707 \pm 0.008$ & $3.121 \pm 0.131$ & $0.136 \pm 0.030$ \\
ViT & frozen & $64$   & $84{,}546$    & $0.547 \pm 0.050$ & $0.702 \pm 0.037$ & $3.387 \pm 0.233$ & $0.132 \pm 0.024$ \\
ViT & joint  & $128$  & $300{,}162$   & $0.546 \pm 0.047$ & $0.706 \pm 0.019$ & $3.149 \pm 0.936$ & $0.039 \pm 0.016$ \\
ViT & frozen & $128$  & $300{,}162$   & $0.567 \pm 0.038$ & $0.679 \pm 0.057$ & $3.302 \pm 0.451$ & $0.025 \pm 0.010$ \\
ViT & joint  & $256$  & $1{,}124{,}610$ & $0.503 \pm 0.050$ & $0.707 \pm 0.023$ & $3.779 \pm 0.547$ & $0.031 \pm 0.020$ \\
ViT & frozen & $256$  & $1{,}124{,}610$ & $0.539 \pm 0.007$ & $0.700 \pm 0.035$ & $3.522 \pm 0.160$ & $0.015 \pm 0.001$ \\
\bottomrule
\end{tabular}
\end{table}

Three observations, all descriptive at $n = 3$. The frozen arm's endpoint
accuracy is ahead of the joint arm's at every width and both heads, by
$+0.014$ to $+0.038$; every one of those gaps is smaller than the sum of the
two seed standard deviations, so none is resolvable here. Peak accuracies are
indistinguishable between arms (the largest difference is $0.028$, at ViT
$D = 128$, in favour of the joint arm). The endpoint test loss of the joint
CNN arm diverges with width, reaching $2.334 \pm 0.169$ at $D = 1024$ against
$1.119 \pm 0.054$ frozen, while its training loss falls to $0.029$; the ViT
runs are in a memorization regime at every width, with training losses of
$0.015$--$0.14$ and test losses above $3$, so gradient-ratio observables are
uninformative there.

\subsection{Experiment 1.2 (v5): residual overlap with the fast kernel subspace}
\label{app:exp12v5}

\paragraph{Model and protocol.} This experiment combines the model of the
verified shallow instance with the training protocol of
Appendix~\ref{app:exp12v4}. The model is the linear per-pixel map
$\R^{64} \to \R^{16}$ followed by a two-layer ReLU head of width $D$ under
theorem initialization with $\sigma_b = 0.5$; that head initialization is
required by the verified instance and is not PyTorch's default. Optimizer,
data, loss and schedule are exactly those of Appendix~\ref{app:exp12v4}, with
$D \in \{128, 512, 2048\}$, seeds $42$--$44$ and both a \texttt{joint} and a
\texttt{frozen} arm. The frozen arm here keeps \texttt{requires\_grad} true on
the spectral parameters but excludes them from the optimizer's parameter list;
the trajectory is identical to a hard freeze, and the logged spectral gradient
is the counterfactual signal a frozen encoder would have received. Artifacts:
\texttt{results/exp1\_2v5\_residual\_export.csv} ($198$ rows) and
\texttt{results/exp1\_2v5\_residual\_export\_REPORT.md}.

\paragraph{Probe and normalization.} All kernel quantities are computed on a
fixed probe of $N = 256$ supervised sites with $C = 2$ classes, so the probe
logit space is $\R^{512}$. The probe is drawn once per run with a fixed seed
from the first $16$ training images and held fixed for the whole trajectory;
non-probe pixels are set to the ignore index, so the GGN machinery sees the
same site set. With $z = N^{-1/2}\,\mathrm{stack}(\hat y)$ and averaged
cross-entropy, $r = N^{-1/2}(\mathrm{softmax}(z) - \mathrm{onehot}(y))$,
$J_b = N^{-1/2}\,\partial \hat y/\partial b$ and $K_b = J_bJ_b^\top$, so
$\nabla_b\loss = J_b^\top r$ holds exactly. We write
$a_{\mathrm{eff}} = r^\top K_\thetap r / \norm{r}^2$ and
$\kappa_{\mathrm{eff}} = r^\top K_\phip r / \norm{r}^2$ for the realised
Rayleigh quotients, $f_k$ for the fraction of $\norm{r}^2$ inside the top-$k$
eigenspace of $K_\phip$, and $\mathrm{share} =
a_{\mathrm{eff}}/(a_{\mathrm{eff}} + \kappa_{\mathrm{eff}})$, which coincides
by construction with the instantaneous gradient-flow attribution bound.
Table~\ref{tab:app_v5_ident} gives the three identities the runner
recomputes at each of the $198$ measurements.

\begin{table}[htbp]
\centering
\caption{Experiment 1.2 (v5): validation identities, maximum relative error
over all $198$ measurements, from
\texttt{results/exp1\_2v5\_residual\_export\_REPORT.md} \S2. All three are at
float64 round-off, so the kernels here are the same objects the GGN machinery
measures, in the same normalization.}
\label{tab:app_v5_ident}
\begin{tabular}{lr}
\toprule
Identity & Max.\ relative error \\
\midrule
$a_{\mathrm{eff}}\norm{r}^2 = \norm{\nabla_\thetap\loss}^2$ (autograd) & $1.30 \times 10^{-15}$ \\
$\kappa_{\mathrm{eff}}\norm{r}^2 = \norm{\nabla_\phip\loss}^2$ (autograd) & $1.62 \times 10^{-14}$ \\
$\lambda_{\max}(H^{1/2}K_\phip H^{1/2}) = \lambda_{\max}(G_{\phip\phip})$ (power iteration) & $2.47 \times 10^{-15}$ \\
\bottomrule
\end{tabular}
\end{table}

\paragraph{Headline table.} Table~\ref{tab:app_v5_head} reproduces the joint
arm of \S3 of the report at four log-spaced checkpoints for every width and
seed. The worst-case-coercivity proxy
$\lambda_{\max}(K_\thetap)/(\lambda_{\max}(K_\thetap)+\lambda_{\min}(K_\phip))$
is between $0.9998$ and $1.0000$ at every one of the $198$ rows and is omitted
from the table; $\lambda_{\min}(K_\phip)$ ranges from $2.0 \times 10^{-6}$ to
$1.8 \times 10^{-4}$ and is likewise omitted. Full trajectories at all eleven
checkpoints for all eighteen runs, and the frozen arm, are in the report.

\begin{table}[htbp]
\centering
\footnotesize
\setlength{\tabcolsep}{2.5pt}
\caption{Experiment 1.2 (v5), joint arm: probe measurements at four
checkpoints, every width and seed, from
\texttt{results/exp1\_2v5\_residual\_export\_REPORT.md} \S3. \emph{share} is
$a_{\mathrm{eff}}/(a_{\mathrm{eff}}+\kappa_{\mathrm{eff}})$; \emph{proxy} is
$\lambda_{\max}(K_\thetap)/(\lambda_{\max}(K_\thetap)+\lambda_{\max}(K_\phip))$.
These are probe quantities at eleven checkpoints, not a certificate about the
intervening steps.}
\label{tab:app_v5_head}
\begin{tabular}{rrrrrrrrrrrr}
\toprule
$D$ & seed & step & $a_{\mathrm{eff}}$ & $\kappa_{\mathrm{eff}}$ & $\lambda_{\max}(K_\thetap)$ & $\lambda_{\max}(K_\phip)$ & $f_1$ & $f_{20}$ & $f_{100}$ & share & proxy \\
\midrule
$128$ & $42$ & $0$    & $0.704$ & $0.858$   & $2.393$  & $19.839$  & $1.80\cdot10^{-2}$ & $0.3441$ & $0.4356$ & $0.4510$ & $0.1076$ \\
$128$ & $42$ & $32$   & $0.017$ & $0.011$   & $4.887$  & $21.590$  & $1.48\cdot10^{-5}$ & $0.0390$ & $0.1223$ & $0.6022$ & $0.1846$ \\
$128$ & $42$ & $426$  & $0.012$ & $0.022$   & $3.764$  & $21.884$  & $9.52\cdot10^{-6}$ & $0.0594$ & $0.1296$ & $0.3513$ & $0.1468$ \\
$128$ & $42$ & $2400$ & $0.025$ & $0.043$   & $6.965$  & $22.678$  & $3.68\cdot10^{-6}$ & $0.0288$ & $0.1778$ & $0.3688$ & $0.2350$ \\
$128$ & $43$ & $0$    & $0.299$ & $4.219$   & $1.532$  & $21.735$  & $1.34\cdot10^{-1}$ & $0.4465$ & $0.5173$ & $0.0662$ & $0.0659$ \\
$128$ & $43$ & $32$   & $0.013$ & $0.060$   & $3.275$  & $23.963$  & $1.10\cdot10^{-3}$ & $0.0375$ & $0.1516$ & $0.1755$ & $0.1202$ \\
$128$ & $43$ & $426$  & $0.009$ & $0.018$   & $2.229$  & $24.267$  & $1.81\cdot10^{-4}$ & $0.0298$ & $0.1404$ & $0.3405$ & $0.0841$ \\
$128$ & $43$ & $2400$ & $0.036$ & $0.037$   & $8.669$  & $26.355$  & $2.49\cdot10^{-5}$ & $0.0525$ & $0.1913$ & $0.4927$ & $0.2475$ \\
$128$ & $44$ & $0$    & $0.711$ & $0.948$   & $3.793$  & $16.447$  & $5.30\cdot10^{-3}$ & $0.3073$ & $0.3813$ & $0.4285$ & $0.1874$ \\
$128$ & $44$ & $32$   & $0.049$ & $0.078$   & $4.917$  & $17.693$  & $1.53\cdot10^{-4}$ & $0.0400$ & $0.1711$ & $0.3826$ & $0.2175$ \\
$128$ & $44$ & $426$  & $0.022$ & $0.026$   & $5.010$  & $17.847$  & $8.86\cdot10^{-6}$ & $0.0307$ & $0.1702$ & $0.4611$ & $0.2192$ \\
$128$ & $44$ & $2400$ & $0.029$ & $0.042$   & $6.428$  & $18.976$  & $1.04\cdot10^{-4}$ & $0.0309$ & $0.2020$ & $0.4073$ & $0.2530$ \\
\midrule
$512$ & $42$ & $0$    & $1.098$ & $4.107$   & $2.208$  & $70.380$  & $1.28\cdot10^{-2}$ & $0.5162$ & $0.5947$ & $0.2109$ & $0.0304$ \\
$512$ & $42$ & $32$   & $0.027$ & $0.071$   & $8.021$  & $72.767$  & $1.22\cdot10^{-4}$ & $0.0513$ & $0.1544$ & $0.2763$ & $0.0993$ \\
$512$ & $42$ & $426$  & $0.021$ & $0.041$   & $7.199$  & $72.745$  & $2.62\cdot10^{-5}$ & $0.0402$ & $0.1718$ & $0.3393$ & $0.0901$ \\
$512$ & $42$ & $2400$ & $0.065$ & $0.094$   & $24.622$ & $72.428$  & $4.01\cdot10^{-4}$ & $0.0154$ & $0.1741$ & $0.4101$ & $0.2537$ \\
$512$ & $43$ & $0$    & $0.464$ & $7.095$   & $1.717$  & $67.270$  & $3.67\cdot10^{-2}$ & $0.4456$ & $0.5469$ & $0.0614$ & $0.0249$ \\
$512$ & $43$ & $32$   & $0.023$ & $0.052$   & $4.829$  & $69.680$  & $1.14\cdot10^{-6}$ & $0.0334$ & $0.1328$ & $0.3050$ & $0.0648$ \\
$512$ & $43$ & $426$  & $0.013$ & $0.307$   & $3.550$  & $70.279$  & $2.90\cdot10^{-3}$ & $0.0369$ & $0.1388$ & $0.0413$ & $0.0481$ \\
$512$ & $43$ & $2400$ & $0.077$ & $0.085$   & $8.980$  & $70.889$  & $9.68\cdot10^{-6}$ & $0.0374$ & $0.2226$ & $0.4732$ & $0.1124$ \\
$512$ & $44$ & $0$    & $0.218$ & $14.482$  & $1.749$  & $61.928$  & $2.33\cdot10^{-2}$ & $0.4084$ & $0.4749$ & $0.0148$ & $0.0275$ \\
$512$ & $44$ & $32$   & $0.050$ & $0.236$   & $5.380$  & $64.643$  & $2.66\cdot10^{-4}$ & $0.0460$ & $0.1826$ & $0.1734$ & $0.0768$ \\
$512$ & $44$ & $426$  & $0.019$ & $0.064$   & $4.804$  & $65.177$  & $2.64\cdot10^{-6}$ & $0.0253$ & $0.1555$ & $0.2328$ & $0.0686$ \\
$512$ & $44$ & $2400$ & $0.073$ & $0.113$   & $9.265$  & $68.294$  & $1.12\cdot10^{-5}$ & $0.0144$ & $0.1853$ & $0.3935$ & $0.1195$ \\
\midrule
$2048$ & $42$ & $0$    & $0.525$ & $6.973$   & $1.805$  & $253.786$ & $8.38\cdot10^{-6}$ & $0.3850$ & $0.4677$ & $0.0701$ & $0.0071$ \\
$2048$ & $42$ & $32$   & $0.044$ & $0.324$   & $10.807$ & $254.624$ & $9.03\cdot10^{-4}$ & $0.0373$ & $0.1713$ & $0.1197$ & $0.0407$ \\
$2048$ & $42$ & $426$  & $0.026$ & $0.361$   & $11.223$ & $255.856$ & $9.58\cdot10^{-4}$ & $0.0385$ & $0.1612$ & $0.0663$ & $0.0420$ \\
$2048$ & $42$ & $2400$ & $0.196$ & $0.320$   & $41.018$ & $252.457$ & $2.56\cdot10^{-7}$ & $0.0206$ & $0.1848$ & $0.3796$ & $0.1398$ \\
$2048$ & $43$ & $0$    & $0.179$ & $148.674$ & $1.579$  & $276.775$ & $2.17\cdot10^{-1}$ & $0.6836$ & $0.7327$ & $0.0012$ & $0.0057$ \\
$2048$ & $43$ & $32$   & $0.049$ & $0.354$   & $9.486$  & $274.651$ & $4.62\cdot10^{-4}$ & $0.0193$ & $0.1151$ & $0.1217$ & $0.0334$ \\
$2048$ & $43$ & $426$  & $0.022$ & $1.213$   & $5.681$  & $276.344$ & $3.30\cdot10^{-3}$ & $0.0270$ & $0.1585$ & $0.0175$ & $0.0201$ \\
$2048$ & $43$ & $2400$ & $0.153$ & $0.179$   & $17.424$ & $281.103$ & $8.36\cdot10^{-6}$ & $0.0299$ & $0.1326$ & $0.4617$ & $0.0584$ \\
$2048$ & $44$ & $0$    & $0.495$ & $63.759$  & $2.389$  & $253.499$ & $4.96\cdot10^{-2}$ & $0.4682$ & $0.5447$ & $0.0077$ & $0.0093$ \\
$2048$ & $44$ & $32$   & $0.018$ & $0.257$   & $8.093$  & $252.495$ & $1.78\cdot10^{-4}$ & $0.0237$ & $0.1879$ & $0.0648$ & $0.0311$ \\
$2048$ & $44$ & $426$  & $0.034$ & $1.025$   & $8.764$  & $257.769$ & $1.40\cdot10^{-3}$ & $0.0264$ & $0.1860$ & $0.0323$ & $0.0329$ \\
$2048$ & $44$ & $2400$ & $0.131$ & $0.576$   & $30.982$ & $257.813$ & $1.10\cdot10^{-3}$ & $0.0263$ & $0.1584$ & $0.1848$ & $0.1073$ \\
\bottomrule
\end{tabular}
\end{table}

\paragraph{Cumulative gradient-energy fractions.} Table~\ref{tab:app_v5_cum}
gives $\sum_t \norm{\nabla_\thetap\loss}^2 / \sum_t (\norm{\nabla_\thetap
\loss}^2 + \norm{\nabla_\phip\loss}^2)$ over the whole $2400$-step run, from
the per-step logs the runner writes.

\begin{table}[htbp]
\centering
\caption{Experiment 1.2 (v5): cumulative gradient-energy fraction over the
full run, from \texttt{results/exp1\_2v5\_residual\_export\_REPORT.md} \S6.
These are mini-batch gradient norms over the $32$-image training batches
($8192$ sites), not probe quantities, and the optimizer is Adam, so this is a
gradient-energy attribution and not a decomposition of the realised loss
decrease. On the frozen arm the spectral parameters never move, so that column
is a counterfactual ratio.}
\label{tab:app_v5_cum}
\begin{tabular}{rrrr}
\toprule
$D$ & seed & joint & frozen \\
\midrule
$128$  & $42$ & $0.2896$ & $0.7256$ \\
$128$  & $43$ & $0.0806$ & $0.1883$ \\
$128$  & $44$ & $0.3432$ & $0.6091$ \\
$512$  & $42$ & $0.1303$ & $0.3308$ \\
$512$  & $43$ & $0.0618$ & $0.1230$ \\
$512$  & $44$ & $0.0415$ & $0.1336$ \\
$2048$ & $42$ & $0.0173$ & $0.0642$ \\
$2048$ & $43$ & $0.0042$ & $0.0059$ \\
$2048$ & $44$ & $0.0124$ & $0.0310$ \\
\bottomrule
\end{tabular}
\end{table}

\paragraph{What the numbers say.} The following are the report's own
summaries, over the joint arm's $99$ rows unless stated otherwise. The spatial
kernel is extremely ill-conditioned on the probe logit space:
$\lambda_{\max}/\lambda_{\min}$ has median $1.28 \times 10^{6}$ (range
$2.42\times10^{5}$ to $1.06\times10^{7}$), so a coercivity hypothesis
$K_\phip \succeq \kappa I$ taken on the whole output space is available only
with $\kappa = \lambda_{\min}$ and yields nothing. The visited residual is not
in that near-null part --- $\kappa_{\mathrm{eff}}$ exceeds $\lambda_{\min}$ by
between $225\times$ and a median $5.0 \times 10^{3}\times$ --- and it is not in
the coercive top either: $\kappa_{\mathrm{eff}}$ is a median $0.0025$ of
$\lambda_{\max}$ (range $2.5\times10^{-4}$ to $0.537$), and after step $75$ the
top eigendirection carries a median $8.6\times10^{-5}$ of $\norm{r}^2$ (down
from a median $0.0233$ at $t=0$), the top-$20$ eigenspace a median $0.0298$
(down from $0.446$) and the top-$100$ of $512$ a median $0.1612$. The spectral
side collapses the same way: the top left singular direction of $J_\thetap$
carries a median $6.2\times10^{-4}$ of $\norm{r}^2$ after step $75$, down from
$0.1513$ at $t=0$. Measured against the realised share, the top-eigenvalue
proxy understates the encoder's share in $85\%$ of joint rows, by a median
factor $2.72$ and up to $11.6\times$, and overstates it by up to $11.1\times$
in the remaining $15\%$; it errs in both directions and bounds the share in
neither. This is how we read those rows: the diagnostic limitation illustrated
by the counterexample to top-eigenvalue certificates is observed on this
instance.

\paragraph{Scope of these measurements.} Four restrictions apply and are not
optional caveats. (i) The cumulative column is a \emph{gradient-energy
fraction} on the same sequence, not an attribution of Adam's loss decrease;
the algebraic weighted-sum identity holds for any optimizer on one sequence,
but the realised displacement under Adam is not proportional to the gradient.
(ii) Quantities (i)--(v) of the report are computed on the fixed $256$-site
probe while the cumulative column uses training mini-batches: these are the
same functional on different sequences, and we report probe results as probe
results. Eleven checkpoints cannot certify an almost-everywhere trajectory
condition or the intervening steps. (iii) The appropriate isotropic reference
for the subspace fractions lives in the per-site class-contrast space: with
$\Pi = I_N \otimes (I_C - \mathbf{1}\mathbf{1}^\top/C)$, a random unit residual
there has expected projected mass $\mathrm{tr}(P\Pi)/[N(C-1)]$ rather than
$\mathrm{rank}(P)/(NC)$, and for $C = 2$ the difference is consequential. That
reference has not been recomputed for the projectors actually used, so we make
no claim that any measured mass falls below an isotropic level; the measured
masses themselves stand. (iv) The cross-entropy Hessian weighting is a
different quantity, not a correction to a squared-loss approximation: for
cross-entropy $r^\top K_b r = \norm{\nabla_b\loss}^2$ already holds exactly.
For completeness, \S5 of the report records that the weighted share moves by
at most a factor $2.721$ and at least $0.348$ relative to the unweighted one on
this instance. Finally, a cumulative $O(1/D)$ law does not follow from these
three widths: under the phase-restricted attribution result of the mathematics
appendix, a total share of order $1/D$ requires the complementary phase to
carry an $O(1/D)$ fraction of the energy, and that fraction has not been
measured here.

\subsection{Experiment 1.3: the equal-information calibration}
\label{app:exp13}

\paragraph{Calibration.} ``Equal information'' between the spectral pathway
(per-pixel content of $X$ along a direction $u$) and the spatial pathway
(content along an orthogonal direction $v \perp u$, modulated by a smooth
per-position function of the class) requires an operational definition. Three
were compared: \texttt{bayes}, which sets $\alpha, \beta$ so that a
single-modality logistic regression on $\alpha$-only data and a
position-majority classifier on $\beta$-only data both reach a target accuracy
of $0.75$; \texttt{ntk}, which equalizes the leading Gram-matrix eigenvalues of
the two modalities' per-pixel features; and \texttt{margin}, which equalizes
the $L_2$ margins of the two single-modality logistic regressions.

\paragraph{Setup and results.} For each mode the CNN composition is trained at
$D = 256$ for $150$ epochs on $N_{\mathrm{train}} = 512$ samples over three
seeds, in four conditions: joint, frozen at random spectral initialization,
spectral-only ($\beta = 0$) and spatial-only ($\alpha = 0$).
Table~\ref{tab:app_calib} compares the three modes.

\begin{table}[htbp]
\centering
\caption{Experiment 1.3: calibration-mode comparison (CNN, $D = 256$, three
seeds). \emph{spec\,/\,spat} are the achievable single-modality accuracies;
\emph{gap} is $a_{\mathrm{frozen}} - a_{\mathrm{joint}}$.}
\label{tab:app_calib}
\begin{tabular}{lccccc}
\toprule
Mode & $\alpha$, $\beta$ & spec\,/\,spat & joint & frozen & gap \\
\midrule
\texttt{bayes}  & $0.82,\ 18.75$ & $0.75\,/\,0.75$ & $0.48$ & $0.61$ & $+0.13$ \\
\texttt{ntk}    & $0.99,\ 4.91$  & $0.78\,/\,0.58$ & $0.53$ & $0.57$ & $+0.04$ \\
\texttt{margin} & $5.00,\ 5.00$  & $0.95\,/\,0.59$ & $0.60$ & $0.69$ & $+0.09$ \\
\bottomrule
\end{tabular}
\end{table}

Under \texttt{bayes} calibration the single-modality baselines are equal by
construction ($\alpha = 0.82$, $\beta = 18.75$ give spectral-only $0.746$ and
spatial-only $0.754$ under the logistic-regression and position-majority
bounds). On the full CNN model, spectral-only training reaches $0.53$ and
spatial-only training reaches $0.53$; frozen-spectral training on the full data
reaches $0.61$; joint training reaches $0.48$.

\paragraph{Two caveats bound this result.} First, the effect size is modest in
absolute terms. The quoted per-condition standard error of roughly $0.02$ is a
pooled figure over three seeds and $128$ test samples and does not include
between-seed variance, so it understates the uncertainty; at three seeds the
intervals are loose, and $0.48$ is not distinguishable from the $0.50$ chance
floor. Second, the CNN reaches only $0.53$ on single-modality data against a
$0.75$ calibration target, so the network is operating far from the Bayes
bound and part of the joint-training deficit could in principle reflect
generic overfitting with more trainable parameters. What the direction of the
result establishes is that joint training here undershoots baselines that
demonstrably use only one modality, which is not what a reading in which the
spectral features are simply irrelevant would produce.

\subsection{Experiment 1.6: comparison with a logit-magnitude penalty}
\label{app:exp16}

\paragraph{Design.} $120$ runs at width $256$ for $100$ epochs with weight
decay $0$, sweeping the conditions \texttt{joint}, \texttt{frozen} and
\texttt{sd} (spectral decoupling, $(\lambda/2)\norm{\text{logits}}^2$ on the
pre-softmax logits, summed over classes and averaged over examples, following
\citet{pezeshki2021gradient}) with $\lambda \in \{10^{-3}, 10^{-2}, 10^{-1},
1\}$, over noise levels $\{0.05, 0.10, 0.15, 0.20\}$ and five seeds.
Artifacts: \texttt{results/exp1\_6\_summary.csv} ($120$ rows),
\texttt{results/exp1\_6\_notes.md} and per-run logs under
\texttt{results/exp1\_6/}.

\paragraph{Results.} Table~\ref{tab:app_exp16} aggregates over the five seeds
and four noise levels; the means were computed by us from the rows of
\texttt{results/exp1\_6\_summary.csv}, with $\lambda$ selected per noise level
by highest mean final test accuracy. Paired Wilcoxon tests matched on
$(\text{noise}, \text{seed})$, $n = 20$, recomputed from the same rows, give
$T = 88$, $p = 0.55$ for joint against the selected-$\lambda$ penalty, and
$T = 6$, $p = 2.7\times10^{-5}$ for joint against frozen.

\begin{table}[htbp]
\centering
\caption{Experiment 1.6: peak and final test accuracy and their difference,
averaged over five seeds and four noise levels, computed from
\texttt{results/exp1\_6\_summary.csv}. The penalty row uses the $\lambda$ with
the highest mean final accuracy at each noise level ($10^{-2}$ at noise
$0.05$, $10^{-1}$ at the other three).}
\label{tab:app_exp16}
\begin{tabular}{lccc}
\toprule
Condition & Peak acc. & Final acc. & Peak $-$ final \\
\midrule
joint                       & $0.658$ & $0.548$ & $0.111$ \\
logit penalty, selected $\lambda$ & $0.660$ & $0.550$ & $0.110$ \\
frozen                      & $0.583$ & $0.517$ & $0.066$ \\
\bottomrule
\end{tabular}
\end{table}

\paragraph{Caveats.} This experiment was asserted in an earlier draft of the
work without being presented; it is set out here in full, and the following
qualifications belong with it. (i) The frozen condition trades peak accuracy
for a smaller peak-to-final difference: its mean peak is $0.583$ against
$0.660$, because a random spectral projection sometimes destroys signal.
(ii) Its seed-to-seed variability is large --- the standard deviation of final
accuracy across seeds is $0.085$--$0.141$ for frozen against $0.018$--$0.040$
for the other two conditions --- and one seed at noise $0.20$ ends at final
accuracy $0.372$, below chance. (iii) At width $256$ the joint condition does
not collapse dramatically ($0.66$ peak to $0.55$ final); larger drops appear at
$D = 2048$ in Experiment 1.2, so this is a moderate-capacity comparison.
(iv) The best penalty strength depends on the noise level, so a single
$\lambda$ across noise levels is not the best available setting for that
intervention. We report the comparison descriptively, as a contrast between
two interventions on one synthetic instance; nothing here is evaluated as a
recommendation for practice.

\subsection{Solvable serial instance: numerical verification}
\label{app:toy}

\paragraph{Zero-context instance.} Table~\ref{tab:app_toy} verifies the closed
forms and the bounds of the solvable serial cross-entropy instance at
$a_0 = 0$, $v_0 = 1$ and margin $m = \log 19 = 2.944$, from
\texttt{results/toy\_serial\_ce.csv}. The solver column is a stiff ODE
integration of the full $(2+M)$-dimensional state; a separate explicit-Euler
gradient-descent path, run at four step sizes, has observed convergence order
$1.0000004$ and agrees with an independent numpy analytic-gradient
implementation to $0$ relative error.

\begin{table}[htbp]
\centering
\small
\caption{Solvable serial instance ($a_0 = 0$, $v_0 = 1$, $m = \log 19$):
closed form against ODE solver, from \texttt{results/toy\_serial\_ce.csv}. The
last block is the two-timescale displacement bound against the realised
encoder displacement. Every relative discrepancy between closed form and
solver in this table is at most $1.2\times10^{-13}$.}
\label{tab:app_toy}
\begin{tabular}{lrrrr}
\toprule
Quantity & $M = 16$ & $M = 64$ & $M = 256$ & $M = 1024$ \\
\midrule
$\lambda_\thetap(0)$ & $0.25$ & $0.25$ & $0.25$ & $0.25$ \\
$\lambda_\phip(0)$ & $4.0$ & $16.0$ & $64.0$ & $256.0$ \\
$\Dcurv(0)$ & $16.0$ & $64.0$ & $256.0$ & $1024.0$ \\
\midrule
$a(T_m)$ & $0.142325$ & $0.042094$ & $0.011216$ & $0.002857$ \\
suppression $m/a_M$ & $20.688$ & $69.949$ & $262.53$ & $1030.7$ \\
$M\,a_M$ (limit $m$) & $2.2772$ & $2.6940$ & $2.8712$ & $2.9253$ \\
$T_m$ & $0.90784$ & $0.28680$ & $0.078763$ & $0.020252$ \\
$T_m$ lower bound $\Psi(m)/(M{+}1{+}2m^2)$ & $0.60992$ & $0.25437$ & $0.076345$ & $0.020094$ \\
$T_m$ upper bound $\Psi(m)/(M{+}1)$ & $1.2320$ & $0.32222$ & $0.081496$ & $0.020434$ \\
\midrule
$\sup_t\,(q_J - q_F)$ & $0.38232$ & $0.13813$ & $0.039802$ & $0.010376$ \\
uniform gap bound & $11.297$ & $2.9547$ & $0.74729$ & $0.18737$ \\
$(M{+}1)\times$ gap & $6.4994$ & $8.9787$ & $10.229$ & $10.635$ \\
\midrule
$\norm{W(T_m) - W(0)}_2$ & $0.219021$ & $0.071050$ & $0.019658$ & $0.005064$ \\
displacement bound & $1.370601$ & $0.358465$ & $0.090662$ & $0.022732$ \\
bound\,/\,displacement & $6.258$ & $5.045$ & $4.612$ & $4.489$ \\
\midrule
normalized-readout control: $a(T_m)$ & $1.026947$ & $1.026947$ & $1.026947$ & $1.026947$ \\
normalized-readout control: $\norm{W(T_m){-}W(0)}$ & $1.177117$ & $1.177117$ & $1.177117$ & $1.177117$ \\
\midrule
residual envelope excess on $[0, 20T_m]$ & $0$ & $0$ & $0$ & $0$ \\
readout spread $\max_j\beta_j - \min_j\beta_j$ at $T_m$ & $0$ & $0$ & $1.4\cdot10^{-17}$ & $0$ \\
invariant $q = Q_M(a)$, max rel.\ error & $2.0\cdot10^{-11}$ & $6.7\cdot10^{-13}$ & $1.5\cdot10^{-13}$ & $6.8\cdot10^{-14}$ \\
\bottomrule
\end{tabular}
\end{table}

The bound stays within $4.5$ to $6.3$ times the realised displacement across
this width range while both fall by more than a factor of $40$. The
normalized-readout control, in which the readout is $M^{-1/2}Uh$ with $O(1)$
entries trained at unit rate, removes the width dependence exactly: the
encoder's movement is the $M = 1$ value at every width, to $1.7\times10^{-14}$
relative error. The eigenvalue constants $\lambda_\thetap(0) = 1/4$ and
$\lambda_\phip(0) = M/4$ hold for the mean-over-examples binary cross-entropy
with a unit-scaled scalar logit; summing instead of averaging, or rescaling the
logit, multiplies both blocks equally and leaves $\Dcurv(0) = M$ unchanged.

\paragraph{Isotropic extension.} Table~\ref{tab:app_toy_iso} collects the
verification of the general-initialization version, in which $a(0) = a_0$,
$v(0) = v_0 \neq 0$ and the readouts start unequal but sum to exactly zero.
The grid is $M \in \{1, 4, 16, 64, 256, 1024, 4096\}$, $a_0 \in \{-1.5, -0.3,
0, 0.4, 1.7\}$, $v_0 \in \{-1.3, -0.2, 0.05, 0.2, 1.0\}$ and $m \in \{2.0,
2.9, 4.0\}$: $525$ non-trivial cases plus $54$ immediate-stop cases with
$a_0 \geq m$, integrated with DOP853 at \texttt{rtol}~$10^{-10}$ and
\texttt{atol}~$10^{-12}$ with a terminal event at $q = m$.

\begin{table}[htbp]
\centering
\small
\setlength{\tabcolsep}{4pt}
\caption{Isotropic serial instance: verification summary and worst-case bound
slacks over the whole grid, from
\texttt{results/toy\_serial\_ce\_isotropic\_REPORT.md} \S1--\S2. Slacks are
$\text{bound} - \text{value}$ for upper bounds and $\text{value} -
\text{bound}$ for lower bounds; a negative entry would be a violation, and
there are none. $\Delta = m - a_0$.}
\label{tab:app_toy_iso}
\begin{tabular}{lr}
\toprule
Quantity & Value \\
\midrule
Max closed-form vs ODE relative error over $u$, $v(T_m)$, $b(T_m)$, $T_m$ (525 cases) & $1.260\times10^{-11}$ \\
Max relative error of the conserved quantity along the whole trajectory & $1.518\times10^{-10}$ \\
Max root accuracy $|Q_{\mathrm{iso}}(u_M) - m|$ & $3.553\times10^{-15}$ \\
Max $|\sum_j \beta_j(0)|$ (zero-sum initialization is bitwise exact) & $0$ \\
Max spread of $\beta_j(T_m) - \beta_j(0)$ & $7.772\times10^{-16}$ \\
Max frozen-comparator ODE vs closed-form relative error & $5.993\times10^{-11}$ \\
Max $M^{-1/2}$-readout control vs its $M = 1$ value & $1.335\times10^{-12}$ \\
Bound violations, all cases and all bounds & $0$ \\
\midrule
Min slack, $u_M \le \Delta/(1 + Mv_0^2)$ & $2.221\times10^{-10}$ \\
Min slack, $u_M > 0$ & $4.333\times10^{-5}$ \\
Min slack, $|v(T_m) - v_0| \le \Delta^2/(2M|v_0|^3)$ & $1.493\times10^{-9}$ \\
Min slack, encoder displacement bound & $3.918\times10^{-10}$ \\
Min slack, $T_m \le \Psi_{a_0}(m)/(1 + Mv_0^2)$ & $1.748\times10^{-9}$ \\
Min slack, $T_m \ge \Psi_{a_0}(m)/(1 + Mv_0^2 + 2\Delta^2/v_0^2)$ & $3.175\times10^{-9}$ \\
Min slack, $q_J(t) \ge q_F(t)$ & $0$ \\
Min slack, $\sup_t (q_J - q_F) \le C_0/(p_0 M v_0^2)$ & $1.495\times10^{-4}$ \\
\midrule
Reversal error at $\sigma = 1$, $m = 2$, $M = 64\,/\,1024\,/\,16384$ & $0.8207\,/\,0.8380\,/\,0.8397$ \\
Monte-Carlo limit $\Phi(m/2\sigma)$ on the same $4000$ draws & $0.8413$ \\
Draws disagreeing with the limit that satisfy $Mv_0^2 > m/(m - 2a_0)$ & $0$ of $48000$ \\
Spectral-only comparator reversal error, every case & $0$ \\
\bottomrule
\end{tabular}
\end{table}

The update fraction $(a(T_m) - a_0)/(m - a_0)$ falls towards zero in $M$ at
every $(a_0, v_0)$ on the grid --- for instance $0.288$, $0.0423$, $0.000964$
at $M = 1, 16, 1024$ with $a_0 = -1.5$, $v_0 = 1$, $m = 2.9$ --- while
$a(T_m)/m$ and $a(T_m)/a_0$ converge to $a_0/m$ and to $1$ instead, so those
two ratios are not the quantities that vanish. Under the $M^{-1/2}$ readout
the suppression is $0.336399$ at every width, against $0.336$, $0.0472$ and
$0.000969$ for the unnormalized readout at $M = 1$, $16$ and $1024$. Every
Monte-Carlo draw that falls short of the reversal limit fails the sufficient
finite-width condition, either because $|v_0|$ is small or because $a_0$ sits
just below $m/2$, where the threshold diverges.


\section{Production model details}
\label{app:production}

\subsection{Pipeline, capacities and data}
\label{app:prod_setup}

The data are $336 \times 336$ pixel tissue-microarray cores collected by
quantum-cascade-laser infrared spectroscopy over the $1000$--$1800$~cm$^{-1}$
fingerprint region, with per-pixel four-class tissue labels. Each pixel carries
an $S = 942$ dimensional vector, obtained by stacking baseline-corrected
absorbance with its first two derivatives. The compositional model is the Block
Vision Transformer of \citet{oleary2026spatial}: a linear per-pixel
$f_\thetap : \R^{942} \to \R^{K}$, then $16 \times 16$ patch tokenization, a
12-layer, 12-head vision transformer of embedding width $M$, a
transposed-convolution decoder and a convolutional segmentation head. Parameter
counts in Table~\ref{tab:app_capacity} are read from the instantiated model.

\begin{table}[htbp]
\centering
\caption{Module capacities of the production pipeline, counted from the
instantiated model. $K$ is the spectral bottleneck width. The last row is the
single-channel field-standard input convention, for comparison.}
\label{tab:app_capacity}
\begin{tabular}{lrrr}
\toprule
Configuration & $\Cf$ & $\Cg$ & $\Cg/\Cf$ \\
\midrule
$K = 64$, three-channel input  & $61{,}108$  & $18{,}271{,}540$ & $299$ \\
$K = 128$, three-channel input & $121{,}588$ & $21{,}472{,}564$ & $177$ \\
$K = 64$, single-channel input & $20{,}916$  & $18{,}271{,}540$ & $874$ \\
\bottomrule
\end{tabular}
\end{table}

\paragraph{Channel convention.} The published architecture operates on the
second-derivative spectrum alone, the field-standard preprocessing in infrared
histopathology. The models analysed here instead stack raw absorbance with its
first two derivatives, following a companion tokenization study (anonymous,
under review). The choice affects only $\Cf$, and the three-channel convention
is the conservative one for the present purpose: under the field standard the
capacity asymmetry is roughly three times larger. Within the spatial module the
dominant blocks are the transformer ($5.34 \times 10^6$ parameters) and the
transposed convolution ($9.44 \times 10^6$), both $\Theta(M^2)$ in the
embedding width.

\paragraph{The head hypothesis is structurally unmet.} The segmentation head is
the chain $\mathrm{Conv}_{3\times3}(M{+}K \to 96) \to \mathrm{BN} \to
\mathrm{ReLU} \to \mathrm{Conv}_{3\times3}(96 \to 48) \to \mathrm{BN} \to
\mathrm{ReLU} \to \mathrm{Conv}_{3\times3}(48 \to 4)$, so the final affine
classifier sees $48 \times 3 \times 3 = 432$ input features at every width we
build ($M \in \{48, 192, 384\}$, $K \in \{64, 128\}$): the $\Theta(M)$ feature
map is separated from the logits by two nonlinearities. The
curvature-disparity theorem requires a dense classifier over penultimate
features of dimension $\Theta(M)$ with no downstream nonlinearity, and that
hypothesis is not satisfied here. Failing a hypothesis removes a guarantee; it
does not reverse the conclusion.

\subsection{Geometry at initialization}
\label{app:prod_geometry}

Gauss--Newton blocks of the instantiated pipeline were measured on real
training batches at random initialization, by Lanczos with full
reorthogonalization on matrix-free GGN block operators, validated to $10^{-12}$
against dense computation, sweeping $M \in \{48, 96, 192, 384\}$ with three
seeds and four batches per width. Table~\ref{tab:app_geom} collects the
results.

\begin{table}[htbp]
\centering
\small
\setlength{\tabcolsep}{4pt}
\caption{Production geometry at initialization. Ratios are mean $\pm$ sd over
three seeds and four batches per width. Slopes are log--log fits over the width
sweep. Effective ranks are out of $S = 942$. Directional entries are measured
at $M = 192$ over the nine measurements whose batch contains both classes
(three seeds $\times$ three batches).}
\label{tab:app_geom}
\begin{tabular}{lcc}
\toprule
Quantity & Breast & Prostate \\
\midrule
$\Dcurv$ at $M = 48$ & $0.33 \pm 0.12$ & --- \\
$\Dcurv$ at $M = 192$ & $0.46 \pm 0.21$ & --- \\
$\Dcurv$ at $M = 384$ & $0.70 \pm 0.19$ & $1.09 \pm 0.52$ \\
log--log slope of $\lambda_{\max}(G_{\thetap\thetap})$ vs $M$ & $-0.003$ & $-0.012$ \\
log--log slope of $\lambda_{\max}(G_{\phip\phip})$ vs $M$ & $0.38$ ($115 \to 258$) & $0.45$ \\
Effective rank of $\Sigma_X$, raw inputs & $1.07$ & $1.02$ \\
Effective rank of $\Sigma_X$, post-normalization & $1.21$ & --- \\
$|\cos|$ between $v_1$ and the mean spectrum (post-norm.\,/\,raw) & $>0.995$\,/\,$>0.9999$ & --- \\
$\thetap$-block top eigenvector's norm fraction in the $v_1$-induced subspace & $0.954 \pm 0.017$ & --- \\
Trace fraction in the top $40$ of $61{,}108$ spectral directions & $0.83 \pm 0.12$ & --- \\
Curvature ratio, $v_1$ vs full class contrast ($|\cos(d,v_1)| = 0.45$) & $3.7\times$ ($2.6$--$4.8$) & --- \\
Curvature ratio, $v_1$ vs $v_1$-orthogonalized contrast & $12$--$28\times$ & --- \\
Squared-norm fraction of the contrast retained after orthogonalization & $76$--$84\%$ & --- \\
Rank at which $\lambda_{\max}(G_{\phip\phip})$ overtakes the $\thetap$ spectrum & $2$--$5$ & --- \\
\bottomrule
\end{tabular}
\end{table}

Two controls bracket the denominator. Removing the spectral normalization
deflates $\lambda_{\max}(G_{\thetap\thetap})$ by $2.4\times$ ($471 \to 194$ at
$M = 192$) and pushes the ratio across parity ($1.18$ at $M = 192$, $1.31$ at
$M = 384$). Input \emph{centering} changes nothing when the architecture
contains internal batch normalization: $\Dcurv$ is identical to three decimals
with and without centering, because a linear projection followed by
normalization absorbs constant input shifts exactly. The near-rank-one input
statistics make the cap side of the curvature-disparity bound uninformative at
buildable widths independently of the head-hypothesis failure: the data
concentrate almost all their second-moment mass in one direction, and that
direction is the mean spectrum. The directional measurements establish an
anisotropy, not its consequence for learning; whether it starves learning is a
hypothesis, not a measurement.

\subsection{What carries the width trend}
\label{app:prod_width}

The measurements in this subsection are made on the production model at random
initialization, at three widths $M \in \{48, 192, 384\}$ and three seeds, on
one core (the densest of the first $24$ breast fold-0 training cores, with
$29{,}664$ labelled pixels of $112{,}896$), with cuDNN TF32 disabled. Four
normalization modes are compared: \texttt{train} (both head batch
normalizations using batch statistics), \texttt{eval\_bn1} and
\texttt{eval\_bn2} (only the first or only the second head normalization at its
initial running statistics) and \texttt{eval\_head} (both). Artifacts:
\texttt{results/exp1\_8c\_interior\_witness.csv} and
\texttt{results/exp1\_8d\_REPORT.md}.

\paragraph{The incoming features are flat in width.} The only downstream layer
whose input dimension grows with $M$ is the first head convolution. The top
eigenvalue of the patch-unfolded second-moment matrix of the features entering
it is unchanged to four significant figures across the three widths, at
$356.3$, $403.7$ and $332.4$ for seeds $0$, $1$ and $2$ respectively; the
leading direction sits inside the width-independent $K = 64$ skip channels.
Table~\ref{tab:app_gram} also records a discrepancy that matters for any
normalization argument: the same Gram computed over the full batch-normalization
group (all positions, ignored pixels included) rather than over labelled pixels
is substantially smaller, and the cosine between the two centered top
eigenvectors is far from one.

\begin{table}[htbp]
\centering
\small
\caption{Incoming $3\times3$-patch Gram of the first head convolution, top
eigenvalue, from \texttt{results/exp1\_8d\_REPORT.md}. Values are identical
across $M \in \{48, 192, 384\}$ and across all four normalization modes, so a
single column suffices. \emph{labelled} uses the labelled-pixel mask;
\emph{full group} uses the batch-normalization group. The last column is the
cosine between the two centered top eigenvectors.}
\label{tab:app_gram}
\begin{tabular}{rrrrr}
\toprule
seed & labelled & labelled, centered & full group & $\cos$(labelled, full) \\
\midrule
$0$ & $356.3$ & $68.19$ & $263.3$ & $0.2793$ \\
$1$ & $403.7$ & $54.60$ & $306.8$ & $0.5502$ \\
$2$ & $332.4$ & $67.14$ & $253.0$ & $0.2176$ \\
\bottomrule
\end{tabular}
\end{table}

\paragraph{Normalization quantities.} Table~\ref{tab:app_bn} gives the
per-channel batch statistics of the two head normalizations, measured over the
group they actually normalize, together with the activation energy entering and
leaving the first one. Under fan-in initialization the pre-normalization
variance of the first head normalization falls roughly as $1/(M+K)$, and its
inverse-standard-deviation gain grows correspondingly. The stabilizer
$\varepsilon = 10^{-5}$ is never within three orders of magnitude of the
variance: the fraction of channels with variance below $10\varepsilon$ is $0$
in every row of both normalizations at every width, seed and mode.

\begin{table}[htbp]
\centering
\small
\caption{Head normalization quantities at initialization, \texttt{train} mode,
from \texttt{results/exp1\_8d\_REPORT.md} Table 1. Per-seed values are
per-channel medians (variance, gain) or labelled-pixel energies; the
\emph{mean} column is the mean of the three seed values, computed by us.
$\mathrm{gain}_i = \gamma_i/\sqrt{\mathrm{var}_i + \varepsilon}$. The last row
is the product of the mean variance with $M + K$, $K = 64$.}
\label{tab:app_bn}
\begin{tabular}{lrrrrr}
\toprule
Quantity & seed $0$ & seed $1$ & seed $2$ & mean & $M$ \\
\midrule
First head BN, median variance & $0.1470$ & $0.1467$ & $0.1587$ & $0.1508$ & $48$ \\
                               & $0.0678$ & $0.0589$ & $0.0656$ & $0.0641$ & $192$ \\
                               & $0.0346$ & $0.0310$ & $0.0398$ & $0.0351$ & $384$ \\
First head BN, median gain     & $2.609$ & $2.611$ & $2.511$ & $2.577$ & $48$ \\
                               & $3.842$ & $4.120$ & $3.905$ & $3.956$ & $192$ \\
                               & $5.375$ & $5.682$ & $5.010$ & $5.356$ & $384$ \\
Second head BN, median gain    & $2.892$ & $3.122$ & $3.039$ & $3.018$ & $48$ \\
                               & $3.110$ & $3.177$ & $3.195$ & $3.161$ & $192$ \\
                               & $2.966$ & $3.060$ & $3.121$ & $3.049$ & $384$ \\
Pre-BN energy, labelled pixels & $25.64$ & $26.28$ & $23.91$ & $25.28$ & $48$ \\
                               & $11.07$ & $10.67$ & $11.19$ & $10.98$ & $192$ \\
                               & $5.765$ & $5.543$ & $6.015$ & $5.774$ & $384$ \\
Post-BN energy, labelled pixels & $132.1$ & $130.8$ & $133.2$ & $132.0$ & $48$ \\
                                & $130.4$ & $131.5$ & $132.5$ & $131.5$ & $192$ \\
                                & $130.7$ & $128.5$ & $130.8$ & $130.0$ & $384$ \\
\midrule
Mean variance $\times\,(M{+}K)$ & --- & --- & --- & $16.9\,/\,16.4\,/\,15.7$ & $48\,/\,192\,/\,384$ \\
\bottomrule
\end{tabular}
\end{table}

\paragraph{Witness Jacobian-vector products.} A unit perturbation of the first
head convolution is propagated by forward-mode differentiation through the
actual head in the actual normalization mode, so that the train-mode derivative
includes the dependence on the batch statistics. The final scalar
$e_{\mathrm{final}} = N^{-1}\sum_n d_n^\top(\mathrm{diag}(p_n) - p_np_n^\top)d_n$
is a lower-bound witness for $\lambda_{\max}(G_{\phip\phip})$ and is verified
against the same matvec the eigensolver uses; all $216$ rows pass at relative
discrepancy below $10^{-3}$, with a maximum of $1.269 \times 10^{-5}$ over the
$198$ rows with nonzero curvature. Table~\ref{tab:app_jvp} gives the stage
energies for two directions in two modes.

\begin{table}[htbp]
\centering
\small
\caption{Witness JVP stage energies, mean over three seeds, from
\texttt{results/exp1\_8d\_REPORT.md} Table 2a. \texttt{ggn\_top} is the full
top GGN eigenvector of the convolution block; \texttt{s9cen} is the centered
leading direction of the incoming patch Gram. $e_{\mathrm{pre}}$ is the
perturbation energy before the first head normalization,
$e_{\mathrm{logit}}$ is measured on labelled pixels.}
\label{tab:app_jvp}
\begin{tabular}{llrrrrrr}
\toprule
Mode & Direction & $M$ & $e_{\mathrm{pre}}$ & $e_{\mathrm{bn1}}$ & $e_{\mathrm{bn2}}$ & $e_{\mathrm{logit}}$ & $e_{\mathrm{final}}$ \\
\midrule
\texttt{train} & \texttt{ggn\_top} & $48$  & $260.6$ & $2875$ & $4453$ & $261.7$ & $60.73$ \\
\texttt{train} & \texttt{ggn\_top} & $192$ & $259.4$ & $5905$ & $7553$ & $505.2$ & $118.5$ \\
\texttt{train} & \texttt{ggn\_top} & $384$ & $256.8$ & $9168$ & $13810$ & $1018$ & $227.6$ \\
\texttt{train} & \texttt{s9cen}    & $48$  & $76.53$ & $1538$ & $1224$ & $29.23$ & $5.761$ \\
\texttt{train} & \texttt{s9cen}    & $192$ & $76.53$ & $2305$ & $1914$ & $50.83$ & $9.862$ \\
\texttt{train} & \texttt{s9cen}    & $384$ & $76.53$ & $3997$ & $3906$ & $124.3$ & $21.95$ \\
\midrule
\texttt{eval\_head} & \texttt{ggn\_top} & $48$  & $261.9$ & $261.9$ & $42.42$ & $9.607$ & $2.366$ \\
\texttt{eval\_head} & \texttt{ggn\_top} & $192$ & $262.8$ & $262.8$ & $42.18$ & $10.35$ & $2.487$ \\
\texttt{eval\_head} & \texttt{ggn\_top} & $384$ & $262.5$ & $262.5$ & $46.53$ & $10.45$ & $2.372$ \\
\texttt{eval\_head} & \texttt{s9cen}    & $48$  & $76.53$ & $76.53$ & $9.103$ & $0.4667$ & $0.1049$ \\
\texttt{eval\_head} & \texttt{s9cen}    & $192$ & $76.53$ & $76.53$ & $8.460$ & $0.4613$ & $0.1013$ \\
\texttt{eval\_head} & \texttt{s9cen}    & $384$ & $76.53$ & $76.53$ & $9.762$ & $0.4450$ & $0.1011$ \\
\bottomrule
\end{tabular}
\end{table}

A spatially constant bias perturbation $\Delta b = e_c$ is annihilated exactly
by the immediately following train-mode normalization: in \texttt{train} and
\texttt{eval\_bn2} modes, at every width and seed, all stage energies after the
centering step are exactly zero and the quadratic form agrees with the matvec
to at most $1.14\times10^{-17}$ in absolute value. In \texttt{eval\_head} the
same perturbation gives $e_{\mathrm{final}}$ between $7.8\times10^{-4}$ and
$2.0\times10^{-3}$, and in \texttt{eval\_bn1} between $0.0154$ and $0.0692$.

\paragraph{Slopes and attribution.} Table~\ref{tab:app_slopes} gives the
log--log slopes against $M$ per seed and pooled, together with the
mode-to-mode ratios of levels.

\begin{table}[htbp]
\centering
\small
\caption{Log--log slopes against $M$ over the three widths, per seed and
pooled, and mode-to-mode ratios of the levels (each mode's value divided by its
train-mode value at the same width and seed, then averaged over seeds and over
the three widths), from
\texttt{results/exp1\_8d\_REPORT.md}. $\lambda_{WW}$ is the top curvature of
the first head convolution's weight block; $\lambda_\phip$ is the whole spatial
block. Three-width slopes are descriptive, not asymptotic exponents.}
\label{tab:app_slopes}
\begin{tabular}{lrrrrrrr}
\toprule
& \multicolumn{4}{c}{slope of $\lambda_{WW}$} & \multicolumn{1}{c}{$\lambda_\phip$} & \multicolumn{2}{c}{level ratio vs \texttt{train}} \\
\cmidrule(lr){2-5}\cmidrule(lr){6-6}\cmidrule(lr){7-8}
Mode & seed $0$ & seed $1$ & seed $2$ & pooled & pooled slope & $\lambda_{WW}$ & $\lambda_\phip$ \\
\midrule
\texttt{train}     & $0.509$ & $0.785$ & $0.350$ & $+0.613$ & $+0.398$ & $1$ & $1$ \\
\texttt{eval\_bn1} & $0.429$ & $0.634$ & $0.542$ & $+0.563$ & $+0.281$ & $0.903$ & $1.03$ \\
\texttt{eval\_bn2} & $0.544$ & $0.632$ & $0.508$ & $+0.571$ & $+0.253$ & $0.156$ & $0.227$ \\
\texttt{eval\_head} & $-0.112$ & $0.181$ & $-0.084$ & $+0.006$ & $-0.331$ & $0.026$ & $0.037$ \\
\bottomrule
\end{tabular}
\end{table}

Three readings, and one non-reading. The positive width trend in the spatial
block's curvature at initialization is sensitive to head batch normalization:
holding both head normalizations at their initial running statistics removes it
($+0.613 \to +0.006$ for the convolution block, $+0.398 \to -0.331$ for the
whole spatial block) while leaving the incoming features unchanged. Either
normalization alone preserves the trend ($+0.563$ and $+0.571$), which is
evidence against attributing it uniquely to the first one, even though most of
the drop in the \emph{level} attaches to the second. The mechanism this
supports is a finite-range normalization gain acting on width-diluted
activations, over the range of widths measured, and not the width-linear
feature-Gram mechanism of the shallow theory. The non-reading: the
\texttt{eval\_head} intervention changes the normalization derivative, the
forward logits, the downstream gates and the cross-entropy metric at the same
time, so a difference between modes establishes sensitivity to the
intervention, not a unique cause; and the observation that the variance never
approaches the stabilizer is a heuristic about the regime, not a guarantee.

\paragraph{Parameterization control.} Table~\ref{tab:app_scale} scales the
weights and bias of the first head convolution by $c$ and, in the exact
invariance, the following stabilizer by $c^2$. The logits are unchanged and the
block's top curvature scales by exactly $c^{-2}$: the raw Euclidean curvature
of a layer feeding a normalization is a parameterization quantity, not an
architectural one \citep{vanlaarhoven2017l2}.

\begin{table}[htbp]
\centering
\small
\caption{Function-preserving scale control at seed $0$, from
\texttt{results/exp1\_8d\_REPORT.md} Table 3. \emph{ratio} is
$c^2\lambda_{WW}(c)/\lambda_{WW}(1)$, which is $1$ under exact invariance.
\emph{logit dev.} is $\max|\text{logits}(c) - \text{logits}(1)|$, reported both
with the stabilizer co-scaled by $c^2$ and with it held fixed (the
implementation's actual behaviour); $\lambda_{WW}$, \emph{ratio} and
$\norm{W}_F^2$ are taken from the fixed-stabilizer variant, and the co-scaled
variant differs from it by at most one unit in the fourth significant figure.}
\label{tab:app_scale}
\begin{tabular}{rrrrrrr}
\toprule
$M$ & $c$ & $\lambda_{WW}$ & ratio & logit dev.\ (co-scaled) & logit dev.\ (fixed $\varepsilon$) & $\norm{W}_F^2$ \\
\midrule
$48$  & $1/4$ & $719.4$ & $1.000$ & $0$ & $9.40\cdot10^{-4}$ & $2.003$ \\
$48$  & $1/2$ & $179.8$ & $1.000$ & $0$ & $1.89\cdot10^{-4}$ & $8.014$ \\
$48$  & $1$   & $44.96$ & $1.000$ & $0$ & $0$ & $32.06$ \\
$48$  & $2$   & $11.24$ & $1.000$ & $0$ & $4.63\cdot10^{-5}$ & $128.2$ \\
$48$  & $4$   & $2.810$ & $1.000$ & $0$ & $5.84\cdot10^{-5}$ & $512.9$ \\
$192$ & $1/4$ & $1935$  & $0.9998$ & $0$ & $3.00\cdot10^{-3}$ & $2.002$ \\
$192$ & $4$   & $7.560$ & $1.000$ & $0$ & $1.89\cdot10^{-4}$ & $512.6$ \\
$384$ & $1/4$ & $1927$  & $0.9991$ & $0$ & $4.69\cdot10^{-3}$ & $1.999$ \\
$384$ & $4$   & $7.535$ & $1.000$ & $0$ & $2.95\cdot10^{-4}$ & $511.8$ \\
\bottomrule
\end{tabular}
\end{table}

\paragraph{Eigensolver accuracy.} Every row of the width sweep terminates on
the combined criterion (relative change of the Rayleigh quotient below
tolerance \emph{and} residual below $10^{-3}$), never on the iteration or
wall-clock budget. Over all $36$ (mode, width, seed) rows the relative residual
$\norm{Gv - \lambda v}/\lambda$ is at most $3.81\times10^{-4}$ for the
convolution weight block, at most $3.81\times10^{-4}$ for the weight--bias
block and at most $9.97\times10^{-4}$ for the whole spatial block, and the
Rayleigh quotient of the returned iterate agrees with the reported eigenvalue
to the printed precision in every row.

\subsection{Controlled retraining: original protocol}
\label{app:prod_e3c}

The pipeline was retrained on breast fold $0$ at reduced depth (six transformer
layers) with full instrumentation. Two protocols were run and both are
reported: the \emph{original} protocol ($39$ runs), which is the pre-registered
one and whose joint-versus-frozen contrast an audit of the training code later
showed to be confounded, and the \emph{matched} protocol ($21$ runs), which
removes the confounds, saves the selected checkpoint and re-runs the
pre-registered comparisons. Sixty logged runs in total. Per-step residuals,
block gradient norms, per-epoch parameter displacements and input-Jacobian
norms are recorded in both; the matched protocol additionally logs the
per-step clipping coefficient and the per-block applied update norms.

The original protocol comprises arms $\{$joint-linear, frozen-random,
frozen-PCA, joint-MLP$\}$ $\times$ widths $M \in \{48, 192\}$ $\times$ three
seeds under AdamW for $60$ epochs, plus a $\{$joint-linear, frozen-random$\}$
$\times$ three-seed pair under momentum SGD at both widths and a frozen-PCA SGD
triple at $M = 48$ as a positive control. The spectral parameters sit in a
zero-weight-decay group in every arm, and frozen arms log counterfactual
spectral gradients. Table~\ref{tab:app_e3c_orig} gives the AdamW arms.

\begin{table}[htbp]
\centering
\small
\caption{Controlled study, original protocol, AdamW arms: validation macro-F1
(mean $\pm$ sd over three seeds) under two checkpoint policies, and relative
spectral displacement at end of training. \emph{Final-5} is the pre-registered
metric. \emph{Best-val} is the maximum over the $60$ per-epoch evaluations on
the same $28$ validation cores; no checkpoint was saved at that maximum and
this study has no test set, so that column is exploratory and
selection-optimistic. The two policies order the arms oppositely. These arms
also differ in clip scope and in normalization-affine treatment, so this
contrast is confounded.}
\label{tab:app_e3c_orig}
\begin{tabular}{lccccr}
\toprule
& \multicolumn{2}{c}{Final-5 (pre-registered)} & \multicolumn{2}{c}{Best-val (exploratory)} & \\
\cmidrule(lr){2-3}\cmidrule(lr){4-5}
Arm & $M{=}48$ & $M{=}192$ & $M{=}48$ & $M{=}192$ & $\norm{\Delta\thetap}/\norm{\thetap_0}$ \\
\midrule
Joint (linear) & $0.538 \pm 0.082$ & $0.646 \pm 0.050$ & $0.873 \pm 0.030$ & $0.858 \pm 0.056$ & $0.54$--$0.64$ \\
Frozen random  & $0.648 \pm 0.063$ & $0.731 \pm 0.023$ & $0.840 \pm 0.025$ & $0.817 \pm 0.017$ & $0$ \\
Frozen PCA     & $0.731 \pm 0.003$ & $0.726 \pm 0.012$ & $0.769 \pm 0.021$ & $0.862 \pm 0.068$ & $0$ \\
Joint (MLP)    & $0.566 \pm 0.070$ & $0.575 \pm 0.184$ & $0.900 \pm 0.000$ & $0.897 \pm 0.001$ & $0.85$--$0.96$ \\
\bottomrule
\end{tabular}
\end{table}

On the pre-registered final-window metric the joint-linear minus
frozen-random gaps are $-0.110$ at $M = 48$ ($90\%$ paired-$t$ interval
$[-0.333, +0.114]$) and $-0.085$ at $M = 192$ ($[-0.185, +0.015]$), neither
resolvable at $n = 3$; against frozen-PCA they are $-0.193$
($[-0.334, -0.052]$) and $-0.081$ ($[-0.175, +0.014]$). Taking each run's
maximum over its $60$ evaluations reverses the ordering: joint minus
frozen-random becomes $+0.033$ and $+0.041$, joint minus frozen-PCA $+0.104$
and $-0.004$. The switch is an arithmetic identity --- the change in every gap
equals the difference in peak-to-final degradation between the two arms ---
and averaged over seeds that degradation is $0.335$ ($M = 48$) and $0.212$
($M = 192$) for joint-linear against $0.192$ and $0.086$ for frozen-random.

\paragraph{Three asymmetries.} First, gradient-norm clipping at $1.0$ was
computed over the optimizer's parameters, that is over $\thetap + \phip$ in the
joint arms but over $\phip$ alone in the frozen arms, so at identical gradients
the frozen arm takes the larger $\phip$ step: the median extra clip factor paid
by the joint arm is $1.10$--$1.12$ under SGD at $M = 48$ and $3.0$--$3.6$ under
AdamW at $M = 192$. Second, the normalization following the spectral projection
contributes $128$ affine parameters that are trained in joint-linear, frozen in
frozen-random and structurally absent in frozen-PCA. Third, no checkpoint was
saved during training, so the best-validation column is a post-hoc curve
maximum that no saved model realizes, taken over $60$ evaluations of the same
$28$ validation cores that define the comparison.

\subsection{Momentum-SGD controls and normalization recalibration}
\label{app:prod_sgd}

The two optimizer families constitute an optimizer-family control; the SGD arm
is not an instantiation of the two-timescale theorem's gradient flow. It uses
momentum $0.9$; weight decay $0.01$ on $\phip$ and $0$ on $\thetap$;
mini-batches of four cores with reshuffling; and global gradient-norm clipping
at $1.0$, which binds on $73$--$89\%$ of steps. Logged gradient norms are
recorded before clipping, so the plotted theory quantities are not the
quantities the trajectory applied. Under momentum SGD the paired joint minus
frozen-random gap is $+0.002$ at $M = 48$ and $+0.011$ at $M = 192$ on
best-validation, with $90\%$ paired-$t$ intervals $[-0.008, +0.011]$ and
$[-0.007, +0.030]$; on the final-window metric the same pairs give $+0.028$ and
$-0.057$, with intervals $[-0.173, +0.228]$ and $[-0.124, +0.010]$ --- gaps
that differ in sign between the widths, with intervals an order of magnitude
wider than the effect they bracket. Both arms lose $0.12$--$0.22$ macro-F1
between their peak and their final window. Descriptively, the spectral share of
the summed squared gradient norms drops from about $0.94$ (AdamW joint) to
about $0.26$ (SGD joint); momentum, clipping, weight decay and adaptivity all
differ between the two families, so no single factor is isolated.

\paragraph{Recalibration of normalization statistics.} On disposable copies of
the final checkpoints, with all weights fixed, the normalization running
statistics were re-estimated on training data only ($29$ batches) and the model
re-evaluated. For this intervention and these checkpoints, the recovered
fraction of the joint arms' peak-to-final drop is $0.17$ at $M = 48$ and
$0.07$ at $M = 192$: recalibration recovers a nonzero part of the degradation
but not most of it. What it does undo is the instability of the endpoint
number: it repairs the two collapsed joint checkpoints ($+0.24$ and $+0.15$
macro-F1) and shrinks the seed standard deviation by roughly $3.7\times$, while
leaving the frozen arms essentially unchanged ($|\Delta| \le 0.02$, and
$\le 0.006$ for frozen-PCA, whose spectral stage carries no normalization).

\subsection{Matched protocol}
\label{app:prod_matched}

The matched protocol equalizes the three asymmetries and re-runs the
pre-registered comparisons at $M = 48$: the reduction-stage normalization
affine parameters are placed in $\phip$ and trained wherever that
normalization exists; the clip norm at $1.0$ is computed over $\phip$ only
under the same clipping rule and scope in every arm; and the best-validation
checkpoint is saved. AdamW at learning rate $10^{-4}$, $60$ epochs, three
seeds, flip and $90^\circ$-rotation augmentation on in all arms. The logged
per-step clip coefficient averages $0.33$--$0.57$ across arms --- the
coefficients and gradients differ by arm, which is an outcome of the matched
rule --- and the logged applied update norms confirm that the spectral share of
the applied update tracks the learning-rate multiplier ($0.02$, $0.17$ and
$1.36$ times the $\phip$ update norm for $\times0.1$, $\times1$ and
$\times10$). Table~\ref{tab:app_e3c_matched} gives all seven arms.

\begin{table}[htbp]
\centering
\small
\caption{Controlled study, matched protocol ($M = 48$, AdamW, three seeds,
mean $\pm$ sd). Identical normalization-affine treatment and the same
$\phip$-only clipping rule and scope in every arm; best-validation checkpoints
saved. \emph{Final-5} is the pre-registered metric; \emph{best-val} is the
maximum over $60$ per-epoch evaluations on the $28$ fold-0 validation cores and
is exploratory and selection-optimistic, since this study has no test set. The
drift column is the relative spectral-parameter displacement at end of
training; its standard deviations were computed by us from
\texttt{results/e3c\_analysis\_runs.csv}. Rows $1$--$2$ and rows $3$--$4$ are
the two matched pairs; see the text.}
\label{tab:app_e3c_matched}
\begin{tabular}{lccc}
\toprule
Arm & Final-5 (pre-registered) & Best-val (exploratory) & $\norm{\Delta\thetap}/\norm{\thetap_0}$ \\
\midrule
Frozen random                   & $0.684 \pm 0.037$ & $0.814 \pm 0.040$ & $0$ \\
Joint (linear)                  & $0.694 \pm 0.056$ & $0.837 \pm 0.029$ & $0.566 \pm 0.026$ \\
Frozen pretrained               & $0.690 \pm 0.033$ & $0.861 \pm 0.022$ & $0$ \\
Fine-tuned pretrained           & $0.710 \pm 0.052$ & $0.857 \pm 0.009$ & $0.020 \pm 0.005$ \\
Joint, spectral lr $\times 0.1$ & $0.626 \pm 0.055$ & $0.841 \pm 0.007$ & $0.083 \pm 0.006$ \\
Joint, spectral lr $\times 10$  & $0.561 \pm 0.114$ & $0.886 \pm 0.037$ & $4.910 \pm 0.402$ \\
Joint, cosine schedule          & $0.660 \pm 0.020$ & $0.825 \pm 0.033$ & $0.359 \pm 0.013$ \\
\bottomrule
\end{tabular}
\end{table}

\paragraph{Two matched pairs, not a single-factor design.} The four core arms
do not all differ in one factor. Joint-linear and frozen-random use a linear
reduction stage with normalization; frozen-pretrained and fine-tuned-pretrained
use a pretrained multilayer spectral encoder without reduction-stage
normalization. The clean freezing comparisons are therefore the two matched
pairs, joint-linear against frozen-random and fine-tuned against
frozen-pretrained. Comparing across the pairs --- frozen-pretrained against
frozen-random, for instance --- changes architecture, initialization and
training history, and validation selection, not only informativeness.

\paragraph{Verdicts.} All comparisons in Table~\ref{tab:app_verdicts} are
paired by seed with $90\%$ $t$-intervals ($t = 2.920$, $n = 3$), conditional on
one validation split, and are descriptive seed summaries rather than a
population or multi-factor causal analysis.

\begin{table}[htbp]
\centering
\small
\setlength{\tabcolsep}{4pt}
\caption{Matched protocol: paired differences with $90\%$ paired-$t$ intervals
($n = 3$). The last two rows compare the two protocols and are reported as
descriptive seed summaries; saving a checkpoint changes what is available to
evaluate, not the final-window trajectory.}
\label{tab:app_verdicts}
\begin{tabular}{lcc}
\toprule
Contrast & Final-5 & Best-val \\
\midrule
Fine-tuned $-$ frozen-pretrained & $+0.021\ [-0.022, +0.063]$ & $-0.004\ [-0.048, +0.039]$ \\
Joint-linear $-$ frozen-random & $+0.011\ [-0.142, +0.163]$ & $+0.023\ [-0.066, +0.112]$ \\
Frozen-pretrained $-$ frozen-random & --- & $+0.047\ [-0.042, +0.136]$ \\
Spectral lr $\times10 - \times1$ & $-0.133\ [-0.419, +0.152]$ & $+0.048\ [-0.016, +0.112]$ \\
Spectral lr $\times0.1 - \times1$ & $-0.068\ [-0.173, +0.037]$ & $+0.004\ [-0.039, +0.046]$ \\
\midrule
Joint/frozen gap: matched $-$ original & $+0.120\ [-0.230, +0.470]$ & $-0.010\ [-0.105, +0.085]$ \\
\bottomrule
\end{tabular}
\end{table}

The pre-registered prediction that fine-tuning a pretrained spectral encoder
falls below leaving it frozen is not supported: neither difference is
resolvable at $n = 3$, and the fine-tuned arm moved $\thetap$ by only $2\%$ of
its initial norm. The logged applied update norms verify nonzero applied
encoder updates under AdamW in that arm; this is a statement about the applied
updates, not an attribution of the loss decrease. The pre-registered
harmful-drift ordering $\times0.1 \ge \times1 \ge \times10$ is not observed: on
best-validation the $\times10$ arm is above $\times1$ in $3/3$ seeds while its
final-window endpoint is the lowest and most variable of any arm at $0.561$.
Three multiplier means do not establish a systematic relation between
displacement and peak score, and the intervals are wide. The cosine schedule
improved neither mean endpoint ($0.825$ best-val and $0.660$ final-5 against
$0.837$ and $0.694$ for the constant learning rate); one schedule that did not
help does not rule out an excessive late learning rate.

\paragraph{Disclosures.} Width $48$ only; $n = 3$; a single validation fold of
$28$ fold-0 cores; no test set. The best-validation column is exploratory and
the final-window column is the pre-specified metric. The pretrained encoders
were selected by pixel validation macro-F1 on the \emph{same} fold-0 validation
cores the runner later scores, which is a selection advantage for both
pretrained arms, and were pretrained with weight decay $0.01$ on $\thetap$.
Pixel pretraining used the unaugmented centre-crop protocol while subsequent
training used random crop plus augmentation; this affects comparisons across
the two pairs, not the frozen-versus-fine-tuned pair that starts from the same
checkpoint. Augmentation is on in all arms. Finally, the summary artifact
\texttt{results/e3c\_analysis\_runs.csv} does not expose per-run configuration
hashes or data-order pairing, so before the pairs are described as completely
matched a compact provenance manifest should accompany it, recording
architecture, initial checkpoint hash, parameter-group membership, clipping and
weight-decay settings, normalization mode and state, augmentation and
data-order seed, scheduler and selection rule; identical seed labels alone do
not establish bit-identical mini-batch order across different model
constructions.

\subsection{What is not certified on this pipeline}
\label{app:prod_scope}

The spatial module's input-Jacobian operator norm grows over training, from
$3.2$ to $157$ on average for joint-linear (up to $414$) and from $1.7$ to
$1.2\times10^{4}$ for joint-MLP, which rejects an initialization-scale
stability approximation for the containment constant of the displacement
theorem, though not the existence of a finite-horizon constant over the
training window. The fitted per-step residual-envelope rate for joint-linear
grows from $\hat\mu = 6.0\times10^{-4}$ ($M = 48$) to $1.9\times10^{-3}$
($M = 192$), a factor $3.2$ for a $4\times$ width increase; the ratio is
invariant to the step-clock-to-flow-time conversion, the absolute values are
not, and two widths do not establish a trend.

The earlier proxy instantiation of the displacement bound is withdrawn. Even
after correcting the time units and using the supremum gain, the proxy exceeds
the measured displacement by roughly $17$--$47\times$ under AdamW and roughly
$450$--$1900\times$ under momentum SGD; the raw per-step residual violates the
fitted envelope on $19$--$28\%$ of joint-linear steps, so the envelope
assumption is verified for the smoothed curve rather than for the logged
quantity. Momentum, gradient clipping and adaptive preconditioning all lie
outside the theorem's hypotheses, and under AdamW the measured displacement
exceeds $\eta\sum_k\norm{\nabla_\thetap\loss_k}$ by more than an order of
magnitude, so no step-size-rescaled flow bound describes that trajectory in
either direction. No certified instantiation of the displacement bound exists
for this pipeline at present.

Finally, the scalar gradient-ratio diagnostic is blind here: it \emph{rises}
during joint AdamW training (joint-linear epoch-1 median $0.6$--$0.8$, rising
to $3$--$5$ by epoch $60$), including in runs whose validation score is
falling. A single global ratio cannot see a direction-wise anisotropy, and the
initialization geometry of Appendix~\ref{app:prod_geometry} is what such a
ratio would have to resolve.
}{\section{Synthetic experiment details}\label{app:synthetic}\emph{[Filled separately.]}\section{Production model details}\label{app:production}\emph{[Filled separately.]}}

\section{Verification record}
\label{app:audits}

Every proof reproduced or stated in this appendix was independently
re-derived line by line, and every displayed inequality was checked
numerically at the parameter values recorded below. The scripts live under
\texttt{code/audits/} in the supplementary repository; each is a standalone
program that reads and writes no repository file, and each records its own
seed and command in its header. Table~\ref{tab:audits} lists them; a seed is
a NumPy generator seed unless marked \texttt{torch}, in which case it is the
global torch seed, and \texttt{---} marks a deterministic script. The
recorded environment is Python 3.12.13, numpy 2.0.2, scipy 1.16.3, torch
2.9.1, CPU only.

\newcommand{\uscore}{\_\allowbreak}
\begin{table}[ht]
\centering
\small
\begin{tabular}{@{}p{0.235\textwidth}p{0.415\textwidth}p{0.115\textwidth}p{0.105\textwidth}@{}}
\toprule
Script & Verifies & Seed & Result \\
\midrule
\texttt{audit\uscore A1.py}
  & Discrete displacement and clipping bounds; monotone comparison and its
    Gr\"onwall form; affine bridge, its counterexample and the attribution
    share; the top-eigenvalue counterexample; the constants of the
    initialization certificate
  & \texttt{np 0} & PASS (18/19) \\
\texttt{audit\uscore a2.py}
  & Interior curvature witness identity and the ReLU interior theorem's
    probability and curvature-floor events; the normalization
    counterexamples; the cross-entropy gradient-gain counterexample
  & \texttt{np 0} & partial \\
\texttt{audit\uscore a2\uscore bn.py}
  & The normalization annihilation counterexample against a real
    \texttt{BatchNorm1d} in both stabilizer regimes; the width-proportional
    weighted feature Gram; the mean-direction witness with non-constant
    features
  & \texttt{torch 1} & PASS \\
\texttt{check\uscore p5.py}
  & Solvable serial instance: invariants, fitting-time bracket, residual
    envelope, Gauss--Newton constants, gradient-norm bound, joint-versus-%
    frozen gap, displacement table, normalized-readout control
  & --- & PASS \\
\texttt{audit\uscore blocks.py}
  & In the paper's own normalization: the gradient identity
    $\nabla_b\loss = J_b^\top r$, the head-weight block identity for a
    shared-weight $1{\times}1$ convolutional head, and the degenerate
    stationary-point counterexample to Corollary~\ref{cor:attribution}
  & \texttt{torch 0} & PASS \\
\texttt{check\uscore theoremN.py}
  & Scoped nonlinear existence result: the containment inequality, the
    three-term Jacobian decomposition, gate-flip containment, perturbation
    and moment bounds, boundary mass, and the attribution share along the
    nonlinear flow
  & \texttt{np 0} & PASS \\
\texttt{astra\uscore math\uscore 02\uscore checks.py}
  & Finite-horizon attribution with non-commuting kernels by exact spectral
    integral; the isotropic serial instance with signed context and unequal
    zero-sum readouts; the normalization derivative and co-scaled
    curvature law; ReLU feature and gate tube bounds
  & \texttt{np 902} & PASS \\
\midrule
\multicolumn{4}{@{}l}{\emph{Subdirectory}
  \texttt{math02\uscore sections\uscore 2\uscore 4\uscore 6/}} \\
\texttt{a1\uscore finite\uscore horizon.py}
  & Finite-horizon attribution over 24 nonlinear runs, sharpness at equality,
    and the old-versus-new bound & \texttt{20260910} & PASS \\
\texttt{a2\uscore chernoff.py}
  & Truncated matrix Chernoff certificate: truncation radius, minimum
    eigenvalue floor, exponent constant, width threshold
  & \texttt{424242} & PASS \\
\texttt{a2\uscore chernoff\uscore biting.py}
  & The same certificate in a second parameterization, with empirical
    failure frequencies across widths & \texttt{909090} & PASS \\
\texttt{a2\uscore chernoff\uscore constant.py}
  & The constant in the certificate, stressed on the extremal
    balls-in-bins family & \texttt{5150} & PASS \\
\texttt{c23\uscore bullets.py}
  & Reduction of the full flow to its three-variable closure, residual
    identities, and the clipping equivalence against a real optimizer
  & \texttt{246810} & PASS \\
\texttt{i\uscore toy\uscore ode.py}
  & The serial instance's update suppression, displacement, fitting time and
    frozen-comparator bracket over 24 high-accuracy solves
  & \texttt{31337} & PASS \\
\texttt{i6\uscore reversal\uscore mc.py}
  & The contextual-reversal error limit, its per-draw form, the sufficient
    condition, and the tail bound & \texttt{77007} & PASS \\
\texttt{b\uscore bn\uscore torch.py}
  & Normalization Jacobian and spectrum against a real \texttt{BatchNorm2d};
    the width-constant curvature ratio; the function-preserving rescaling
    law; uniform versus channelwise rescaling & \texttt{13579} & PASS \\
\texttt{b\uscore bn\uscore crosschannel.py}
  & Exact block-diagonality of the normalization derivative across channels
  & \texttt{13579} & PASS \\
\bottomrule
\end{tabular}
\caption{Numerical verification record for the results of
Appendices~\ref{app:proofs} and~\ref{app:math}. Both non-clean outcomes are
explained in the text; neither is a mathematical failure. Scripts are kept
byte-identical to what was run.}
\label{tab:audits}
\end{table}

Two entries are not clean, and both are artefacts of the audit code rather
than of the mathematics. In \texttt{audit\_A1.py}, the single failing line
is the sub-check ``share $\leq 1$'' of the affine bridge, whose integral the
script evaluates with a rectangle rule that over-estimates a decaying
integrand (16 of 200 trials exceed one, maximum $1.73$); recomputing the
same integral in closed form gives $0$ of $300$ trials above one with
maximum $0.919$, and the identical claim passes exactly, over $800$
non-commuting-kernel instances, in \texttt{astra\_math\_02\_checks.py}. In
\texttt{audit\_a2.py}, every check passes until the script aborts inside a
functional Jacobian transform applied to a normalization \emph{module},
because the installed torch version forbids the in-place update of that
module's batch counter under the transform; the aborted and subsequent
checks are duplicated in \texttt{audit\_a2\_bn.py}, which constructs the
same layer without running statistics and passes them all, so nothing is
left unverified. Both scripts are retained unmodified so that the record
matches the run.
